\documentclass[11pt]{amsart}
\usepackage[margin=1in]{geometry}
\usepackage{amsmath,amssymb,amsthm,mathtools,booktabs,longtable,array,microtype}
\usepackage{url}
\usepackage{hyperref}
\hypersetup{colorlinks=true,linkcolor=blue,citecolor=blue,urlcolor=blue,
 pdftitle={An Exponent of 1.04273 for the Unit Distance Problem},
 pdfauthor={Eric Naslund}}
\numberwithin{equation}{section}
% Theorems and equations share one counter, so every number is unique.
\newtheorem{theorem}[equation]{Theorem}
\newtheorem{proposition}[equation]{Proposition}
\newtheorem{lemma}[equation]{Lemma}
\newtheorem{corollary}[equation]{Corollary}
\theoremstyle{definition}
\newtheorem{definition}[equation]{Definition}
\theoremstyle{remark}
\newtheorem{remark}[equation]{Remark}
\DeclareMathOperator{\Gal}{Gal}
\DeclareMathOperator{\Cl}{Cl}
\DeclareMathOperator{\rd}{rd}
\DeclareMathOperator{\Tr}{Tr}
\DeclareMathOperator{\rank}{rank}
\DeclareMathOperator{\Br}{Br}
\DeclareMathOperator{\Pic}{Pic}
\DeclareMathOperator{\Frob}{Frob}
\DeclareMathOperator{\inv}{inv}
\DeclareMathOperator{\Res}{Res}
\DeclareMathOperator{\covol}{covol}
\newcommand{\Ftwo}{\mathbb F_2}
\newcommand{\Q}{\mathbb Q}
\newcommand{\R}{\mathbb R}
\newcommand{\C}{\mathbb C}
\newcommand{\Z}{\mathbb Z}
\newcommand{\deltastar}{0.04273}
\newcommand{\thetastar}{65535/131072}
\newcommand{\ceiling}{0.04871285}
\setcounter{tocdepth}{1}
\raggedbottom
\title[The Unit Distance Problem]{An Exponent of 1.04273\\for the Unit Distance Problem}
\author[Eric Naslund]{Eric Naslund\textsuperscript{*}}
\email{naslund.math@gmail.com}
\date{September 30, 2026}
\begin{document}
\begin{abstract}
Let $u(U)$ count the unordered pairs at distance one in a finite planar
set $U$. We construct finite sets $U_j$ with $|U_j|\to\infty$ and
$u(U_j)/|U_j|^{1.04273}\to\infty$. The largest exponent previously
claimed, $1.0358$, is in the author's unpublished manuscript. The
method is the number-field construction of OpenAI and Sawin: unit
distances come from elements of relative norm one in quadratic
extensions, and the fields come from an infinite pro-$2$ class field
tower. The new ingredients are quadratic extensions of mixed signature,
with an exact average over their norm-one units, and a tower over the
real quadratic field $\Q(\sqrt{241})$, in which $2$, $3$ and $5$ split.
Its Golod--Shafarevich function contains two copies of the local
conditions at these primes but only one constant term, and the extra room
lets the tower be ramified only above $2$, $3$ and $5$; the root
discriminant of its fields is about $286$. The relative zeta value is bounded through the zeta function of the
degree-$512$ field generated over $\Q(\sqrt{241})$ by the square roots of
its $\{2,3,5\}$-units, which is the product of the Dedekind zeta function
of $\Q(\sqrt{241})$ and $255$ quadratic Hecke $L$-functions. Finite facts and numerical
inequalities are certified by exact computation and interval arithmetic. A significant portion of this work was verified in Lean, reducing the result with exponent $1.0427$ to an explicit zeta function inequality.
\end{abstract}
\maketitle

% AI Methodology statement, referenced by the asterisk on the author's name.
{\noindent\small\textsuperscript{*}\textbf{AI Methodology.} This work was done with extensive AI use, with agents developing the mathematics and writing the paper. I supervised and prompted the AI agents that produced this work. While this has been edited to improve readability, I recommend the reader also make use of their AI model of choice to ask the paper questions. The
construction over $\Q(\sqrt{241})$, the computations that certify it, the
Lean formalization of the reduction of the proof of exponent $1.0427$ to one
explicit zeta inequality (Remark~\ref{an:lean-hypothesis}), and this paper
were produced with a combination of Claude Opus 5.5, ChatGPT Sol 5.6, ChatGPT Astra 6, working as agents in Claude Code and Codex. The construction adapts the author's previous manuscript
\cite{Naslund0358}, which was also written with AI tools.\par}
\medskip

\section{Introduction}\label{intro:section}
For a finite set $U\subset\R^2$, let
\[
 u(U)=\#\{\{x,y\}\subset U:|x-y|=1\},\qquad
 u(n)=\max_{|U|=n}u(U),
\]
where the pairs are unordered and the distance is Euclidean. The unit
distance problem of Erd\H{o}s \cite{Erdos1946} asks for the order of
growth of $u(n)$. Our main result is the following.

\begin{theorem}\label{thm:main}
There are finite sets $U_j\subset\R^2$, with $|U_j|\to\infty$, such that
\[
 \frac{u(U_j)}{|U_j|^{1.04273}}\longrightarrow\infty.
\]
In particular, $u(n)\ge n^{1.04273}$ for arbitrarily large integers $n$.
\end{theorem}

The largest exponent previously claimed, $1.0418235$, is in the author's
unpublished manuscript \cite{Naslund0418235}. Theorem~\ref{thm:main}
uses the same method with one new ingredient, together with the
corresponding change in the analytic estimate: the class field tower
that supplies the number fields is built over the real quadratic field
$\Q(\sqrt{241})$ instead of over $\Q$. The theorem concerns an unbounded
sequence of cardinalities. Its proof uses no unproved hypothesis, such
as the generalized Riemann hypothesis, but it is computer-assisted:
finite facts about the tower and several numerical inequalities are
certified by exact computation and interval arithmetic, as described in
Section~\ref{intro:code}.

\subsection{Background}\label{intro:background}
Erd\H{o}s \cite{Erdos1946} showed that a suitable section of a scaled
integer lattice gives
\[
 u(n)\ge n^{1+c/\log\log n}
\]
for some $c>0$ and arbitrarily large $n$, and he conjectured that this is
essentially optimal \cite{Erdos1946,Erdos1994}. The best upper bound,
\[
 u(n)=O(n^{4/3}),
\]
is due to Spencer, Szemer\'edi and Trotter
\cite{SpencerSzemerediTrotter1984}; Sz\'ekely \cite{Szekely1997} gave a
proof by crossing numbers.

In $2026$ a team at OpenAI \cite{OpenAI2026} disproved Erd\H{o}s's
conjecture. They constructed sets with $u(U)\ge|U|^{1+\delta}$ for a
fixed $\delta>0$ and arbitrarily large $|U|$, replacing the integer
lattice by ideal lattices in number fields of large degree and small
root discriminant. Alon, Bloom, Gowers, Litt, Sawin, Shankar, Tsimerman,
Wang and Wood \cite{Alon2026} simplified the argument and obtained
$\delta\approx6\cdot10^{-38}$. Sawin \cite{Sawin2026} made every step
explicit and sharpened it, and reached the exponent $1.014114$.

Improvements followed within days in online discussions. On the
Erd\H{o}s Problems forum, mlewko \cite{ErdosProblems90} reached
$1.031849$ in Sawin's criterion with finite data found by ChatGPT, and
on MathOverflow spiderduckpig \cite{SpiderduckpigMO2026} adjusted these
data to reach $1.031883$. In a comment the next day, spiderduckpig
proposed a structural change: in the Golod--Shafarevich step of Sawin's
argument \cite[Lemma~11]{Sawin2026}, replace the unweighted count of
relations by a count weighted by degree in the Zassenhaus filtration.
This was reported to give $\delta\ge0.0333487$
\cite{SpiderduckpigMO2026}. Weighted counts of this kind underlie the
towers of \cite{NaslundLegacy2026}, \cite{Naslund0418235} and the
present paper. The author's MathOverflow answer \cite{NaslundMO2026}
then combined this count with covolume counting, dyadic ramification and
Euler products over a fixed subfield, and reported $\delta\ge0.03583$;
the author's archived manuscript \cite{NaslundLegacy2026} records this
construction with $\delta=0.0358324$. The change log of that answer, of
June~9, 2026, reports a bookkeeping issue affecting the covolume step and
gives $\delta>0.0357$ instead. Table~\ref{intro:table-bounds} lists the
explicit exponents, and the repository \cite{TaoConstants84a} records
further bounds posted online.

\begin{table}[htbp]
\centering
\caption{Explicit lower bounds $u(n)\ge n^{1+\delta}$ for arbitrarily
large $n$, in increasing order. The exponents are rounded down and
constant factors are omitted. The last column gives the kind of source;
apart from Erd\H{o}s's, none of these bounds had appeared in a refereed
journal when this paper was written.}
\label{intro:table-bounds}
\small
\begin{tabular}{@{}llll@{}}
\toprule
Exponent & Source & Main new input & Kind of source\\
\midrule
$1+c/\log\log n$ & Erd\H{o}s \cite{Erdos1946} & integer lattice & journal\\
$1+\delta$, $\delta>0$ & OpenAI \cite{OpenAI2026} & ideal lattices from class field towers & manuscript\\
$1+6\cdot10^{-38}$ & Alon et al.\ \cite{Alon2026} & simplified argument & preprint\\
$1.014114$ & Sawin \cite{Sawin2026} & explicit and sharpened criterion & preprint\\
$1.015263$ & Emmerich \cite{Emmerich2026} & optimized parameters & preprint\\
$1.031185$ & Emmerich, Cordella \cite{EmmerichCordella2026} & extended range of primes & certificate\\
$1.031849$ & mlewko \cite{ErdosProblems90} & finite data found by ChatGPT & forum post\\
$1.0333487$ & spiderduckpig \cite{SpiderduckpigMO2026} & weighted Zassenhaus count & online post\\
$1.0357$ & Naslund \cite{NaslundMO2026,NaslundLegacy2026} & covolume counting, dyadic & online post,\\
 & & ramification, fixed-subfield Euler products & manuscript\\
$1.0418235$ & Naslund \cite{Naslund0418235} & mixed signature, norm-one units, & manuscript\\
 & & nonabelian dyadic local group, profiles & \\
$1.04273$ & this paper & tower over $\Q(\sqrt{241})$ & \\
\bottomrule
\end{tabular}
\end{table}

\subsection{The construction in outline}\label{intro:outline}
The construction carries Erd\H{o}s's lattice argument to number fields
of large degree. Let $K$ be such a field, of degree $2d$, with an
involution $\iota$ whose fixed field is $F$. We take the points of a
translated ideal lattice of $K$ that lie in a bounded region of
$K\otimes\R$. A complex embedding in which $\iota$ acts as complex
conjugation maps them injectively to the plane. It sends every
difference of relative norm one over $F$ to a unit vector. The exponent
is then set by a balance, per unit of $d$. The gain is the number of
equal-norm choices: a prime of $F$ that splits in $K$ and is used with
exponent $k$ gives $k+1$ ideals of the same relative norm. Selecting a
common ideal class loses part of this gain. Through the class-number
formula, that loss is governed by the root discriminant $\lambda$ of $K$
and the value at $s=1$ of $L_F(s)=\zeta_K(s)/\zeta_F(s)$. The geometric
window decides how many of the resulting displacements keep both
endpoints in the region, and its volume fixes the number of points.
Using a prime $\mathfrak p$ with exponent $k$ also makes the ideal
lattice denser by the factor $N\mathfrak p^k$, and the exponent $\delta$
charges the gain for the points this adds.

Quantitatively, suppose that such fields exist with $d\to\infty$, a fixed
root discriminant $\lambda$, fixed local types at the selected primes,
and $\theta=c/d$ bounded below, where $(b,c)$ is the signature of $F$ and
$d=b+2c$; thus $2\theta$ is the proportion of non-real embeddings of $F$.
With hard windows at the selected primes, the exponent $1+\delta$ is
attained when
\begin{equation}\label{intro:margin}
 \sum_q\frac{\log(k_q+1)-\delta k_qf_q\log q}{e_qf_q}
 -\Bigl(\frac12-\delta\Bigr)\log\lambda-C+\mathcal A_\delta(\theta)>0.
\end{equation}
Here $q$ runs over the rational primes below the selected primes, $e_q$
and $f_q$ are their absolute ramification indices and residue degrees in
$F$, $k_q$ is the exponent used at the primes above $q$, $C$ is an upper
bound for $d^{-1}\log L_F(1)$, and
\[
 \mathcal A_\delta(\theta)=(1-\delta)\log2+(1-\theta)\log\pi
 +(1-2\theta)J_\R+\theta J_\C
\]
is the archimedean contribution, in which $J_\R$ and $J_\C$ measure the
overlap of the archimedean profiles at the real and at the complex
places of $F$ under unit displacements (Proposition~\ref{geo:transfer}).
The term $\theta J_\C$ comes from mixed signature: above a complex place
of $F$, the two coordinates of a norm-one element have reciprocal moduli.
Averaging over the norm-one units counts the resulting one-parameter
family of displacements exactly, and the class-number formula cancels the
regulator of these units (Proposition~\ref{geo:mass}).
Section~\ref{fw:section} replaces the hard windows by weighted shells
(Proposition~\ref{fw:transfer}), and Section~\ref{sec:certificate}
verifies the resulting inequality at $\delta=\deltastar$.

The fields must form an infinite family of growing degree with fixed
$\lambda$ and fixed local types. Such families come from infinite
pro-$2$ class field towers with prescribed local Galois groups. The
Golod--Shafarevich criterion \cite{GolodShafarevich} proves that a tower
is infinite when a certain power series in a variable $t$ takes a
negative value; the use of Frobenius conditions to control splitting in
such towers goes back to Hajir, Maire and Ramakrishna
\cite{HajirMaireRamakrishna2021Cutting}.

\subsection{A tower over $\Q(\sqrt{241})$}\label{intro:new}
The fields of \cite{Naslund0418235} come from a pro-$2$ tower over $\Q$
that is unramified outside $2,3,5,7,11,13$, with prescribed local Galois
groups at these primes and Frobenius conditions at the selected primes
up to $31$. Its Golod--Shafarevich series has the form
\[
 1-7t+(\text{a cost for each local condition})-s_D(t).
\]
Here $7$ is the number of generators of the Galois group of the maximal
pro-$2$ extension of $\Q$ unramified outside these six primes, and
$s_D(t)$ is a saving: by Hilbert reciprocity, one dyadic relation is
implied by the others and need not be paid for. Tame ramification at $7$,
$11$ and $13$ supplies generators to pay for the local conditions, and it
costs a factor $\sqrt{1001}$ in the root discriminant.

Now let $B$ be a totally real field of degree $m$ in which every prime
used by the construction splits completely. The generators and the local
conditions of a construction over $\Q$ are then repeated $m$ times over
$B$, while the constant term $1$ and the saving $s_D$ are not; if $P(t)$
is the series over $\Q$, the series over $B$ is roughly
$mP(t)-(m-1)(1-s_D(t))$. The extra room can be spent on ramifying at
fewer primes. The price is the discriminant of $B$, which enters the
root discriminant of the fields. The field $B=\Q(\sqrt{241})$ is the
real quadratic field of smallest discriminant in which $2$, $3$ and $5$
all split, and over it the tower needs no ramification at $7$, $11$ or
$13$. Section~\ref{tw:tower} gives the exact series.

Over $B$ we prescribe, at each of the two primes above $2$, the same
nonabelian dyadic local Galois group of order $32$ as in
\cite{Naslund0418235}; at each of the four primes above $3$ and $5$, a
tame local group $C_2\times C_2$; and conditions that make the Frobenius
elements have order dividing $4$ at the two primes above $29$ and at the
prime $7$, which is inert in $B$. The tower is ramified only above $2$,
$3$ and $5$. Its Galois group has $8$ generators, and the root
discriminant of its fields is
\[
 \lambda=2^{9/4}\sqrt{3615}\approx286.0,
 \qquad 3615=3\cdot5\cdot241,
\]
against $2^{9/4}\sqrt{15015}\approx582.9$ over $\Q$: the base contributes
$\sqrt{241}$, and $\sqrt{1001}$ is saved. The price is paid in the
selected primes: only $2$, $3$, $5$, $7$ and $29$ remain, against the
eleven primes up to $31$ over $\Q$. Theorem~\ref{tw:field-family} states
the resulting family of fields.

The analytic estimate changes accordingly. In \cite{Naslund0418235} the
bound $C$ for $d^{-1}\log L_F(1)$ came from a complete Hecke-factor
calculation in a nonabelian subfield of degree $16384$ of the tower.
Here every field $K$ of Theorem~\ref{tw:field-family} contains the
Kummer field $E=B(\sqrt V)$, of degree $512$ over $\Q$, where $V$ is the
group of $\{2,3,5\}$-units of $B$ modulo squares. Its Dedekind zeta
function is the product of $\zeta_B$ and $255$ quadratic Hecke
$L$-functions of $B$, which we evaluate with certified approximate
functional equations (Sections~\ref{gf:section} and~\ref{an:section}).
Together with refinements from the exact local types of the tower, this
gives $C=\ceiling$ (Theorem~\ref{an:ceiling}).
Table~\ref{intro:table-constructions} compares the two constructions.

\begin{table}[htbp]
\centering
\caption{The construction of \cite{Naslund0418235} and the present one.
The pairs $(e,f)$ are absolute ramification indices and residue degrees
of the selected primes of $F$.}
\label{intro:table-constructions}
\begin{tabular}{lll}
\toprule
 & \cite{Naslund0418235} & this paper\\
\midrule
base field & $\Q$ & $\Q(\sqrt{241})$\\
primes ramified in the tower & $2,3,5,7,11,13$ & above $2,3,5$\\
selected primes, $(e,f)$ & $2{:}\,(8,4)$;\ $3,5{:}\,(2,2)$; & $2{:}\,(8,4)$;\ $3,5{:}\,(2,2)$;\\
 & $7,11,13{:}\,(2,4)$;\ $17,\dots,31{:}\,(1,4)$ & $7{:}\,(1,8)$;\ $29{:}\,(1,4)$\\
root discriminant $\lambda$ & $2^{9/4}\sqrt{15015}\approx582.9$ & $2^{9/4}\sqrt{3615}\approx286.0$\\
$\ell=\log\lambda$ & $6.3680$ & $5.6560$\\
lower bound $\theta_*$ for $\theta$ & $4095/8192$ & $65535/131072$\\
ceiling $C$ & $0.042161819$ & $0.04871285$\\
exponent $1+\delta$ & $1.0418235$ & $1.04273$\\
certified margin & $5.3\cdot10^{-6}$ & $1.77\cdot10^{-4}$\\
\bottomrule
\end{tabular}
\end{table}

\subsection{Why the exponent improves}\label{intro:why}
In \eqref{intro:margin}, lowering $\log\lambda$ from $6.3680$ to
$5.6560$ is worth $(1/2-\delta)\cdot0.712\approx0.33$ per unit of $d$ at
$\delta\approx0.043$. Three terms pay for it. The first sum, the gain
from equal-norm choices net of their cost in vertices, falls from $1.034$
to $0.721$, because the tower over $B$ selects only five rational primes,
and two of them, $7$ and $29$, with large residue degrees. The ceiling
$C$ rises by $0.0066$, because the Kummer field carries less Frobenius
information than the field of degree $16384$ used over $\Q$. The
archimedean term falls by $0.013$, mostly because $\delta$ is larger.
Table~\ref{intro:table-budget} shows the certified terms of the margin
for both constructions, each at its own exponent.

\begin{table}[htbp]
\centering
\caption{The terms of the certified margin, per unit of $d$, rounded to
seven decimals. The first row is the weighted-shell form of the first sum
in \eqref{intro:margin} (Section~\ref{fw:section}), and the archimedean
term is computed with the certified lower bound $J_\C^*$ for $J_\C$. The
margin is the quantity $M_*(\theta_*)$ of Section~\ref{sec:certificate}
and of the certificate of \cite{Naslund0418235}.}
\label{intro:table-budget}
\begin{tabular}{lrr}
\toprule
Term & \cite{Naslund0418235}, $\delta=0.0418235$ & this paper, $\delta=\deltastar$\\
\midrule
finite sum & $1.0335669$ & $0.7214955$\\
$-(1/2-\delta)\log\lambda$ & $-2.9176605$ & $-2.5863213$\\
$-C$ & $-0.0421618$ & $-0.0487129$\\
$\mathcal A_\delta(\theta_*)$ & $1.9262607$ & $1.9137154$\\
\midrule
margin & $0.0000053$ & $0.0001767$\\
\bottomrule
\end{tabular}
\end{table}

The balance leaves a certified margin of about $1.77\cdot10^{-4}$ per unit
of $d$ at $\delta=\deltastar$. With the same profiles, and shell weights
optimized in the same way, the margin at $\delta=0.0428$ is negative.

\subsection{Relation to Sawin's criterion}\label{intro:comparison}
All bounds in Table~\ref{intro:table-bounds} from Sawin's onward use his
framework \cite{Sawin2026}: a quadratic extension $K/F$, ideals of equal
relative norm built from split primes, a pigeonhole over ideal classes,
and fields from Golod--Shafarevich towers with Frobenius conditions.
Section~1.6 of \cite{Naslund0418235} compares that criterion with
\eqref{intro:margin} term by term. In brief, counting vertices by
covolume instead of packing returns $\delta\log\lambda$. The exact
class-number formula, with an explicit Euler-product bound $C$, replaces
Louboutin's bound for the relative class number \cite{Louboutin2000}
together with the factor $2^d$ that Sawin allows for units; their
combined cost $1+\log\log\lambda$ would be about $2.73$ here. Mixed
signature contributes the term $\theta(J_\C-2J_\R-\log\pi)$, worth more
than $0.31$ per unit of $d$ here. Optimized profiles and shells replace
the sup-norm window. The present paper keeps all of these and changes the
tower. Sawin suggested that towers satisfying different local conditions
might improve the exponent \cite{Sawin2026}. The base field is a change
of a different kind: it adds no local conditions at new primes, but
repeats the conditions at the small primes over two primes of $B$ each.

\subsection{Organization}\label{intro:organization}
Section~\ref{tw:tower} constructs the tower over $B$ and proves
Theorem~\ref{tw:field-family}, which supplies the fields.
Section~\ref{gf:section} describes the Kummer field and its quadratic
$L$-functions, and Section~\ref{an:section} proves the ceiling
$d^{-1}\log L_F(1)<C$ (Theorem~\ref{an:ceiling}).
Sections~\ref{geo:section}--\ref{pf:section} carry out the geometric
construction: the norm-one mass formula and the transfer to planar sets,
the weighted windows at the selected primes, and the archimedean
profiles. These sections hold for any quadratic extension with the
stated properties and follow \cite{Naslund0418235} closely.
Section~\ref{sec:certificate} evaluates the final inequality at
$\delta=\deltastar$ and completes the proof of Theorem~\ref{thm:main}.

\subsection{Computations and data}\label{intro:code}
The proof uses finite computations: linear algebra over $\Ftwo$ for the
tower, exact rational arithmetic for the Golod--Shafarevich series,
certified evaluations of $L$-functions in the ball arithmetic of Arb,
now part of FLINT \cite{Johansson2017,FLINT}, and directed interval
bounds for the geometric margin \cite{Mpmath}. PARI/GP \cite{PARI} is
used for the arithmetic of number fields and for independent checks.
Each finite fact is verified by a supplementary program, named in the
text where the fact is used. These certificates concern exact arithmetic
and rigorous inequalities, not the statistical accuracy of
floating-point approximations.

The programs and their data are in the directory
\texttt{papers/0.04273/certificates} of the public repository
\cite{NaslundCode2026}. Some certified routines, for the degree-two
approximate functional equations, the archimedean profiles and the
finite windows, are imported unchanged from the supplementary archive of
\cite{Naslund0418235}, which is in the same repository together with a
manifest of the SHA-256 hashes of its members. A replay program runs
every step and records its results. The computations were run with
PARI/GP~2.17.2, python-flint~0.9.0 (FLINT~3.6.0) and mpmath~1.3.0.

A separate Lean~4 development \cite{NaslundLean2026}, built on Mathlib
\cite{Lean4_2021,Mathlib2020}, proves the planar conclusion with the
slightly smaller exponent $1.0427$ from one explicit numerical inequality
for the Dedekind zeta function of $E$
(Remark~\ref{an:lean-hypothesis}). That inequality has not been proved in
Lean, so the formalization is conditional, and it is not used as a
substitute for the arguments here.

\tableofcontents
\section{\texorpdfstring{The Tower over $\Q(\sqrt{241})$}{The Tower over Q(sqrt 241)}}\label{tw:tower}

This section constructs the number fields of Theorem~\ref{tw:field-family}
as finite layers of an infinite pro-$2$ extension of $B=\Q(\sqrt{241})$
with prescribed local Galois groups. The construction adapts the tower over
$\Q$ of \cite[Section~2]{Naslund0418235}, and we repeat its arguments in the
form needed over $B$. Subsection~\ref{tw:base-section} describes $B$ and its
$\{2,3,5\}$-units. Subsection~\ref{tw:local-section} recalls the local Galois
groups involved and constructs the dyadic local field.
Subsection~\ref{tw:global-section} shows that the Galois group
$\mathcal G_B$ of the maximal pro-$2$ extension of $B$ unramified outside
$2,3,5$ has a presentation with eight generators and seven local relators.
Subsection~\ref{tw:prescribed-section} imposes the local conditions and
defines the quotient $G_B$. Subsection~\ref{tw:retained-section} computes the
first three layers of the Zassenhaus filtration of $G_B$: they retain the
prescribed local groups, and complex conjugation has many conjugates.
Subsections~\ref{tw:fox-section} and~\ref{tw:infinite-section} prove that
$G_B$ is infinite, by the filtered Fox calculus of \cite{Naslund0418235} and
the Golod--Shafarevich method \cite{GolodShafarevich}.
Subsection~\ref{tw:fields-section} proves Theorem~\ref{tw:field-family}.

We use the following conventions. Subgroups of profinite groups are closed,
generation is topological, and the normal closure of a set is the smallest
closed normal subgroup containing it. The commutator is
$[g,h]=g^{-1}h^{-1}gh$. For a finite $2$-group $\Delta$ with augmentation
ideal $\mathfrak I\subset\Ftwo[\Delta]$, the Zassenhaus filtration is
$D_n\Delta=\{g:g-1\in\mathfrak I^n\}$; for a pro-$2$ group it is defined
through finite quotients. Then $[D_m,D_n]\subset D_{m+n}$, $g^2\in D_{2n}$
for $g\in D_n$, a surjection maps $D_n$ onto $D_n$, and
$\operatorname{gr}\Delta=\bigoplus_{n\ge1}D_n\Delta/D_{n+1}\Delta$ is a
restricted Lie algebra over $\Ftwo$ whose bracket is induced by commutators
and whose restricted square $v\mapsto v^{[2]}$ is induced by squaring
\cite{Jennings,Quillen,Ershov2012}. By Lazard's formula
\cite[Proposition~3.2]{Ershov2012}, $D_n=\prod_{i2^j\ge n}\gamma_i^{2^j}$,
where $\gamma_i$ is the lower central series. In particular $D_1$ is the
whole group, $D_2$ is the Frattini subgroup, and $D_3=D_2^{\,2}[D_2,D_1]$. For
a free pro-$2$ group $\mathcal F$ of rank $n$, $\operatorname{gr}\mathcal F$
is the free restricted Lie algebra on $\mathcal F/D_2\mathcal F\cong\Ftwo^n$
\cite{Quillen,Ershov2012}. We realize it inside the free associative
algebra $\Ftwo\langle X_1,\dots,X_n\rangle$, with bracket $[a,b]=ab+ba$ and
restricted square $a^{[2]}=a^2$; its layers of degrees one, two and three
have dimensions $n$, $n(n+1)/2$ and $(n^3-n)/3$.

\subsection{The Base Field}\label{tw:base-section}
Let $B=\Q(\sqrt{241})$, with real places $v_1$ ($\sqrt{241}\mapsto+\sqrt{241}$)
and $v_2$ ($\sqrt{241}\mapsto-\sqrt{241}$). Put
\begin{equation}\label{tw:alpha}
\begin{gathered}
 \alpha_0=-1,\qquad \alpha_1=\varepsilon=-71011068+4574225\sqrt{241},\\
 \alpha_2=\tfrac{-6101-393\sqrt{241}}2,\qquad
 \alpha_3=\tfrac{6101-393\sqrt{241}}2,\\
 \alpha_4=31-2\sqrt{241},\quad \alpha_5=31+2\sqrt{241},\quad
 \alpha_6=326-21\sqrt{241},\quad \alpha_7=326+21\sqrt{241}.
\end{gathered}
\end{equation}
Their norms $N_{B/\Q}(\alpha_i)$ are $1,-1,-2,-2,-3,-3,-5,-5$, and
$\alpha_2\alpha_3=2$, $\alpha_4\alpha_5=-3$, $\alpha_6\alpha_7=-5$. For a
nonzero ideal $\mathfrak a$ of $\mathcal O_B$ we write
$N\mathfrak a=|\mathcal O_B/\mathfrak a|$ for its absolute norm.

\begin{lemma}\label{tw:base}
\begin{enumerate}
\item[(a)] $\mathcal O_B=\Z[(1+\sqrt{241})/2]$, the discriminant of $B$ is
 $241$, and $B$ has class number one.
\item[(b)] The unit $\varepsilon$ has norm $N_{B/\Q}(\varepsilon)=-1$, and the
 classes of $-1$ and $\varepsilon$ form a basis of
 $\mathcal O_B^\times/\mathcal O_B^{\times2}$.
\item[(c)] In $\mathcal O_B$ we have $2=\mathfrak p_1\mathfrak p_2$,
 $3=\mathfrak q_1\mathfrak q_2$, $5=\mathfrak r_1\mathfrak r_2$ and
 $29=\mathfrak t_1\mathfrak t_2$, where
 \[
  \mathfrak p_j=\alpha_{1+j}\mathcal O_B,\quad
  \mathfrak q_j=\alpha_{3+j}\mathcal O_B,\quad
  \mathfrak r_j=\alpha_{5+j}\mathcal O_B,\quad
  \mathfrak t_j=(-14127\pm910\sqrt{241})\mathcal O_B,
 \]
 with the sign $+$ for $j=1$. These eight primes are distinct and of degree
 one. The prime $7$ is inert, and we write $\mathfrak t_0=7\mathcal O_B$, so
 that $N\mathfrak t_0=49$. The prime $241$ ramifies. At each of the eight
 primes of degree one, above $p\in\{2,3,5,29\}$, the completion of $B$ is
 $\Q_p$, and $\sqrt{241}$ maps to the square root of $241$ in $\Z_p$ given in
 Table~\ref{tw:squareclass-table}.
\item[(d)] Let $S=\{\mathfrak p_1,\mathfrak p_2,\mathfrak q_1,\mathfrak q_2,\mathfrak r_1,\mathfrak r_2\}$
 and $V=\mathcal O_B[1/30]^\times/\mathcal O_B[1/30]^{\times2}$. The classes
 of $\alpha_0,\dots,\alpha_7$ form a basis of $V$, and $V$ is the subgroup of
 $B^\times/B^{\times2}$ of classes with even valuation at every prime outside
 $S$. Moreover $\Pic(\mathcal O_B[1/30])=0$.
\end{enumerate}
\end{lemma}

\begin{proof}
(a) Since $241\equiv1\pmod 4$, the ring of integers and the discriminant are
as stated. The Minkowski bound is $\sqrt{241}/2<8$, so every ideal class
contains an integral ideal of norm at most $7$, which is a product of primes
of norm at most $7$. By (c), these are the six primes above $2,3,5$, and they
are principal.

(b) A direct computation gives $71011068^2-241\cdot4574225^2=-1$. By
Dirichlet's theorem $\mathcal O_B^\times=\{\pm1\}\times\eta^{\Z}$ for a
fundamental unit $\eta$, and $\varepsilon=\pm\eta^k$. Since
$N_{B/\Q}(\varepsilon)=-1$, the exponent $k$ is odd, so $-1$ and
$\varepsilon$ generate $\mathcal O_B^\times$ modulo squares. Neither $-1$ nor
$\pm\varepsilon$ is a square, because $B$ is real and
$N_{B/\Q}(\pm\varepsilon)=-1$.

(c) We have $241\equiv1\pmod8$, $241\equiv1\pmod3$, $241\equiv1\pmod5$,
$241\equiv9\pmod{29}$ and $241\equiv3\pmod7$, and $3$ is not a square modulo
$7$. Hence $2,3,5,29$ split, $7$ is inert and $241$ ramifies. The listed
generators have norms $\pm2,\pm3,\pm5,29$, so they generate primes of degree
one, and the displayed products show that the two generators above each $p$
generate the two distinct primes above $p$; for $29$ the product of the two
generators is $29$. A prime of degree one above $p$ is the preimage of
$p\Z_p$ under an embedding $B\to\Q_p$, which sends $\sqrt{241}$ to one of
the two square roots of $241$ in $\Z_p$; the completion there is $\Q_p$, and
the prime contains the stated generator exactly for the square root in
Table~\ref{tw:squareclass-table}.

(d) Every $S$-unit $u$ satisfies
$u\mathcal O_B=\prod_{i\ge2}(\alpha_i\mathcal O_B)^{a_i}$, so
$u\prod_{i\ge2}\alpha_i^{-a_i}$ is a unit, and by (b) the classes of
$\alpha_0,\dots,\alpha_7$ generate $V$. If $\prod_i\alpha_i^{a_i}$ with
$a_i\in\{0,1\}$ is a square, the valuations at $S$ give $a_i=0$ for $i\ge2$,
and then (b) gives $a_0=a_1=0$. If $b\in B^\times$ has even valuation outside
$S$, then
$b\mathcal O_B=\mathfrak a^2\prod_{\mathfrak p\in S}\mathfrak p^{a_{\mathfrak p}}$
with $\mathfrak a$ principal by (a), so $b$ is an $S$-unit times a square.
Finally $\Pic(\mathcal O_B[1/30])$ is a quotient of $\Pic(\mathcal O_B)=0$.
\end{proof}

\begin{table}[htbp]
\centering
\caption{The image of $\sqrt{241}$ in $\Z_p$ at each prime of degree one,
given by its residue, and in $\R$ at the real places; the classes of
$\alpha_0,\dots,\alpha_7$ in $\Q_p^\times/\Q_p^{\times2}$ at these primes, and
their signs at the real places. At $p=2$ the classes are written as
$\pm1,\pm5,\pm2,\pm10$; at odd $p$ as $1,u,p,up$, where $u=-1,2,2$ is a
nonsquare unit for $p=3,5,29$.}
\label{tw:squareclass-table}
\begin{tabular}{llrrrrrrrr}
\toprule
place & $\sqrt{241}\mapsto$ & $\alpha_0$ & $\alpha_1$ & $\alpha_2$ & $\alpha_3$ & $\alpha_4$ & $\alpha_5$ & $\alpha_6$ & $\alpha_7$\\
\midrule
$\mathfrak p_1$ & $7\bmod 32$ & $-1$ & $-5$ & $-10$ & $-5$ & $1$ & $5$ & $-5$ & $1$\\
$\mathfrak p_2$ & $25\bmod 32$ & $-1$ & $5$ & $5$ & $10$ & $5$ & $1$ & $1$ & $-5$\\
$\mathfrak q_1$ & $2\bmod 3$ & $-1$ & $1$ & $-1$ & $1$ & $3$ & $-1$ & $-1$ & $-1$\\
$\mathfrak q_2$ & $1\bmod 3$ & $-1$ & $-1$ & $-1$ & $1$ & $-1$ & $3$ & $-1$ & $-1$\\
$\mathfrak r_1$ & $1\bmod 5$ & $1$ & $2$ & $2$ & $1$ & $1$ & $2$ & $10$ & $2$\\
$\mathfrak r_2$ & $4\bmod 5$ & $1$ & $2$ & $1$ & $2$ & $2$ & $1$ & $2$ & $10$\\
$\mathfrak t_1$ & $3\bmod 29$ & $1$ & $1$ & $2$ & $1$ & $1$ & $2$ & $2$ & $2$\\
$\mathfrak t_2$ & $26\bmod 29$ & $1$ & $1$ & $1$ & $2$ & $2$ & $1$ & $2$ & $2$\\
$v_1$ & $+\sqrt{241}$ & $-$ & $+$ & $-$ & $-$ & $-$ & $+$ & $-$ & $+$\\
$v_2$ & $-\sqrt{241}$ & $-$ & $-$ & $+$ & $+$ & $+$ & $-$ & $+$ & $-$\\
\bottomrule
\end{tabular}
\end{table}

The classes in Table~\ref{tw:squareclass-table} follow from the residues of
$\sqrt{241}$ given there; at $p=2$ they need its image modulo $32$. For
example, at $\mathfrak p_1$ we have $4574225\equiv1$ and $-71011068\equiv4$
modulo $8$, so $\varepsilon\equiv4+7\equiv3\equiv-5\pmod 8$. At the inert
prime $\mathfrak t_0$, an element $a\in\mathcal O_B$ prime to $\mathfrak t_0$
is a square modulo $\mathfrak t_0$ exactly when $a^{24}\equiv1$, that is, when
the reduction of $N_{B/\Q}(a)\equiv a^8$ is a square in $\mathbb F_7$. Hence
$\alpha_i$ is a square modulo $\mathfrak t_0$ exactly for $i=0,4,5,6,7$.

\subsection{Local Galois Groups}\label{tw:local-section}
For a place $\nu$ of $B$ let $\mathcal G_\nu$ be the Galois group of the
maximal pro-$2$ extension of the completion $B_\nu$; for a real place,
$\mathcal G_\nu=\Gal(\C/\R)$. We write $(\cdot,\cdot)$ for the Hilbert symbol
of order two of a local field \cite[Chapter~III, Section~4]{MilneCFT}. For
odd $p$ and $p$-adic units $u,w$, and for $2$-adic units $u,w$,
\begin{equation}\label{tw:hilbert}
 (p^au,p^bw)_p=(-1)^{ab(p-1)/2}\Bigl(\frac up\Bigr)^b\Bigl(\frac wp\Bigr)^a,
 \qquad
 (2^au,2^bw)_2=(-1)^{\frac{u-1}2\frac{w-1}2+a\frac{w^2-1}8+b\frac{u^2-1}8}
\end{equation}
\cite[Chapter~III, Theorem~1]{SerreCourse}, \cite[Chapter~VIII,
Section~5]{MilneCFT}. A defining relator of a pro-$2$ group with respect to
generators $g_1,\dots,g_k$ is an element of the free pro-$2$ group on $k$
letters whose normal closure is the kernel of the map to the group; its
quadratic initial is its class in the degree-two layer of the Zassenhaus
filtration of the free group.

\begin{lemma}\label{tw:local-groups}
\begin{enumerate}
\item[(a)] If $\nu$ is real, then $\mathcal G_\nu=\langle\iota\rangle\cong C_2$
 with the single defining relator $\iota^2$, and $\iota(\sqrt b)=-\sqrt b$
 exactly when $b<0$.
\item[(b)] Let $p$ be an odd prime and $\varpi$ a uniformizer of $\Q_p$. Then
 $\mathcal G_{\Q_p}$ is generated by a generator $\tau$ of its inertia group
 and a Frobenius lift $\varphi$ with $\varphi(\sqrt\varpi)=\sqrt\varpi$, with
 the single defining relator $\varphi\tau\varphi^{-1}\tau^{-p}$. For
 $b\in\Q_p^\times$, $\tau(\sqrt b)=-\sqrt b$ exactly when $v_p(b)$ is odd;
 for a unit $b$, $\varphi(\sqrt b)=-\sqrt b$ exactly when $b$ is not a square
 modulo $p$. The relator has quadratic initial
 $[\tau,\varphi]+\frac{p-1}2\tau^{[2]}$.
\item[(c)] Let $E_2=\Q_2(\sqrt{-1},\sqrt2,\sqrt5)$, and let
 $x,y,z\in\mathcal G_{\Q_2}$ multiply the triple $(\sqrt2,\sqrt{-1},\sqrt5)$
 by the signs $(-,+,+)$, by $(+,-,+)$ and by $(+,-,-)$. Then $x,y,z$
 generate $\mathcal G_{\Q_2}$. If $b\equiv(-1)^{a_1}2^{a_2}5^{a_3}$ modulo
 squares, they multiply $\sqrt b$ by $(-1)^{a_2}$, $(-1)^{a_1}$ and
 $(-1)^{a_1+a_3}$, that is, by the Hilbert symbols $(5,b)_2$, $(-1,b)_2$ and
 $(-2,b)_2$. The group $\mathcal G_{\Q_2}$ has a single defining relator $r$
 with respect to $x,y,z$, and every such $r$ has quadratic initial
 $y^{[2]}+[x,y]+[x,z]$.
\end{enumerate}
\end{lemma}

\begin{proof}
(a) is clear.

(b) By Iwasawa's theorem \cite[Theorem~7.5.3]{NeukirchSchmidtWingberg2008},
the Galois group of the maximal tamely ramified extension of $\Q_p$ is
generated by a Frobenius lift $\varphi'$ and a generator $\tau$ of tame
inertia, with the single relation $\varphi'\tau\varphi'^{-1}=\tau^p$. Since
$p$ is odd,
every $2$-extension of $\Q_p$ is tamely ramified, so $\mathcal G_{\Q_p}$ is
the pro-$2$ group with the same presentation. The extension
$\Q_p(\sqrt\varpi)$ is ramified, so $\tau$ moves $\sqrt\varpi$; replacing
$\varphi'$ by $\varphi'\tau$ if necessary gives $\varphi$. This substitution
does not change the relator, since
$(\varphi'\tau)\tau(\varphi'\tau)^{-1}=\varphi'\tau\varphi'^{-1}$.
The actions on square roots hold because $\Q_p(\sqrt b)$ is ramified exactly
when $v_p(b)$ is odd, and because $\varphi$ acts on unramified extensions as
the Frobenius. Modulo $D_3$, the element $\varphi\tau\varphi^{-1}\tau^{-1}$
has class $[\varphi,\tau]=[\tau,\varphi]$, and $\tau^{1-p}$ has class
$\frac{p-1}2\tau^{[2]}$.

(c) By Kummer theory the maximal elementary abelian quotient of
$\mathcal G_{\Q_2}$ is $\Gal(E_2/\Q_2)$, since $\Q_2^\times/\Q_2^{\times2}$
has basis $-1,2,5$. The images of $x,y,z$ form a basis of it, so $x,y,z$
generate \cite[Proposition~3.9.1]{NeukirchSchmidtWingberg2008}. The action on
$\sqrt b$ is read off from the actions on $\sqrt{-1},\sqrt2,\sqrt5$, and it
agrees with the stated Hilbert symbols by \eqref{tw:hilbert}. By
\cite[Theorem~7.5.11]{NeukirchSchmidtWingberg2008}, $\mathcal G_{\Q_2}$ is a
Demu\v skin group of rank $3$: it has a single defining relator $r$, and
$H^2(\mathcal G_{\Q_2},\Ftwo)\cong\Ftwo$. For $b\in\Q_2^\times$ let $\chi_b$
be the character with $g(\sqrt b)=(-1)^{\chi_b(g)}\sqrt b$. The characters
dual to $x,y,z$ are $\chi_2,\chi_{-5},\chi_5$. Inflation identifies
$H^2(\mathcal G_{\Q_2},\Ftwo)$ with $H^2$ of the absolute Galois group of
$\Q_2$ \cite[Proposition~7.5.8]{NeukirchSchmidtWingberg2008}, and there the
invariant of $\chi_a\cup\chi_b$ is given by the Hilbert symbol $(a,b)_2$
\cite[Chapter~III, Section~4]{MilneCFT}. By \eqref{tw:hilbert}, the matrix of
the indicator of $(a,b)_2=-1$ on $2,-5,5$ is
\[
 \begin{pmatrix}0&1&1\\1&1&0\\1&0&0\end{pmatrix}.
\]
The evaluation $\operatorname{tr}_r$ of the transgression at $r$ is an
isomorphism $H^2(\mathcal G_{\Q_2},\Ftwo)\to\Ftwo$
\cite[Proposition~3.9.12(iii)]{NeukirchSchmidtWingberg2008}, and so is the
invariant, so the two coincide. We now apply the relation--cup identity
\cite[Proposition~3.9.13(ii)]{NeukirchSchmidtWingberg2008}, see also
\cite[Proposition~3.2]{Quadrelli2021}. Write $g_1,g_2,g_3$ for $x,y,z$. The
identity writes the relator as
$r=\prod_jg_j^{qa_j}\prod_{k<l}[g_k,g_l]^{a_{kl}}\,r'$, with $r'$ in the
third term of the descending $q$-central series, and states that
$\operatorname{tr}_r$ of the cup product of the characters dual to $g_k$ and
$g_l$ is $-a_{kl}$ for $k<l$ and $-\binom q2a_k$ for $k=l$. Here $q$, the
smallest elementary divisor of the
abelianization of $\mathcal G_{\Q_2}$, is $2$: by local class field theory
this abelianization is the pro-$2$ completion of $\Q_2^\times$, which is
$\Z_2^2\times\Z/2$. The third term of the $2$-central series is
$(D_2)^2[D_2,D_1]=D_3$ by Lazard's formula, the class of $g_j^2$ in
$\operatorname{gr}_2$ is $g_j^{[2]}$, and the signs disappear modulo $2$,
with $\binom22=1$. Hence the off-diagonal entries of the matrix are the
coefficients of $[x,y],[x,z],[y,z]$ in the quadratic initial of $r$, and the
diagonal entries, which are the values of the cup squares $\chi\cup\chi$,
are the coefficients of $x^{[2]},y^{[2]},z^{[2]}$. For $p=2$ the cup square
of a character is its Bockstein, so the diagonal coefficients are also given
by \cite[Proposition~3.9.14]{NeukirchSchmidtWingberg2008}.
\end{proof}

We now construct the dyadic local field. It is the field of
\cite[Lemma~2.2]{Naslund0418235}, and we repeat the proof. Let $U_4$ be the
unramified extension of $\Q_2$ of degree four, with Frobenius $\mathrm{Fr}$;
let $i=\sqrt{-1}$, $k=\Q_2(i)$, $\alpha=1+2i$, and
\[
 M_8=\Q_2(i,\sqrt5,\sqrt\alpha),\qquad M_{16}=M_8(\sqrt2),\qquad
 L_2=M_{16}U_4 .
\]
Since $\alpha\bar\alpha=5$ with $\bar\alpha=1-2i$, the field $M_8$ contains
$\sqrt{\bar\alpha}=\sqrt5/\sqrt\alpha$. Let $\sigma_y$ and $\sigma_z$ be the
automorphisms of $M_8$ given by
\[
 \sigma_y:\ i\mapsto-i,\ \sqrt5\mapsto\sqrt5,\ \sqrt\alpha\mapsto\sqrt5/\sqrt\alpha;
 \qquad
 \sigma_z:\ i\mapsto-i,\ \sqrt5\mapsto-\sqrt5,\ \sqrt\alpha\mapsto\sqrt5/\sqrt\alpha.
\]

\begin{lemma}\label{tw:local-two}
\begin{enumerate}
\item[(a)] The extension $M_8/\Q_2$ is Galois of degree $8$, and its Galois
 group $\langle\sigma_y,\sigma_z\rangle$ is dihedral.
\item[(b)] The extension $L_2/\Q_2$ is Galois of degree $32$, with
 ramification index $8$ and residue degree $4$. There are unique
 $x,y,z\in\Gal(L_2/\Q_2)$ such that $x$ is trivial on $M_8$ and $U_4$ and
 negates $\sqrt2$; $y$ restricts to $\sigma_y$ on $M_8$, fixes $\sqrt2$ and
 is trivial on $U_4$; and $z$ restricts to $\sigma_z$ on $M_8$, fixes $\sqrt2$
 and restricts to $\mathrm{Fr}$ on $U_4$. They act on $E_2\subset L_2$ as in
 Lemma~\ref{tw:local-groups}(c), and the map sending the generators of
 \[
  D=\langle x,y,z\mid x^2,\ y^2,\ z^4,\ [y,z]^2,\ [x,y],\ [x,z],\ [[y,z],z]\rangle
 \]
 to $x,y,z$ is an isomorphism $D\cong\Gal(L_2/\Q_2)$. The inertia group is
 $\langle x,y,[y,z]\rangle\cong C_2^3$.
\item[(c)] The Zassenhaus layers of $D$ are
 $\operatorname{gr}_1D=\langle\bar x,\bar y,\bar z\rangle\cong\Ftwo^3$ and
 $\operatorname{gr}_2D=\langle\overline{z^2},\overline{[y,z]}\rangle\cong\Ftwo^2$,
 and $D_3(D)=1$, so the augmentation Hilbert polynomial of $\Ftwo[D]$ is
 $P_D(t)=(1+t)^3(1+t^2)^2$.
\item[(d)] The different of $L_2/\Q_2$ has valuation $18$ in $L_2$;
 equivalently, its normalized exponent is $18/8=9/4$.
\end{enumerate}
\end{lemma}

\begin{proof}
\emph{Quadratic extensions of $k$.} Put $\varpi=1+i$, a uniformizer of $k$,
so that $v_k(2)=2$, and for $b\in k^\times$ let $f(b)$ be the exponent of the
conductor of $k(\sqrt b)/k$, which for a quadratic extension equals the
exponent of its discriminant. Then $f(5)=0$, since $\Q_2(\sqrt5)/\Q_2$ is
unramified. Since $\alpha=1+\varpi^2$, the element
$\varpi_\alpha=(\sqrt\alpha-1-\varpi)/\varpi$ is a root of the Eisenstein
polynomial $X^2+\frac{2(1+\varpi)}\varpi X+\frac2\varpi$; its derivative at
$\varpi_\alpha$ is the sum of terms of valuations $5$ and $2$ in
$k(\sqrt\alpha)$, so $f(\alpha)=2$. Next $2\alpha=\varpi^2\alpha'$ with
$\alpha'=2-i$, and $\sqrt{\alpha'}-1$ is a root of $X^2+2X+1-\alpha'$, which
is Eisenstein since $v_k(1-\alpha')=1$; its derivative $2\sqrt{\alpha'}$ has
valuation $4$, so $f(2\alpha)=4$. Similarly $2=-i\varpi^2$, $-i\equiv i$
modulo squares, and $\sqrt i-1$ is a root of $X^2+2X+1-i$, so $f(2)=4$.
Finally $f(5b)=f(b)$ when $b\notin k^{\times2}\cup5k^{\times2}$: the field
$k(\sqrt5,\sqrt b)$ is unramified over both $k(\sqrt b)$ and $k(\sqrt{5b})$,
so the two ways of computing its discriminant over $k$ give $2f(b)=2f(5b)$.

(a) Since $k(\sqrt5)/k$ is unramified and $k(\sqrt\alpha)/k$ is ramified,
$\alpha$ is not a square in $k(\sqrt5)$, so $[M_8:\Q_2]=8$. The field $M_8$
contains the conjugate $\sqrt{\bar\alpha}$ of $\sqrt\alpha$, so it is Galois
over $\Q_2$, and an automorphism is determined by the images $\pm i$,
$\pm\sqrt5$ and a square root of the image of $\alpha$; in particular
$\sigma_y$ and $\sigma_z$ exist. Now $\sigma_y^2=1$, $\sigma_z^2$ fixes
$i,\sqrt5$ and negates $\sqrt\alpha$, $\sigma_z$ has order four, and
$\sigma_y\sigma_z\sigma_y=\sigma_z^{-1}$. So
$\Gal(M_8/\Q_2)=\langle\sigma_y,\sigma_z\rangle$ is dihedral of order eight,
with center $\langle\sigma_z^2\rangle$ and maximal abelian subextension
$\Q_2(i,\sqrt5)$.

(b) The quadratic subfields of $M_8$ are $\Q_2(\sqrt b)$ for $b=-1,5,-5$, so
$\sqrt2\notin M_8$, and $M_{16}$ has degree $16$ and group
$\Gal(M_8/\Q_2)\times C_2$. This group has elementary abelian
abelianization, so the maximal unramified subextension of $M_{16}$, which is
cyclic over $\Q_2$ and contains $\Q_2(\sqrt5)$, is $\Q_2(\sqrt5)$. Hence
$M_{16}$ has ramification index $8$, $M_{16}\cap U_4=\Q_2(\sqrt5)$, and
$[L_2:\Q_2]=16\cdot4/2=32$. Since $L_2$ contains $U_4$ and has ramification
index at least $8$, its residue degree is $4$ and its ramification index is
$8$. Restriction identifies $\Gal(L_2/\Q_2)$ with the set of triples
$(\sigma,\epsilon,\mu)$, where $\sigma\in\Gal(M_8/\Q_2)$, $\epsilon=\pm1$ is
the sign by which the element multiplies $\sqrt2$, and $\mu\in\Gal(U_4/\Q_2)$
agrees with $\sigma$ on $\sqrt5$; both sets have $32$ elements. Since $\mathrm{Fr}$
negates $\sqrt5$, the triples
\[
 x=(1,-1,1),\qquad y=(\sigma_y,1,1),\qquad z=(\sigma_z,1,\mathrm{Fr}),\qquad
 w=[y,z]=(\sigma_z^2,1,1)
\]
exist and are the unique elements described in the statement. They act on
$E_2$ as required. The elements $x$ and $w$ are central involutions, $y$ is
an involution, and $z$ has order four. Hence the relations of $D$ hold. Conversely, in
$D$ these relations make $w=[y,z]$ central: $w^y=w^{-1}=w$, and $[w,z]=1$ is
imposed. Since $zy=yzw$, every element of $D$ is
$x^{a_1}y^{a_2}z^{a_3}w^{a_4}$ with $a_1,a_2,a_4\in\Ftwo$ and
$a_3\in\Z/4$, and
\[
 (a_1,a_2,a_3,a_4)(a'_1,a'_2,a'_3,a'_4)
 =(a_1+a'_1,\,a_2+a'_2,\,a_3+a'_3,\,a_4+a'_4+a_3a'_2).
\]
These $32$ normal forms have distinct images in $\Gal(L_2/\Q_2)$: the third
component determines $a_3$, the second $a_1$, and
$\sigma_y^{a_2}\sigma_z^{a_3+2a_4}$ determines $a_2$ and $a_4$. So
$D\cong\Gal(L_2/\Q_2)$. The inertia group consists of the elements trivial
on $U_4$, namely $\langle x,y,w\rangle$.

(c) The Frattini subgroup of $D$ is $\langle z^2,w\rangle$, which is central
of exponent two, and $D$ has exponent four. So Lazard's formula gives
$D_3(D)=[[D,D],D]\,[D,D]^2D^4=1$, and the Hilbert polynomial follows from
Jennings' theorem \cite{Jennings}.

(d) The extension $M_{16}/k$ is abelian of type $(2,2,2)$, generated by
$\sqrt5,\sqrt\alpha,\sqrt2$. We apply the conductor--discriminant formula for
abelian extensions \cite[Chapter~V, Theorem~3.27]{MilneCFT} to the extension
$\Q(i)(\sqrt5,\sqrt{1+2i},\sqrt2)$ of $\Q(i)$, whose completion at the unique
prime above $1+i$ is $M_{16}/k$. Taking $(1+i)$-adic valuations gives
\[
 v_k(\mathfrak d_{M_{16}/k})=\sum_bf(b)=0+0+2+2+4+4+4+4=20,
\]
the sum over the classes $b=1,5,\alpha,5\alpha,2,10,2\alpha,10\alpha$. Hence
\[
 v_2(\mathfrak d_{M_{16}/\Q_2})=[M_{16}:k]\,v_2(\mathfrak d_{k/\Q_2})+20=8\cdot2+20=36.
\]
Since $M_{16}$ has residue degree two over $\Q_2$, its different has
valuation $36/2=18$. The extension $L_2/M_{16}$ is unramified, so the
different of $L_2/\Q_2$ also has valuation $18$, and $18/8=9/4$.
\end{proof}

The field $E_2$ is the maximal elementary abelian subextension of $L_2$, so
$\Gal(L_2/E_2)$ is the Frattini subgroup $\langle z^2,[y,z]\rangle$, which is
central of exponent two. Hence any triple in $\Gal(L_2/\Q_2)$ that acts on
$E_2$ as $x,y,z$ do differs from $x,y,z$ by central elements of order at
most two; it satisfies the relations of $D$ and generates, so it also
defines an isomorphism $D\cong\Gal(L_2/\Q_2)$. The description of the
inertia group, however, refers to the triple of
Lemma~\ref{tw:local-two}(b).

\subsection{The Global Group and Its Complete Presentation}\label{tw:global-section}
Let $B_S$ be the maximal pro-$2$ extension of $B$, in a fixed algebraic
closure, that is unramified at every finite prime outside $S$; no condition
is imposed at $v_1,v_2$. Let $\mathcal G_B=\Gal(B_S/B)$, and let
$E=B(\sqrt{\alpha_0},\dots,\sqrt{\alpha_7})=B(\sqrt V)$ be the Kummer field.
Every prime above $2$ lies in $S$, so a quadratic extension $B(\sqrt b)$ lies
in $B_S$ exactly when $b$ has even valuation outside $S$. By
Lemma~\ref{tw:base}(d), $E$ is the maximal elementary abelian extension of
$B$ in $B_S$, and $[E:B]=256$. For $g\in\mathcal G_B$ define its elementary
image $\bar g\in\Ftwo^8$ by $g(\sqrt{\alpha_i})=(-1)^{\bar g_i}\sqrt{\alpha_i}$.
This identifies $\mathcal G_B/D_2\mathcal G_B=\Gal(E/B)$ with $\Ftwo^8$. We
write vectors as strings $v_0v_1\cdots v_7$. For $a\in V$ let $\chi_a$ be the
Kummer character, $g(\sqrt a)=(-1)^{\chi_a(g)}\sqrt a$; then $a\mapsto\chi_a$
identifies $V$ with the dual of $\Ftwo^8$, and $\chi_{\alpha_i}(g)=\bar g_i$.

For a place $\nu$ of $B$, a choice of a place of $B_S$ above $\nu$ gives a
homomorphism $\mathcal G_\nu\to\mathcal G_B$, the local map at $\nu$; another
choice changes it by an inner automorphism of $\mathcal G_B$. For a finite
place $\nu\notin S$, the local map kills inertia, and its image is generated
by a Frobenius element $\Frob_\nu$. Let $\Sigma=S\cup\{v_1,v_2\}$. For
$\nu\in\Sigma$ we fix generators of $\mathcal G_\nu$ as follows, and use the
same names for their images in $\mathcal G_B$: $\iota_j$ at $v_j$
(Lemma~\ref{tw:local-groups}(a)), so that $\iota_1,\iota_2$ are complex
conjugations; $\tau_\nu,\varphi_\nu$ at
$\nu\in\{\mathfrak q_1,\mathfrak q_2,\mathfrak r_1,\mathfrak r_2\}$, with the
generator $\alpha_i$ of $\nu$ as uniformizer $\varpi$ in
Lemma~\ref{tw:local-groups}(b); and $x_j,y_j,z_j$ at $\mathfrak p_j$, the
images of one fixed triple $x,y,z\in\mathcal G_{\Q_2}=\mathcal G_{\mathfrak p_j}$
lifting the elements $x,y,z$ of Lemma~\ref{tw:local-two}(b). This triple
satisfies the hypothesis of Lemma~\ref{tw:local-groups}(c). We also write
$\Frob_{\mathfrak t_k}$ for $k=0,1,2$.

\begin{table}[htbp]
\centering
\caption{Elementary images in $\Ftwo^8$ of the local generators. Entry $i$
is $1$ exactly when the element negates $\sqrt{\alpha_i}$. The third and
sixth columns give the action on $\sqrt b$ for $b$ in the completion.}
\label{tw:vector-table}
\begin{tabular}{lll@{\qquad}lll}
\toprule
element & image & action & element & image & action\\
\midrule
$\iota_1$ & $10111010$ & sign at $v_1$ & $\iota_2$ & $11000101$ & sign at $v_2$\\
$x_1$ & $00100000$ & $(5,b)_{\mathfrak p_1}$ & $x_2$ & $00010000$ & $(5,b)_{\mathfrak p_2}$\\
$y_1$ & $11110010$ & $(-1,b)_{\mathfrak p_1}$ & $y_2$ & $10000001$ & $(-1,b)_{\mathfrak p_2}$\\
$z_1$ & $10000100$ & $(-2,b)_{\mathfrak p_1}$ & $z_2$ & $11111000$ & $(-2,b)_{\mathfrak p_2}$\\
$\tau_{\mathfrak q_1}$ & $00001000$ & $(-1,b)_{\mathfrak q_1}$ & $\varphi_{\mathfrak q_1}$ & $10100111$ & $(-\alpha_4,b)_{\mathfrak q_1}$\\
$\tau_{\mathfrak q_2}$ & $00000100$ & $(-1,b)_{\mathfrak q_2}$ & $\varphi_{\mathfrak q_2}$ & $11101011$ & $(-\alpha_5,b)_{\mathfrak q_2}$\\
$\tau_{\mathfrak r_1}$ & $00000010$ & $(2,b)_{\mathfrak r_1}$ & $\varphi_{\mathfrak r_1}$ & $01100101$ & $(-\alpha_6,b)_{\mathfrak r_1}$\\
$\tau_{\mathfrak r_2}$ & $00000001$ & $(2,b)_{\mathfrak r_2}$ & $\varphi_{\mathfrak r_2}$ & $01011010$ & $(-\alpha_7,b)_{\mathfrak r_2}$\\
$\Frob_{\mathfrak t_1}$ & $00100111$ & $(29,b)_{\mathfrak t_1}$ & $\Frob_{\mathfrak t_2}$ & $00011011$ & $(29,b)_{\mathfrak t_2}$\\
$\Frob_{\mathfrak t_0}$ & $01110000$ & $(7,b)_{\mathfrak t_0}$ & & &\\
\bottomrule
\end{tabular}
\end{table}

Table~\ref{tw:vector-table} lists the elementary images. They follow from
Lemma~\ref{tw:local-groups} and Table~\ref{tw:squareclass-table}. For
example, at the places above $3$ and $5$ the vector of $\tau_\nu$ marks the
odd valuations, and that of $\varphi_\nu$ marks the units that are not
squares modulo $\nu$, with a zero at the generator of $\nu$. The columns
headed ``action'' express the same actions as Hilbert symbols, that is, as
local Artin symbols \cite[Chapter~III, Remark~4.5]{MilneCFT}. The
supplementary program \texttt{kummer241.gp} checks the norms of the
$\alpha_i$, the generators of the primes in Lemma~\ref{tw:base}(c) and the
residues of $\sqrt{241}$ in Table~\ref{tw:squareclass-table}, and recomputes
Table~\ref{tw:vector-table} from these Hilbert symbols.

Fix a minimal presentation $1\to R\to\mathcal F\to\mathcal G_B\to1$, where
$\mathcal F$ is a free pro-$2$ group of rank $8$. Minimality means
$R\subset D_2\mathcal F$, so $\mathcal F/D_2\mathcal F=\Ftwo^8$ with the
coordinates above. Put $\mathfrak L=\operatorname{gr}\mathcal F$, so that
$\dim\mathfrak L_1=8$, $\dim\mathfrak L_2=36$ and $\dim\mathfrak L_3=168$.
For $v,w\in\mathfrak L_1$ we have $(v+w)^{[2]}=v^{[2]}+w^{[2]}+[v,w]$, and the
square of an element of $\mathcal F$ with image $v$ has class $v^{[2]}$. For
each local generator $g$ above choose a lift $\hat g\in\mathcal F$, and define
the local relators
\begin{equation}\label{tw:relators}
 \rho_{v_j}=\hat\iota_j^{\,2},\qquad
 \rho_\nu=\hat\varphi_\nu\hat\tau_\nu\hat\varphi_\nu^{-1}\hat\tau_\nu^{-p}
 \ (\nu\in\{\mathfrak q_1,\mathfrak q_2,\mathfrak r_1,\mathfrak r_2\}),\qquad
 \rho_{\mathfrak p_j}=r(\hat x_j,\hat y_j,\hat z_j),
\end{equation}
where $p\in\{3,5\}$ is the rational prime below $\nu$, and $r$ is a
defining relator of $\mathcal G_{\Q_2}$ with respect to the fixed triple
$x,y,z$, the same for $j=1,2$. They lie in $R$, because each local relation
holds in $\mathcal G_\nu$. Their classes in $\mathfrak L_2$ are obtained by
substituting elementary images into the quadratic initials of
Lemma~\ref{tw:local-groups}:
\begin{equation}\label{tw:local-initials}
 \bar\iota_j^{\,[2]},\qquad
 [\bar\tau_\nu,\bar\varphi_\nu]+\tfrac{p-1}2\bar\tau_\nu^{\,[2]},\qquad
 \bar y_j^{\,[2]}+[\bar x_j,\bar y_j]+[\bar x_j,\bar z_j].
\end{equation}
The results below hold for every minimal presentation and every choice of
lifts.

The classes \eqref{tw:local-initials} have an arithmetic meaning. The
elements $e_i^{[2]}$ and $[e_i,e_j]$ ($i<j$) of the standard basis form a
basis of $\mathfrak L_2$, so there is a linear isomorphism $\beta$ from
$\mathfrak L_2$ onto the space of symmetric bilinear forms on $V$ with
\[
 \beta(v^{[2]})(a,b)=\chi_a(v)\chi_b(v),\qquad
 \beta([v,w])(a,b)=\chi_a(v)\chi_b(w)+\chi_a(w)\chi_b(v)
 \qquad(v,w\in\mathfrak L_1).
\]
It is defined by these formulas for basis vectors $v,w$, which send the basis
of $\mathfrak L_2$ to a basis of the symmetric bilinear forms. The formulas
then hold for all $v,w$, because the bracket is bilinear,
$(v+w)^{[2]}=v^{[2]}+w^{[2]}+[v,w]$, and squaring is the identity on $\Ftwo$.
Equivalently, if $\xi\in\mathfrak L_2$ is written in the free algebra
$\Ftwo\langle X_0,\dots,X_7\rangle$ as $\sum_{i,j}n_{ij}X_iX_j$, then
$\beta(\xi)(\alpha_i,\alpha_j)=n_{ij}$.
For $\nu\in\Sigma$, the local Hilbert form $H_\nu$ on $V$ is the symmetric
bilinear form with $H_\nu(a,b)=1$ if $(a,b)_\nu=-1$ and $H_\nu(a,b)=0$
otherwise.

\begin{lemma}\label{tw:local-forms}
For every $\nu\in\Sigma$, $\beta$ maps the class of $\rho_\nu$ to $H_\nu$. The
eight classes \eqref{tw:local-initials} sum to zero, and the seven classes
with $\nu\ne\mathfrak p_1$ are linearly independent.
\end{lemma}

\begin{proof}
For a local generator $g$ at $\nu$, $\chi_a(\bar g)$ is the value at $g$ of
the local Kummer character of the image of $a$ in $B_\nu$. Hence
$\beta$ of the class of $\rho_\nu$ is the pullback to $V$ of a symmetric
bilinear form on $B_\nu^\times/B_\nu^{\times2}$, as is $H_\nu$, and it
suffices to compare the two on a basis of $B_\nu^\times/B_\nu^{\times2}$. At a
real place the only class is $-1$, and both forms take the value $1$ on
$(-1,-1)$. At $\nu$ above $p\in\{3,5\}$, take the basis $\varpi,u$ with
$\varpi$ the generator of $\nu$ and $u$ a nonsquare unit. By
Lemma~\ref{tw:local-groups}(b), $\chi_\varpi(\tau)=\chi_u(\varphi)=1$ and
$\chi_\varpi(\varphi)=\chi_u(\tau)=0$, so the form of
$[\tau,\varphi]+\frac{p-1}2\tau^{[2]}$ takes the values $\frac{p-1}2$, $1$,
$0$ on $(\varpi,\varpi)$, $(\varpi,u)$, $(u,u)$; by \eqref{tw:hilbert} these
are the indicators of $(\varpi,\varpi)_p=(-1)^{(p-1)/2}$, $(\varpi,u)_p=-1$
and $(u,u)_p=1$. At $\mathfrak p_j$ the characters $\chi_2,\chi_{-5},\chi_5$
are dual to $x,y,z$, as in the proof of Lemma~\ref{tw:local-groups}(c). So
$y^{[2]}$ contributes the entry at $(-5,-5)$, $[x,y]$ the entries at $(2,-5)$
and $(-5,2)$, and $[x,z]$ those at $(2,5)$ and $(5,2)$: the form of
$y^{[2]}+[x,y]+[x,z]$ on the basis $2,-5,5$ is the matrix of the Hilbert
symbol displayed in that proof, from which this initial was obtained.

For $a,b\in V$ and a place $\nu'\notin\Sigma$, the place $\nu'$ is finite
and not above $2$, and $a,b$ are units at $\nu'$, so $(a,b)_{\nu'}=1$
\cite[Chapter~VIII, 5.8]{MilneCFT}. By the
product formula for the Hilbert symbol \cite[Chapter~V, Example~5.4]{MilneCFT},
$\sum_{\nu\in\Sigma}H_\nu=0$, so the eight classes sum to zero.

Finally, by Table~\ref{tw:squareclass-table} and \eqref{tw:hilbert}, the
seven pairs
\[
 (-1,\alpha_2),\ (-1,\varepsilon),\ (-1,\alpha_7),\ (-1,\alpha_4),\
 (-1,\alpha_5),\ (\alpha_2,\alpha_6),\ (\varepsilon,\alpha_7)
\]
have Hilbert symbol $-1$ exactly at $\{\mathfrak p_1,v_1\}$,
$\{\mathfrak p_1,v_2\}$, $\{\mathfrak p_2,v_2\}$, $\{\mathfrak q_1,v_1\}$,
$\{\mathfrak q_2,v_2\}$, $\{\mathfrak r_1,v_1\}$ and $\{\mathfrak r_2,v_2\}$;
the supplementary program \texttt{cup241.gp} recomputes these symbols.
Evaluating a relation $\sum_{\nu\ne\mathfrak p_1}n_\nu H_\nu=0$, with
$n_\nu\in\Ftwo$, at the first two pairs gives $n_{v_1}=n_{v_2}=0$, and then
at the other five pairs gives
$n_{\mathfrak p_2}=n_{\mathfrak q_1}=n_{\mathfrak q_2}=n_{\mathfrak r_1}=n_{\mathfrak r_2}=0$.
\end{proof}

\begin{lemma}\label{tw:complete-global-presentation}
The group $\mathcal G_B$ has generator rank $8$ and relation rank $7$. The
seven relators $\rho_\nu$ with $\nu\in\Sigma\setminus\{\mathfrak p_1\}$ form
a minimal system of defining relations: their normal closure in
$\mathcal F$ is $R$. In particular the dyadic local relator
$\rho_{\mathfrak p_1}$ lies in the normal closure of the other seven.
\end{lemma}

\begin{proof}
The generator rank is $\dim H^1(\mathcal G_B,\Ftwo)=8$
\cite[Proposition~3.9.1]{NeukirchSchmidtWingberg2008}.

Next we bound $H^2$. Let $X=\operatorname{Spec}\mathcal O_B[1/30]$. Since
$2$ is invertible on $X$, the sheaf $\mu_2$ is the constant sheaf $\Ftwo$.
The finite subextensions of $B_S/B$ correspond to the connected finite étale
Galois $2$-covers $Y\to X$, and the Hochschild--Serre spectral sequence of
this tower \cite[Theorem~14.9 and Remark~14.10]{MilneEtale} gives an exact
sequence
\[
 0\to H^1(\mathcal G_B,\Ftwo)\to H^1(X,\Ftwo)\to
 \bigl(\varinjlim_Y H^1(Y,\Ftwo)\bigr)^{\mathcal G_B}\to
 H^2(\mathcal G_B,\Ftwo)\to H^2(X,\Ftwo).
\]
A class in $H^1(Y,\Ftwo)$ is an étale double cover of $Y$. Its Galois closure
over $X$ is a finite étale Galois cover whose group embeds in the wreath
product of $\Gal(Y/X)$ with $C_2$, hence a $2$-cover in the tower, over which
the class vanishes. So the direct limit is zero, and
$H^2(\mathcal G_B,\Ftwo)$ injects into $H^2(X,\mu_2)$. The Kummer sequence and
$\Pic(X)=0$ identify $H^2(X,\mu_2)$ with $\Br(X)[2]$, where
$\Br(X)=H^2(X,\mathbb G_m)$. By \cite[Theorem~6.8.3 and
Example~6.9.4]{Poonen}, $\Br(X)$ is the kernel of the sum of the local
invariants $\bigoplus_{\nu\in\Sigma}\Br(B_\nu)\to\Q/\Z$. Each of the eight
groups $\Br(B_\nu)[2]$ has order two, and the sum of the invariants maps their
direct sum onto $\frac12\Z/\Z$, so $\dim\Br(X)[2]=7$. We use only the
resulting bound $\dim H^2(\mathcal G_B,\Ftwo)\le\dim\Br(X)[2]\le7$.

By Lemma~\ref{tw:local-forms}, the classes in $\mathfrak L_2$ of the seven
relators $\rho_\nu$, $\nu\ne\mathfrak p_1$, are linearly independent. Since
$R\subset D_2\mathcal F$, we have $R^2[R,\mathcal F]\subset D_3\mathcal F$, so
the classes of these seven relators in $R/R^2[R,\mathcal F]$ are independent.
On the other hand
$\dim R/R^2[R,\mathcal F]=\dim H^2(\mathcal G_B,\Ftwo)\le7$ by the five-term
exact sequence of the presentation
\cite[Proposition~3.9.5]{NeukirchSchmidtWingberg2008},
\cite[(3.3)--(3.4)]{Quadrelli2021}. Hence the seven classes form a basis, the
relation rank is $7$, and by the pro-$2$ Nakayama lemma
\cite[Corollary~3.9.3]{NeukirchSchmidtWingberg2008} the seven relators
generate $R$ as a normal subgroup. Finally $\rho_{\mathfrak p_1}\in R$.
\end{proof}

By Lemma~\ref{tw:local-forms}, the only linear relation among the eight
classes \eqref{tw:local-initials} is that their sum vanishes, and this
relation is Hilbert reciprocity. Hence any seven of the eight classes are
linearly independent, and the proof of
Lemma~\ref{tw:complete-global-presentation} shows that any seven of the eight
local relators generate $R$ as a normal subgroup. We omit the dyadic relator
at $\mathfrak p_1$.

\begin{remark}\label{tw:cup-remark}
By Lemma~\ref{tw:complete-global-presentation}, the bound
$\dim H^2(\mathcal G_B,\Ftwo)\le\dim\Br(X)[2]=7$ is attained, so
$H^2(\mathcal G_B,\Ftwo)\to\Br(X)[2]$ is an isomorphism. This can also be seen
as in the proof of \cite[Lemma~2.1]{Naslund0418235}: the cup product
$\chi_a\cup\chi_b$ maps to the class of the quaternion algebra $(a,b)$, whose
invariant at $\nu$ is given by $(a,b)_\nu$ \cite[Chapter~III,
Remark~4.7]{MilneCFT}, and the seven pairs in the proof of
Lemma~\ref{tw:local-forms} give seven classes with independent invariant
vectors. The supplementary program \texttt{cup241.gp} computes the invariant
vectors at the eight places of $\Sigma$ of all $36$ algebras
$(\alpha_i,\alpha_j)$.
\end{remark}

\subsection{The Prescribed Local Groups}\label{tw:prescribed-section}
We now cut $\mathcal G_B$ down to prescribed local groups. At the dyadic
places we use the restriction
$\mathcal G_{\mathfrak p_j}=\mathcal G_{\Q_2}\to\Gal(L_2/\Q_2)=D$ of
Lemma~\ref{tw:local-two}. At a place $\nu$ above $p\in\{3,5\}$, the maximal
elementary abelian quotient of $\mathcal G_\nu$ is
$\Gal(\Q_p(\sqrt u,\sqrt p)/\Q_p)\cong C_2\times C_2$ for a nonsquare unit $u$;
this quotient has ramification index and residue degree two. At
$\mathfrak t_1,\mathfrak t_2$, whose completions are $\Q_{29}$, and at
$\mathfrak t_0$, whose residue field is $\mathbb F_{49}$, we bound the order
of the Frobenius by four.

\begin{definition}\label{tw:GB}
Let $N$ be the normal closure in $\mathcal G_B$ of the images of
$\ker(\mathcal G_{\mathfrak p_j}\to D)$ for $j=1,2$, of the Frattini subgroups
$D_2\mathcal G_\nu$ for $\nu\in\{\mathfrak q_1,\mathfrak q_2,\mathfrak r_1,\mathfrak r_2\}$,
and of $\Frob_{\mathfrak t_k}^4$ for $k=0,1,2$. Put $G_B=\mathcal G_B/N$ and
$\overline G_B=G_B/D_4G_B$, and let $B_G\subset B_S$ be the fixed field of
$N$, so that $\Gal(B_G/B)=G_B$.
\end{definition}

This does not depend on the choice of places of $B_S$ above each $\nu$. The
image of $\mathcal G_{\mathfrak p_j}$ in $G_B$ is a quotient of $D$, the image
of $\mathcal G_\nu$ for $\nu$ above $3$ or $5$ is a quotient of
$C_2\times C_2$ in which inertia maps to the image of $\tau_\nu$, and
$\Frob_{\mathfrak t_k}$ has order dividing four in $G_B$.
Subsection~\ref{tw:retained-section} shows that these are equalities. Since
$D$ and $\mathcal G_{\Q_2}$ both have generator rank three,
$\ker(\mathcal G_{\mathfrak p_j}\to D)$ lies in the Frattini subgroup of
$\mathcal G_{\mathfrak p_j}$. Hence $N\subset D_2\mathcal G_B$,
$G_B/D_2G_B=\Gal(E/B)$, and $E\subset B_G$ is the fixed field of $D_2G_B$.

\begin{lemma}\label{tw:GB-presentation}
The kernel $N_B$ of $\mathcal F\to\mathcal G_B\to G_B$ is the normal closure
of the following thirty elements:
\begin{enumerate}
\item[(i)] the seven relators $\rho_\nu$, $\nu\in\Sigma\setminus\{\mathfrak p_1\}$;
\item[(ii)] for $j=1,2$: $\hat x_j^{\,2}$, $[\hat x_j,\hat y_j]$,
 $[\hat x_j,\hat z_j]$, $[[\hat y_j,\hat z_j],\hat z_j]$, $\hat z_j^{\,4}$,
 $[\hat y_j,\hat z_j]^2$;
\item[(iii)] for $\nu\in\{\mathfrak q_1,\mathfrak q_2,\mathfrak r_1,\mathfrak r_2\}$:
 $\hat\tau_\nu^{\,2}$ and $\hat\varphi_\nu^{\,2}$;
\item[(iv)] for $k=0,1,2$: $\hat F_k^{\,4}$, where $\hat F_k$ is the chosen
 lift of $\Frob_{\mathfrak t_k}$.
\end{enumerate}
Moreover $N_B\subset D_2\mathcal F$, and $\rho_{\mathfrak p_1}\in N_B$.
\end{lemma}

\begin{proof}
By Lemma~\ref{tw:complete-global-presentation}, it suffices to show that the
images in $\mathcal G_B$ of the elements (ii)--(iv) generate $N$ as a normal
subgroup. Let $\mathcal F_3$ be free on $x,y,z$, and let $N_D$ be the normal
closure of the seven relators of $D$ in Lemma~\ref{tw:local-two}, so that
$\mathcal F_3/N_D=D$. The relator $r$ maps to $1$ in $D$, so $r\in N_D$, and
$\ker(\mathcal G_{\Q_2}\to D)$ is the image of $N_D$. Since
$N_D\subset D_2\mathcal F_3$, we have
$N_D^2[N_D,\mathcal F_3]\subset D_3\mathcal F_3$. The seven relators of $D$
span $N_D/N_D^2[N_D,\mathcal F_3]$
\cite[Corollary~3.9.3]{NeukirchSchmidtWingberg2008}; write the class of $r$ as
a combination of them. The classes in $\operatorname{gr}_2\mathcal F_3$ of
$x^2,y^2,[x,y],[x,z]$ are independent, the relators $z^4$, $[y,z]^2$ and
$[[y,z],z]$ lie in $D_3\mathcal F_3$, and $r$ has class
$y^{[2]}+[x,y]+[x,z]$; so $y^2$ has coefficient one. Replacing $y^2$ by $r$
therefore still gives a spanning set, and by
\cite[Corollary~3.9.3]{NeukirchSchmidtWingberg2008} the elements $r$, $x^2$,
$[x,y]$, $[x,z]$, $[[y,z],z]$, $z^4$, $[y,z]^2$ generate $N_D$ as a normal
subgroup. Since $r=1$ in $\mathcal G_{\Q_2}$, the other six generate
$\ker(\mathcal G_{\Q_2}\to D)$. For $\nu$ above $p\in\{3,5\}$, the quotient
of $\mathcal G_\nu$ by the normal closure of $\tau^2,\varphi^2$ is generated
by two involutions which commute, because
$\varphi\tau\varphi^{-1}=\tau^p=\tau$ there; so that normal closure is
$D_2\mathcal G_\nu$. Finally $N\subset D_2\mathcal G_B$ and
$R\subset D_2\mathcal F$ give $N_B\subset D_2\mathcal F$, and
$\rho_{\mathfrak p_1}\in R\subset N_B$.
\end{proof}

\subsection{The Retained Quotient}\label{tw:retained-section}
For $n\ge1$ let $R_n\subset\mathfrak L_n$ be the image of
$N_B\cap D_n\mathcal F$. Since $N_B\subset D_2\mathcal F$ and $D_nG_B$ is
the image of $D_n\mathcal F$,
\[
 \operatorname{gr}_1G_B=\Ftwo^8,\qquad
 \operatorname{gr}_nG_B\cong\mathfrak L_n/R_n .
\]
Consider the following $22$ elements of $\mathfrak L_2$, written with the
vectors of Table~\ref{tw:vector-table}:
\begin{equation}\label{tw:initials-list}
\begin{aligned}
 &\bar\iota_1^{\,[2]},\ \bar\iota_2^{\,[2]};\qquad
 [\bar\tau_\nu,\bar\varphi_\nu]+\bar\tau_\nu^{\,[2]}\ (\nu=\mathfrak q_1,\mathfrak q_2),\quad
 [\bar\tau_\nu,\bar\varphi_\nu]\ (\nu=\mathfrak r_1,\mathfrak r_2);\\
 &\bar y_j^{\,[2]}+[\bar x_j,\bar y_j]+[\bar x_j,\bar z_j],\quad
 \bar x_j^{\,[2]},\quad [\bar x_j,\bar y_j],\quad [\bar x_j,\bar z_j]\quad(j=1,2);\\
 &\bar\tau_\nu^{\,[2]},\ \bar\varphi_\nu^{\,[2]}\quad
 (\nu=\mathfrak q_1,\mathfrak q_2,\mathfrak r_1,\mathfrak r_2).
\end{aligned}
\end{equation}
They are the classes of the eight local relators \eqref{tw:local-initials}
and of the fourteen quadratic elements $\hat x_j^{\,2}$,
$[\hat x_j,\hat y_j]$, $[\hat x_j,\hat z_j]$, $\hat\tau_\nu^{\,2}$,
$\hat\varphi_\nu^{\,2}$ of Lemma~\ref{tw:GB-presentation}.

\begin{lemma}\label{tw:quadratic-layer}
The space $R_2$ is spanned by the $22$ elements \eqref{tw:initials-list}, and
$\dim R_2=21$; their only linear relation is that the eight classes of local
relators sum to zero. Hence $\dim\operatorname{gr}_2G_B=15$. Consequently, if
$g\in G_B$ has nonzero elementary image $v$, then $g$ has order at least four
in $G_B/D_3G_B$ exactly when $v^{[2]}\notin R_2$; otherwise $g$ has order two
there.
\end{lemma}

\begin{proof}
Conjugates of an element of $D_2\mathcal F$ have the same class modulo
$D_3\mathcal F$, classes of products add, and $\mathfrak L_2$ is finite.
Hence, by Lemma~\ref{tw:GB-presentation}, $R_2$ is spanned by the classes of
the thirty generators of $N_B$. Those in (i) have the classes
\eqref{tw:local-initials}, and the quadratic elements of (ii) and (iii) have
the classes $\bar x_j^{\,[2]}$, $[\bar x_j,\bar y_j]$, $[\bar x_j,\bar z_j]$,
$\bar\tau_\nu^{\,[2]}$, $\bar\varphi_\nu^{\,[2]}$. The remaining generators,
$[[\hat y_j,\hat z_j],\hat z_j]$, $\hat z_j^{\,4}$, $[\hat y_j,\hat z_j]^2$ and
the fourth powers in (iv), lie in $D_3\mathcal F$. The class of
$\rho_{\mathfrak p_1}$ also lies in $R_2$, since $\rho_{\mathfrak p_1}\in N_B$.
The supplementary program \texttt{lie241.py} computes the rank of the $22$
elements in the $36$-dimensional space $\mathfrak L_2$: it is $21$. The sum of
the eight local classes vanishes by Lemma~\ref{tw:local-forms}, that is, by
Hilbert reciprocity, so this is the only relation; none of the fourteen
quadratic elements of (ii) and (iii) is involved in it. For the last statement, $g^2\in D_2G_B$ has class $v^{[2]}$ in
$\operatorname{gr}_2G_B=\mathfrak L_2/R_2$, so $g^2\in D_3G_B$ exactly when
$v^{[2]}\in R_2$.
\end{proof}

\begin{lemma}\label{tw:retained-quotient}
\begin{enumerate}
\item[(a)] $\dim\operatorname{gr}_3G_B=26$, so $|\overline G_B|=2^{8+15+26}=2^{49}$.
\item[(b)] For $j=1,2$, the local map induces an injection
 $D\to G_B/D_3G_B$, and $\operatorname{gr}_nD\to\operatorname{gr}_nG_B$ is
 injective for $n=1,2$. For $\nu$ above $3$ or $5$, the image of
 $\mathcal G_\nu$ in $G_B$ is
 $\langle\tau_\nu\rangle\times\langle\varphi_\nu\rangle\cong C_2\times C_2$,
 and it injects into $G_B/D_2G_B$. For $k=0,1,2$, $\Frob_{\mathfrak t_k}$ has
 order exactly four in $G_B$ and in $G_B/D_3G_B$. For $j=1,2$, $\iota_j$ has
 order two in $G_B/D_2G_B$.
\item[(c)] The conjugacy class of $\iota_1$ in $\overline G_B$ has $2^{15}$
 elements.
\item[(d)] The elementary images $\bar\iota_1\ne\bar\iota_2$ lie outside the
 span of the elementary images of the local generators at each of
 $\mathfrak p_1,\mathfrak p_2,\mathfrak q_1,\mathfrak q_2,\mathfrak r_1,\mathfrak r_2,\mathfrak t_0,\mathfrak t_1,\mathfrak t_2$.
\end{enumerate}
\end{lemma}

\begin{proof}
(a) Work in $\mathcal F/D_4\mathcal F$, in which $D_2\mathcal F/D_4\mathcal F$
is elementary abelian and $D_3\mathcal F/D_4\mathcal F$ is central. Let
$\rho_1,\dots,\rho_{21}$ be the generators of $N_B$ of degree two, namely
(i), the first three elements of (ii) for each $j$, and (iii). Their classes
are the elements of \eqref{tw:initials-list} other than the class of
$\rho_{\mathfrak p_1}$, and these are independent by
Lemma~\ref{tw:quadratic-layer}. The cubic generators are
$\omega_j=[[\hat y_j,\hat z_j],\hat z_j]$, and the remaining generators lie
in $D_4\mathcal F$. For $\rho\in D_2\mathcal F$ and $g\in\mathcal F$ we have
$g^{-1}\rho g=\rho[\rho,g]$, and modulo $D_4\mathcal F$ the commutator
$[\rho,g]$ is central and depends only, and linearly, on the classes of
$\rho$ in $\mathfrak L_2$ and of $g$ in $\mathfrak L_1$. Hence the image of
$N_B$ in $\mathcal F/D_4\mathcal F$ is the subgroup generated by the $\rho_i$,
the $[\rho_i,f]$ for $f$ in a basis of $\mathcal F$, and the $\omega_j$. An
element $\prod_i\rho_i^{a_i}$ times an element of $D_3\mathcal F$ lies in
$D_3\mathcal F$ only if all $a_i=0$. Therefore $R_3$ is spanned by the
brackets $[\bar\rho_i,\xi]$, for $\xi$ in a basis of $\mathfrak L_1$, and by
$[[\bar y_j,\bar z_j],\bar z_j]$. The supplementary program
\texttt{lie241c.py} computes these elements in the free associative algebra,
where $[a,\xi]=a\xi+\xi a$: they span a space of dimension $142$ in the
$168$-dimensional space $\mathfrak L_3$. So $\dim\operatorname{gr}_3G_B=26$.

(b) The local map $\mathcal G_{\mathfrak p_j}\to G_B$ factors through $D$ by
Definition~\ref{tw:GB}, and the induced map $D\to G_B$ preserves the
Zassenhaus filtrations. It is injective on $\operatorname{gr}_1$ because
$\bar x_j,\bar y_j,\bar z_j$ are independent, and on $\operatorname{gr}_2$
because $\bar z_j^{\,[2]}$ and $[\bar y_j,\bar z_j]$, the images of
$\overline{z^2}$ and $\overline{[y,z]}$, are independent modulo $R_2$; the
program \texttt{lie241.py} checks both facts. If $1\neq g\in D$, let $n\le2$
be maximal with $g\in D_n(D)$; then the image of $g$ is nonzero in
$\operatorname{gr}_nG_B$, so it is not in $D_3G_B$. For $\nu$ above $3$ or
$5$, the vectors $\bar\tau_\nu$ and $\bar\varphi_\nu$ are independent, and the
image of $\mathcal G_\nu$ is a quotient of $C_2\times C_2$. For the Frobenius
elements, \texttt{lie241.py} checks that
$\overline{\Frob}_{\mathfrak t_k}^{\,[2]}\notin R_2$, so
Lemma~\ref{tw:quadratic-layer} gives order at least four in $G_B/D_3G_B$,
while Definition~\ref{tw:GB} gives order at most four in $G_B$. Finally
$\bar\iota_j\neq0$.

(c) Write $\overline G=\overline G_B$ and $\iota=\iota_1$. The map
$\mathfrak L_1\to\operatorname{gr}_2\overline G$, $v\mapsto[\bar\iota,v]$, has
rank $7$ by \texttt{lie241.py}, so its kernel is the line through
$\bar\iota$. The map
$\operatorname{gr}_2\overline G\to\operatorname{gr}_3\overline G$,
$u\mapsto[u,\bar\iota]$, has rank $8$ by \texttt{lie241c.py}. Since
$D_2\overline G$ is abelian and $D_3\overline G$ is central in $\overline G$,
the map $h\mapsto[h,\iota]$ is a homomorphism $D_2\overline G\to D_3\overline G$;
it factors through $\operatorname{gr}_2\overline G$ and equals the second map.
Its kernel $C_{\overline G}(\iota)\cap D_2\overline G$ therefore has order
$2^{15+26-8}=2^{33}$. If $g\in C_{\overline G}(\iota)$, then
$[\bar\iota,\bar g]=0$ in $\operatorname{gr}_2\overline G$, so
$\bar g\in\{0,\bar\iota\}$; conversely $\iota\in C_{\overline G}(\iota)$.
Hence $|C_{\overline G}(\iota)|=2^{34}$, and the class of $\iota$ has
$2^{49-34}=2^{15}$ elements.

(d) This is a finite check on Table~\ref{tw:vector-table}, carried out by
\texttt{lie241.py}.
\end{proof}

\subsection{Filtered Fox Images}\label{tw:fox-section}
We now recall the filtered Fox calculus of \cite[Section~2]{Naslund0418235}.
Let $\Delta$ be a finite $2$-group, $\Lambda=\Ftwo[\Delta]$, and
$\mathfrak I$ its augmentation ideal. Filtrations are decreasing, with
$F^0M=M$ and $F^nM=0$ for large $n$. For a finite-dimensional filtered space
put $h_M(t)=\sum_n\dim(F^nM/F^{n+1}M)\,t^n$. We filter $\Lambda$ by
$F^n\Lambda=\mathfrak I^n$, and $\Lambda^m(-1)$ denotes $\Lambda^m$ with
$F^n=(\mathfrak I^{n-1})^m$ for $n\ge1$, so that
$h_{\Lambda^m(-1)}=mt\,h_\Lambda$. Subspaces carry the induced filtration
unless stated otherwise. For $0<t<1$,
\begin{equation}\label{tw:cumulative}
 h_M(t)=(1-t)\sum_{n\ge1}\dim(M/F^nM)\,t^{n-1}.
\end{equation}
Hence, for $0<t<1$: (h1) if $f\colon M\to M'$ is surjective with
$f(F^nM)\subset F^nM'$, then $h_{M'}(t)\le h_M(t)$; (h2) if $U_1\subset U_2$
are subspaces of the same filtered space, then $h_{U_1}(t)\le h_{U_2}(t)$;
(h3) if $f$ is surjective and strict, $f(F^nM)=F^nM'$, then
$h_M=h_{\ker f}+h_{M'}$; and (h4) for subspaces $U_1,U_2$ of a filtered
space,
\begin{equation}\label{tw:filtered-sum}
 h_{U_1+U_2}(t)\le h_{U_1}(t)+h_{U_2}(t)-h_{U_1\cap U_2}(t).
\end{equation}
For (h4), apply (h3) to $U_1\oplus U_2\to U_1+U_2$, $(u_1,u_2)\mapsto u_1+u_2$,
whose kernel is $U_1\cap U_2$, with the image filtration $F^nU_1+F^nU_2$ on
$U_1+U_2$; this filtration is contained in the induced one, and (h1) applies
to the identity map.

Let $\mathcal F$ be a free pro-$2$ group with basis $f_1,\dots,f_n$ and
$\pi\colon\mathcal F\to\Delta$ a surjection. Since $\Lambda^n\rtimes\Delta$,
with $\Delta$ acting by left multiplication, is a finite $2$-group, there is
a unique homomorphism $\mathcal F\to\Lambda^n\rtimes\Delta$ sending $f_i$ to
$(e_i,\pi(f_i))$, where $e_1,\dots,e_n$ is the standard basis of $\Lambda^n$.
We write it as $w\mapsto(\partial w,\pi(w))$ and call $\partial w$ the row of
$w$. For $w$ in the discrete free group on the $f_i$, the row
$\partial w\in\Lambda^n$ is the image of the vector of free derivatives of $w$
\cite[Section~2]{Fox1953}. We use the following properties.
\begin{enumerate}
\item[(F1)] $\partial(ww')=\partial w+\pi(w)\partial w'$.
\item[(F2)] $\sum_i\partial_iw\,(\pi(f_i)-1)=\pi(w)-1$.
\item[(F3)] If $\pi(\rho)=1$, then $\partial(u\rho u^{-1})=\pi(u)\partial\rho$;
 and $\partial$ restricts to a continuous homomorphism on $\ker\pi$.
\item[(F4)] If $\mathcal F'$ is free on $f'_1,\dots,f'_m$,
 $u_1,\dots,u_m\in\mathcal F$, and $\partial'$ denotes the rows for the map
 $\mathcal F'\to\Delta$, $f'_l\mapsto\pi(u_l)$, then
 $\partial(W(u_1,\dots,u_m))=\sum_l\partial'_lW\cdot\partial u_l$ for
 $W\in\mathcal F'$.
\end{enumerate}
(F1) is the multiplication of the semidirect product. For (F2), the pairs
$(a,g)$ with $\sum a_i(\pi(f_i)-1)=g-1$ form a subgroup containing the
images of the $f_i$. (F3) follows from (F1). For (F4), both sides are the first
components of homomorphisms $\mathcal F'\to\Lambda^n\rtimes\Delta$ that agree
on the basis.

\begin{lemma}\label{tw:fox-rows}
Let $\mathcal M=\{a\in\Lambda^n:\sum_ia_i(\pi(f_i)-1)=0\}$, with the
filtration induced from $\Lambda^n(-1)$. Then
$h_{\mathcal M}=1-(1-nt)h_\Lambda$ and $\mathcal M=\partial(\ker\pi)$. If
$\rho$ lies in the normal closure of a set $Y\subset\ker\pi$, then
$\partial\rho$ lies in the $\Lambda$-span of $\partial Y$. In particular,
$\mathcal M$ is the $\Lambda$-span of the rows of any set of normal
generators of $\ker\pi$.
\end{lemma}

\begin{proof}
Since the $\pi(f_i)$ generate $\Delta$, we have
$\mathfrak I=\sum_i\Lambda(\pi(f_i)-1)\subset\sum_i\Ftwo(\pi(f_i)-1)+\mathfrak I^2$,
hence $\mathfrak I^m=\sum_i\mathfrak I^{m-1}(\pi(f_i)-1)$ for all $m\ge1$, by
induction and because $\mathfrak I$ is nilpotent. So
$\Lambda^n(-1)\to\mathfrak I$, $a\mapsto\sum a_i(\pi(f_i)-1)$, is strict and
surjective with kernel $\mathcal M$, and (h3) gives
$h_{\mathcal M}=nt\,h_\Lambda-(h_\Lambda-1)$.

By (F2), $\partial(\ker\pi)\subset\mathcal M$. Conversely, let
$\mathcal F^\circ\subset\mathcal F$ be the dense subgroup generated by the
$f_i$, which is a free group on them. Then $\pi(\mathcal F^\circ)=\Delta$, and
the kernel of $\Ftwo[\mathcal F^\circ]\to\Lambda$ is
$\Ftwo[\mathcal F^\circ]\,(\mathcal F^\circ\cap\ker\pi-1)$. Lift
$a\in\mathcal M$ to $\tilde a\in\Ftwo[\mathcal F^\circ]^n$. Then
$\sum_i\tilde a_i(f_i-1)=\sum_l\zeta_l(n_l-1)$ with
$n_l\in\mathcal F^\circ\cap\ker\pi$ and $\zeta_l\in\Ftwo[\mathcal F^\circ]$.
By Fox's fundamental formula
$n_l-1=\sum_i(\partial n_l/\partial f_i)(f_i-1)$, where the
$\partial/\partial f_i$ are the free derivatives \cite[Section~2]{Fox1953},
and the coefficients of the $f_i-1$ are unique, because $\partial/\partial f_j$
maps $\sum_i\zeta'_i(f_i-1)$ to $\zeta'_j$. So
$\tilde a_i=\sum_l\zeta_l\,\partial n_l/\partial f_i$. On $\mathcal F^\circ$
the rows are the images of the vectors of free derivatives, so
$a=\sum_l\bar\zeta_l\,\partial n_l$ lies in $\Lambda\,\partial(\ker\pi)$, which
equals $\partial(\ker\pi)$ by (F3). For the remaining statements, let $M_Y$ be
the $\Lambda$-span of $\partial Y$. By (F3), the elements $n\in\ker\pi$ with
$\partial n\in M_Y$ form a closed normal subgroup of $\mathcal F$ containing
$Y$.
\end{proof}

The next lemma computes the cost of imposing a full local group. The groups
$\Gamma$ below are $C_2$, $C_2\times C_2$, $C_4$ and $D$, with standard
generators $\iota$; $\tau,\varphi$; a generator; and $x,y,z$. Let $k$ be the
number of standard generators, so $k=1,2,1,3$. The augmentation Hilbert
polynomials are $P_\Gamma(t)=1+t$, $(1+t)^2$, $(1+t)(1+t^2)$ and $P_D(t)$,
and $D_3(\Gamma)=1$ in each case. Put
\[
 \psi_\Gamma(t)=kt-1+P_\Gamma(t)^{-1}.
\]

\begin{lemma}\label{tw:filtered-local-freeness}
Let $\Gamma\le\Delta$ be a subgroup isomorphic to one of these groups, with
$g_1,\dots,g_k\in\Gamma$ corresponding to its standard generators, and
suppose that $\operatorname{gr}_n\Gamma\to\operatorname{gr}_n\Delta$ is
injective for $n=1,2$. Let $\Lambda_\Gamma=\Ftwo[\Gamma]\subset\Lambda$, with
augmentation ideal $\mathfrak I_\Gamma$, and let
$\mathcal M_\Gamma\subset\Lambda_\Gamma^k(-1)$ be the kernel of
$b\mapsto\sum_lb_l(g_l-1)$.
\begin{enumerate}
\item[(a)] There are elements $\eta_j\in\Lambda$ with weights $w_j$ such that
 $\Lambda=\bigoplus_j\eta_j\Lambda_\Gamma$ and
 $\mathfrak I^m=\bigoplus_j\eta_j\mathfrak I_\Gamma^{m-w_j}$ for every $m$,
 where $\mathfrak I_\Gamma^{m'}=\Lambda_\Gamma$ for $m'\le0$; moreover
 $\sum_jt^{w_j}=h_\Lambda/P_\Gamma$.
\item[(b)] The $\Lambda$-module $\Lambda\mathcal M_\Gamma\subset\Lambda^k(-1)$
 has $h_{\Lambda\mathcal M_\Gamma}=\psi_\Gamma h_\Lambda$.
\item[(c)] Let $\mathcal F$ and $\pi$ be as above, let
 $u_1,\dots,u_k\in\mathcal F$ with $\pi(u_l)=g_l$, and let $\mathcal F_k$ be
 free on $k$ letters, mapping to $\Gamma$ by the standard generators. Let
 $U\subset\Lambda^n(-1)$ be the $\Lambda$-span of the rows
 $\partial W(u_1,\dots,u_k)$ for all $W$ in the kernel of
 $\mathcal F_k\to\Gamma$. Then $h_U(t)\le\psi_\Gamma(t)h_\Lambda(t)$ for
 $0<t<1$.
\end{enumerate}
\end{lemma}

\begin{proof}
(a) Since $D_3(\Gamma)=1$ and $\operatorname{gr}\Gamma\to\operatorname{gr}\Delta$
is injective, $\Gamma\cap D_m\Delta=D_m(\Gamma)$ for all $m$. Choose elements
of $\Gamma$ whose classes form a homogeneous basis of $\operatorname{gr}\Gamma$,
and extend them by elements of $\Delta$ to a family whose classes form a
homogeneous basis of $\operatorname{gr}\Delta$. By the restricted
Poincar\'e--Birkhoff--Witt form of Jennings' theorem \cite{Jennings,Quillen},
the ordered products of the elements $g-1$ of this family, each used at most
once and with the elements of $\Gamma$ last, form a basis of $\Lambda$, and
those of weight at least $m$ form a basis of $\mathfrak I^m$. The same holds
for $\Lambda_\Gamma$ and the elements of $\Gamma$ alone. Let the $\eta_j$ be
the ordered products of the other elements. Comparing Hilbert series gives
the last assertion.

(b) As in the proof of Lemma~\ref{tw:fox-rows}, the map
$\Lambda_\Gamma^k(-1)\to\mathfrak I_\Gamma$ is strict and surjective, so
$h_{\mathcal M_\Gamma}=ktP_\Gamma-(P_\Gamma-1)$. By (a),
$\Lambda^k(-1)=\bigoplus_j\eta_j\Lambda_\Gamma^k(-1)$ with
$F^m=\bigoplus_j\eta_jF^{m-w_j}$, and
$\Lambda\mathcal M_\Gamma=\bigoplus_j\eta_j\mathcal M_\Gamma$. Hence
$h_{\Lambda\mathcal M_\Gamma}=(h_\Lambda/P_\Gamma)\,h_{\mathcal M_\Gamma}=\psi_\Gamma h_\Lambda$.

(c) For $W$ in the kernel, the local row $\partial'W\in\Lambda_\Gamma^k$ lies in
$\mathcal M_\Gamma$ by (F2), and by (F4)
$\partial W(u)=\sum_l\partial'_lW\,\partial u_l$. Thus $U$ is contained in the
image of $\Lambda\mathcal M_\Gamma$ under the map
$T_u\colon\Lambda^k(-1)\to\Lambda^n(-1)$, $a\mapsto\sum_la_l\partial u_l$,
which preserves filtrations because $\partial u_l\in\Lambda^n=F^0$. By (h2),
(h1) and (b),
$h_U\le h_{T_u(\Lambda\mathcal M_\Gamma)}\le h_{\Lambda\mathcal M_\Gamma}=\psi_\Gamma h_\Lambda$.
\end{proof}

For $\Gamma=D$ we write $d_D=\psi_D=3t-1+P_D(t)^{-1}$. In the proof of
Proposition~\ref{tw:infinite}, the module of the dyadic place $\mathfrak p_1$
and the modules of the other places share the row of one relator, and the
next lemma computes the module generated by this row.

\begin{lemma}\label{tw:local-fox}
In the situation of Lemma~\ref{tw:filtered-local-freeness} with $\Gamma=D$,
suppose that $u_1,u_2,u_3$ are the basis elements $f_1,f_2,f_3$ of
$\mathcal F$. Let $r\in\mathcal F_3$ map to $1$ in $D$, with class
$y^{[2]}+[x,y]+[x,z]$ in $\operatorname{gr}_2\mathcal F_3$, and let
$\rho=r(f_1,f_2,f_3)$. Then
\[
 h_{\Lambda\partial\rho}=s_Dh_\Lambda,\qquad
 s_D(t)=t^2\Bigl(1-\frac{t^7}{P_D(t)}\Bigr).
\]
\end{lemma}

\begin{proof}
By (F4), $\partial\rho=(\partial'r,0,\dots,0)$, where
$\partial'r\in\Lambda_D^3$ is the local row, so $\Lambda\partial\rho$ is the
module $\Lambda\partial'r\subset\Lambda^3(-1)$ placed in the first three
coordinates. By Lemma~\ref{tw:filtered-local-freeness}(a),
$\Lambda^3=\bigoplus_j\eta_j\Lambda_D^3$ with
$F^n(\Lambda^3(-1))=\bigoplus_j\eta_jF^{n-w_j}(\Lambda_D^3(-1))$, and
$b\mapsto\eta_jb$ is injective on $\Lambda_D^3$. The submodule
$\Lambda\partial'r=\bigoplus_j\eta_j\Lambda_D\partial'r$ decomposes in the
same way, so
$F^n(\Lambda\partial'r)=\bigoplus_j\eta_jF^{n-w_j}(\Lambda_D\partial'r)$ for
the induced filtrations, and
$h_{\Lambda\partial\rho}=(h_\Lambda/P_D)\,h_{\Lambda_D\partial'r}$. It
therefore suffices to show that $h_{\Lambda_D\partial'r}=t^2(P_D-t^7)$, for
the filtration induced from $\Lambda_D^3(-1)$. Write $\mathfrak I_D$ for the augmentation ideal of
$\Lambda_D=\Ftwo[D]$, and $X_1,X_2,X_3,X_4$ for $x-1,y-1,z-1,w-1$, where
$w=[y,z]$.

\emph{The linear part of $\partial'r$.} Here $\pi$ also denotes the map
$\mathcal F_3\to D$. For $g,h\in\mathcal F_3$, (F1) gives
$\partial'(g^2)=(1+\pi(g))\partial'g$ and
$\partial'[g,h]=\pi(g)^{-1}(\pi(h)^{-1}-1)\partial'g+\pi(g)^{-1}\pi(h)^{-1}(\pi(g)-1)\partial'h$.
Moreover $\partial'g$ is congruent modulo $\mathfrak I_D$ to the image
$\bar g\in\Ftwo^3$ of $g$ in $\mathcal F_3/D_2\mathcal F_3$, and
$\pi(g)-1\equiv\sum_l\bar g_lX_l$ modulo $\mathfrak I_D^2$ by (F2). Hence, for
$q\in D_2\mathcal F_3$, the row $\partial'q$ has entries in $\mathfrak I_D$,
and modulo $\mathfrak I_D^2$ it is additive in $q$, because
$\pi(q)-1\in\mathfrak I_D^2$. It vanishes modulo $\mathfrak I_D^2$ on
$D_3\mathcal F_3=(D_2\mathcal F_3)^2[D_2\mathcal F_3,\mathcal F_3]$ (Lazard's
formula): for $q\in D_2\mathcal F_3$ the formulas above give
$\partial'(q^2),\partial'[q,h]\in\mathfrak I_D^2$, since $\bar q=0$ and
$\pi(q)-1\in\mathfrak I_D^2$. So the row modulo $\mathfrak I_D^2$ depends
only on the class of $q$ in $\operatorname{gr}_2\mathcal F_3$, the classes
$v^{[2]}$ and $[v,v']$ giving the rows with $l$-th entries $v_l\sum_mv_mX_m$
and $v_l\sum_mv'_mX_m+v'_l\sum_mv_mX_m$.
For the class $y^{[2]}+[x,y]+[x,z]$ this gives
\[
 \partial'r\equiv(X_2+X_3,\ X_2+X_1,\ X_1)\pmod{\mathfrak I_D^2}.
\]

\emph{Three facts about $\Lambda_D$.} Let $A=\bigoplus_m\mathfrak I_D^m/\mathfrak I_D^{m+1}$,
and write $X,Y,Z$ and $W$ for the classes of $X_1,X_2,X_3$ in degree one and
of $X_4$ in degree two. By Jennings' theorem \cite{Jennings} and
Lemma~\ref{tw:local-two}(c), the products
$X_1^{a_1}X_2^{a_2}X_3^{a_3}X_4^{a_4}$ with $a_1,a_2,a_4\in\{0,1\}$ and
$0\le a_3\le3$, of weight $a_1+a_2+a_3+2a_4\ge m$, form a basis of
$\mathfrak I_D^m$; here $X_3^2=z^2-1$. Hence the monomials
$X^{a_1}Y^{a_2}Z^{a_3}W^{a_4}$ form a basis of $A$, and, since the
coefficients of $P_D$ are $1,3,5,7,7,5,3,1$,
\[
 \dim\mathfrak I_D^m=32,31,28,23,16,9,4,1,0\qquad(m=0,\dots,8).
\]
Because $x$ and $w$ are central and $x^2=y^2=w^2=z^4=1$, the elements $X$
and $W$ are central in $A$ and $X^2=Y^2=W^2=Z^4=0$. Moreover $ZY=YZ+W$,
because $X_3X_2+X_2X_3=zy+yz=zy(w-1)\equiv X_4\pmod{\mathfrak I_D^3}$; by
induction $Z^jY=YZ^j+jZ^{j-1}W$ for $j\ge1$.

First, $\mathfrak I_D^7=\Ftwo\Omega_D$, where $\Omega_D=\sum_{g\in D}g$:
$\mathfrak I_D^7$ is one-dimensional and $(g-1)\mathfrak I_D^7\subset\mathfrak I_D^8=0$,
so its nonzero element is invariant under left multiplication by $D$, and
hence equal to $\Omega_D$.

Second, for $m<7$ the map $A_m\to A_{m+1}^3$, $\xi\mapsto(\xi X,\xi Y,\xi Z)$,
is injective. Let $\xi=\sum_\mu b_\mu\mu$ over the basis monomials
$\mu=X^{a_1}Y^{a_2}Z^{a_3}W^{a_4}$ of degree $m$. Right multiplication by $X$
maps the monomials with $a_1=0$ to distinct basis monomials and kills the
others, and right multiplication by $Z$ maps those with $a_3\le2$ to
distinct basis monomials and kills those with $a_3=3$. So $\xi X=\xi Z=0$
forces $\xi=\sum_{a_2,a_4}b_{a_2a_4}XY^{a_2}Z^3W^{a_4}$. Then
\[
 \xi Y=b_{00}(XYZ^3+XZ^2W)+b_{01}XYZ^3W+b_{10}XYZ^2W,
\]
a combination of distinct basis monomials, so $\xi Y=0$ forces
$\xi=b_{11}XYZ^3W$, which has degree $7$. Hence $\xi=0$ when $m<7$.

Third, the same holds for $\xi\mapsto(\xi(Y+Z),\xi(Y+X),\xi X)$, because this
triple is obtained from $(X,Y,Z)$ by the invertible matrix
\[
 \begin{pmatrix}0&1&1\\1&1&0\\1&0&0\end{pmatrix}.
\]
This is the coefficient matrix of $y^{[2]}+[x,y]+[x,z]$, that is, the matrix
of the $2$-adic Hilbert symbol on $2,-5,5$ in the proof of
Lemma~\ref{tw:local-groups}(c). So the injectivity used here, and with it
the saving $s_D$, comes from the nondegeneracy of the local Hilbert pairing,
that is, from local duality. The supplementary program \texttt{dyadic241.py}
confirms these three facts by direct computation in $\Lambda_D$.

\emph{The filtration.} If $\xi\in\mathfrak I_D^m\setminus\mathfrak I_D^{m+1}$
with $m<7$, then by the third fact $\xi\,\partial'r$ lies in $F^{m+2}$ but not
in $F^{m+3}$, whereas $\Omega_D\partial'r=0$ because $\Omega_D\mathfrak I_D=0$.
Now let $\xi\in\Lambda_D$ with $\xi\,\partial'r\in F^{m+2}$, and suppose that
$\xi\notin\mathfrak I_D^m+\Ftwo\Omega_D$. Let $m'<m$ be maximal with
$\xi\in\mathfrak I_D^{m'}+\Ftwo\Omega_D$, and write $\xi=\xi'+a\,\Omega_D$
with $\xi'\in\mathfrak I_D^{m'}$ and $a\in\Ftwo$. Then
$\xi'\notin\mathfrak I_D^{m'+1}$, and $m'<7$ because
$\mathfrak I_D^7=\Ftwo\Omega_D$; so $\xi\,\partial'r=\xi'\partial'r$ is not in
$F^{m'+3}\supset F^{m+2}$, a contradiction. Therefore
$F^{m+2}(\Lambda_D\partial'r)=\mathfrak I_D^m\partial'r$ for every $m\ge0$,
the map $\xi\mapsto\xi\,\partial'r$ has kernel $\Ftwo\Omega_D$, and
$h_{\Lambda_D\partial'r}=\sum_{m<7}\dim(\mathfrak I_D^m/\mathfrak I_D^{m+1})\,t^{m+2}=t^2(P_D-t^7)$.
\end{proof}

\subsection{Infinitude}\label{tw:infinite-section}
Put
\begin{equation}\label{tw:costs}
 \psi_1(t)=\frac{t^2}{1+t},\quad
 \psi_2(t)=2t-1+\frac1{(1+t)^2},\quad
 \psi_4(t)=\frac{t^4}{(1+t)(1+t^2)},\quad
 d_D(t)=3t-1+\frac1{P_D(t)},
\end{equation}
so that $\psi_1,\psi_2,\psi_4,d_D$ are the functions $\psi_\Gamma$ for
$\Gamma=C_2,C_2\times C_2,C_4,D$; they are the functions used for the same
local groups in \cite[Section~2]{Naslund0418235}. With $s_D$ as in
Lemma~\ref{tw:local-fox}, let
\begin{equation}\label{tw:PB}
 P_B(t)=1-8t+2\psi_1(t)+4\psi_2(t)+2d_D(t)-s_D(t)+3\psi_4(t).
\end{equation}
Exact rational arithmetic, carried out by the supplementary program
\texttt{gs241.py}, gives
\[
 P_B\Bigl(\frac{34}{117}\Bigr)=-\frac{187433948535241}{88772225460489675}<0.
\]

\begin{proposition}\label{tw:infinite}
The group $G_B$ is infinite.
\end{proposition}

\begin{proof}
Suppose that $G_B$ is finite, and put $\Delta=G_B$ and
$\Lambda=\Ftwo[\Delta]$. By Table~\ref{tw:vector-table},
$\bar x_1,\bar y_1,\bar z_1$ are independent; choose
$g_4,\dots,g_8\in\mathcal G_B$ such that
$\bar x_1,\bar y_1,\bar z_1,\bar g_4,\dots,\bar g_8$ is a basis of
$\Ftwo^8$. Let $\mathcal F$ be free on $f_1,\dots,f_8$, and send
$f_1,f_2,f_3$ to $x_1,y_1,z_1$ and $f_l$ to $g_l$ for $4\le l\le8$. This is a
minimal presentation of
$\mathcal G_B$; we take the lifts $\hat x_1=f_1$, $\hat y_1=f_2$,
$\hat z_1=f_3$, and arbitrary lifts of the other local generators. Let
$\pi\colon\mathcal F\to\Delta$ be the composite and $\mathcal M$ as in
Lemma~\ref{tw:fox-rows}, so that $h_{\mathcal M}=1-(1-8t)h_\Lambda$.

By Lemma~\ref{tw:retained-quotient}(b), which concerns only
$\operatorname{gr}_1$ and $\operatorname{gr}_2$, the image in $\Delta$ of each
local group is the prescribed group $\Gamma$ and satisfies the hypothesis of
Lemma~\ref{tw:filtered-local-freeness}. We use eleven places: the real
places $v_1,v_2$, with $\Gamma=C_2$; the four places above $3$ and $5$, with
$\Gamma=C_2\times C_2$; the dyadic places $\mathfrak p_1,\mathfrak p_2$, with
$\Gamma=D$; and $\mathfrak t_0,\mathfrak t_1,\mathfrak t_2$, with $\Gamma=C_4$
generated by the Frobenius. For each of these places $\nu$, let $U_\nu$ be
the module $U$ of Lemma~\ref{tw:filtered-local-freeness}(c) for this $\Gamma$,
formed with the lifts of the local generators at $\nu$. Every generator of
$N_B$ in Lemma~\ref{tw:GB-presentation} is of the form $W(u_1,\dots,u_k)$ for
one of these places and some $W$ in the kernel of $\mathcal F_k\to\Gamma$, and
conversely every such element lies in $N_B$. By Lemma~\ref{tw:fox-rows},
$\mathcal M=\sum_\nu U_\nu$. Let $U'$ be the sum of the $U_\nu$ over the ten
places $\nu\ne\mathfrak p_1$. The row of $\rho_{\mathfrak p_1}$ lies in
$U_{\mathfrak p_1}$. Since $\rho_{\mathfrak p_1}$ lies in the normal closure of
the relators (i) of Lemma~\ref{tw:GB-presentation}, its row also lies in the
$\Lambda$-span of their rows (Lemma~\ref{tw:fox-rows}), which is contained in
$U'$. Hence $\Lambda\partial\rho_{\mathfrak p_1}\subset U_{\mathfrak p_1}\cap U'$.
By \eqref{tw:filtered-sum}, (h2), Lemma~\ref{tw:filtered-local-freeness}(c)
and Lemma~\ref{tw:local-fox}, for $0<t<1$,
\[
 1-(1-8t)h_\Lambda(t)=h_{\mathcal M}(t)
 \le h_{U_{\mathfrak p_1}}(t)+\sum_{\nu\ne\mathfrak p_1}h_{U_\nu}(t)
 -h_{\Lambda\partial\rho_{\mathfrak p_1}}(t)
 \le\bigl(2\psi_1+4\psi_2+2d_D+3\psi_4-s_D\bigr)(t)\,h_\Lambda(t),
\]
that is, $1\le h_\Lambda(t)P_B(t)$. Since $h_\Lambda(t)>0$, this contradicts
$P_B(34/117)<0$.
\end{proof}

\subsection{The Fields}\label{tw:fields-section}
Recall that $\lambda=2^{9/4}\sqrt{3615}$, where $3615=3\cdot5\cdot241$, and
note that $\thetastar=1/2-2^{-17}$.

\begin{theorem}\label{tw:field-family}
For every power of two $m\ge2^{49}$ there is a finite Galois extension
$K/B$ of degree $m$, contained in $B_G$, such that $G_B\to\overline G_B$
factors through $\Gal(K/B)$. In particular $K$ contains the fixed field of
$D_4G_B$, hence the fixed field of $D_3G_B$ and the Kummer field $E$. Every
such $K$ is totally imaginary. Let $\iota_1$ be complex conjugation at a
place of $K$ above $v_1$, let $F=K^{\langle\iota_1\rangle}$ and
$d=[F:\Q]=m$, and let $(b,c)$ be the signature of $F$ and $\theta=c/d$. Then
$F$ has a real place,
\[
 \rd(K)=\rd(F)=\lambda,\qquad \thetastar\le\theta<\frac12 .
\]
The extension $K/F$ is quadratic, unramified at every finite place, and
split at every place of $F$ above $2,3,5,7$ and $29$. At these primes all
places of $K$, and of $F$, have the absolute ramification indices and residue
degrees of Table~\ref{tw:types-table}.
\end{theorem}

\begin{table}[htbp]
\centering
\caption{Absolute ramification indices and residue degrees of the places of
$K$ and $F$ above the selected primes, the corresponding local extensions of
the completions of $B$, and their contributions to $\rd(K)$. The last row is
the ramified prime of $B$, at which $K/B$ is unramified.}
\label{tw:types-table}
\begin{tabular}{llll}
\toprule
$p$ & $(e_p,f_p)$ & local extension over $B$ & factor of $\rd(K)$\\
\midrule
$2$ & $(8,4)$ & $L_2/\Q_2$, group $D$ & $2^{9/4}$\\
$3$ & $(2,2)$ & $\Q_3(\sqrt{-1},\sqrt3)/\Q_3$, group $C_2\times C_2$ & $3^{1/2}$\\
$5$ & $(2,2)$ & $\Q_5(\sqrt2,\sqrt5)/\Q_5$, group $C_2\times C_2$ & $5^{1/2}$\\
$7$ & $(1,8)$ & unramified of degree $4$ over $B_{\mathfrak t_0}$ & $1$\\
$29$ & $(1,4)$ & unramified of degree $4$ over $\Q_{29}$ & $1$\\
\midrule
$241$ & $e_p=2$ & unramified & $241^{1/2}$\\
\bottomrule
\end{tabular}
\end{table}

\begin{proof}
\emph{Existence.} By Proposition~\ref{tw:infinite}, $G_B$ has open normal
subgroups of arbitrarily large index, and $D_4G_B$ is open because $G_B$ is
finitely generated. Given $m$, choose an open normal subgroup
$U'\subset D_4G_B$ of $G_B$ of index at least $m$, and put $\Delta'=G_B/U'$
and $N'=D_4G_B/U'$. A nontrivial normal subgroup of a finite $2$-group meets
its center, so $N'$ contains normal subgroups of $\Delta'$ of every order
dividing $|N'|$. Take one of order $|\Delta'|/m$, which divides
$|N'|=|\Delta'|/2^{49}$. Its preimage $U$ is open and normal in $G_B$,
contained in $D_4G_B$, and of index $m$. Let $K$ be its fixed field in $B_G$.
Conversely, every $K$ as in the statement is the fixed field of an open
normal subgroup $U\subset D_4G_B$. We fix such a $K$ and put
$H=\Gal(K/B)=G_B/U$.

\emph{Complex places.} Since $\sqrt{-1}=\sqrt{\alpha_0}\in E\subset K$, the
field $K$ is totally imaginary, and $[K:F]=2$. Let $\sigma_1\colon K\to\C$
induce the chosen place, so that $\bar\sigma_1=\sigma_1\circ\iota_1$ and
$\sigma_1(F)\subset\R$. The embeddings above $v_1$ are $\sigma_1\circ h$ with
$h\in H$, and two of them agree on $F$ exactly when they differ by right
multiplication by an element of $\langle\iota_1\rangle$. The restriction of
$\sigma_1\circ h$ to $F$ is real exactly when $\sigma_1\circ\iota_1h$ and
$\sigma_1\circ h$ agree on $F$, that is, when
$h^{-1}\iota_1h\in\langle\iota_1\rangle$, that is, when
$h\in C_H(\iota_1)$. Similarly, if $\iota'$ is complex conjugation for an
embedding $\sigma_2$ above $v_2$, the restriction of $\sigma_2\circ h$ to $F$
is real exactly when $h^{-1}\iota'h=\iota_1$. This never happens, because
conjugation preserves elementary images and $\bar\iota'=\bar\iota_2\ne\bar\iota_1$
(Lemma~\ref{tw:retained-quotient}(d)). Hence $b=|C_H(\iota_1)|/2\ge1$, so $F$
has a real place and $\theta=(1-b/d)/2<1/2$. Moreover $d=|H|$, so
$b/d=|C_H(\iota_1)|/(2|H|)$ is half the reciprocal of the size of the
conjugacy class of $\iota_1$ in $H$. This class maps onto the class of
$\iota_1$ in $\overline G_B$, which has $2^{15}$ elements by
Lemma~\ref{tw:retained-quotient}(c). Hence $b/d\le2^{-16}$ and
$\theta\ge(1-2^{-16})/2=\thetastar$.

\emph{Local structure and splitting.} Let $\mathfrak P$ be a prime of $K$
above $\mathfrak p\in S\cup\{\mathfrak t_0,\mathfrak t_1,\mathfrak t_2\}$. Its
decomposition group in $H$ is conjugate to the image of the local map at
$\mathfrak p$. Since $U\subset D_3G_B$, Lemma~\ref{tw:retained-quotient}(b)
shows that this image is $D$, $C_2\times C_2$ or $C_4$, the last one cyclic on
the Frobenius, and that the prescribed local group injects into $H$.
Therefore the completion $K_{\mathfrak P}$ is $L_2$ over
$B_{\mathfrak p_j}=\Q_2$, the field $\Q_p(\sqrt u,\sqrt p)$ over
$B_{\mathfrak p}=\Q_p$ for $p=3,5$, and the unramified extension of degree
four of $B_{\mathfrak t_k}$. Since $B_{\mathfrak t_0}$ is unramified of
degree two over $\Q_7$ and $B_{\mathfrak t_1}=B_{\mathfrak t_2}=\Q_{29}$, the
places of $K$ have the types of Table~\ref{tw:types-table}. Every element of
the decomposition group has elementary image in the span of the images of the
local generators at $\mathfrak p$, whereas every conjugate of $\iota_1$ has
elementary image $\bar\iota_1$. By Lemma~\ref{tw:retained-quotient}(d),
$\iota_1$ lies in no decomposition group of a prime above
$S\cup\{\mathfrak t_0,\mathfrak t_1,\mathfrak t_2\}$, so the decomposition
group of $\mathfrak P$ in $\Gal(K/F)=\langle\iota_1\rangle$ is trivial. Hence
the places of $F$ above $2,3,5,7,29$ split in $K$, and they have the same
local types.

\emph{Ramification and root discriminant.} The extension $K/B$, and hence
$K/F$, is unramified at every prime outside $S$, in particular at
$(\sqrt{241})$. At primes above $S$ the inertia group of $K/F$ is trivial,
because it lies in the decomposition group. So $K/F$ is unramified at every
finite place, its relative discriminant is trivial, $|d_K|=|d_F|^2$, and
$\rd(K)=\rd(F)$. To compute $\rd(K)$ we use
$|d_K|=241^{[K:B]}N\mathfrak d_{K/B}$. Since $K/B$ is Galois, for
$\mathfrak p\in S$ all primes $\mathfrak P$ above $\mathfrak p$ have the same
ramification index $e$, residue degree and different exponent
$v_{\mathfrak P}(\mathfrak D_{K/B})$, and
$v_{\mathfrak p}(\mathfrak d_{K/B})=[K:B]\,v_{\mathfrak P}(\mathfrak D_{K/B})/e$.
For $\mathfrak p$ above $3$ or $5$ we have $e=2$ and
$v_{\mathfrak P}(\mathfrak D_{K/B})=1$, since over the unramified extension
$\Q_p(\sqrt u)$ the different of $\Q_p(\sqrt u,\sqrt p)$ is generated by
$2\sqrt p$. For $\mathfrak p_1,\mathfrak p_2$, Lemma~\ref{tw:local-two}(d)
gives $v_{\mathfrak P}(\mathfrak D_{K/B})/e=9/4$. Since $[K:\Q]=2[K:B]$, the
base contributes $241^{1/2}$, each of the two primes above $3$ contributes
$3^{1/4}$, each of the two above $5$ contributes $5^{1/4}$, and each dyadic
prime contributes $2^{9/8}$. Hence
$\rd(K)=2^{9/4}\,3^{1/2}\,5^{1/2}\,241^{1/2}=\lambda$.
\end{proof}

\begin{remark}\label{tw:lean-remark}
The Lean development \cite{NaslundLean2026} formalizes a version of this
construction; Remark~\ref{an:lean-hypothesis} states what it proves and the
numerical hypothesis it assumes.
\end{remark}

\section{\texorpdfstring{The Kummer Field and Its Quadratic $L$-Functions}{The Kummer Field and Its Quadratic L-Functions}}\label{gf:section}

The analytic estimate of Section~\ref{an:section} compares every field of
Theorem~\ref{tw:field-family} with one fixed subfield, the \emph{Kummer
field}
\[
 E=B(\sqrt{\alpha_0},\ldots,\sqrt{\alpha_7})=B(\sqrt V)
\]
of Subsection~\ref{tw:global-section}, with $\alpha_0,\ldots,\alpha_7$ as in
\eqref{tw:alpha} and $V$ as in Lemma~\ref{tw:base}(d). It is the maximal
elementary abelian $2$-extension of $B$ unramified at the finite primes
outside $S$, that is, the fixed field of the Frattini subgroup of $G_B$
(Lemma~\ref{gf:in-tower}); the Lean formalization calls it the genus field.
This section shows that $E$ lies in every field of
Theorem~\ref{tw:field-family}. It then writes $\zeta_E$ as the product of
$\zeta_B$ and $255$ quadratic Hecke $L$-functions of $B$, and determines
their conductors, gamma factors, root numbers and Euler factors.

For $v,e\in\Ftwo^8$ put $\langle v,e\rangle=\sum_{i=0}^7v_ie_i\in\Ftwo$; as
in Section~\ref{tw:tower}, vectors are written as words $v_0v_1\cdots v_7$.
For $e\in\Ftwo^8$ put
\[
 \alpha_e=\prod_{i=0}^7\alpha_i^{e_i},\qquad B_e=B(\sqrt{\alpha_e}).
\]
For a nonzero ideal $\mathfrak a$ of a number field, $\mathrm N\mathfrak a$
denotes its absolute norm. The number of positive divisors of an integer
$n\ge1$ is $\tau(n)$.

\subsection{The Characters of the Kummer Field}
By Subsection~\ref{tw:global-section}, $[E:B]=256$, and the elementary image
$\bar g\in\Ftwo^8$, defined by $g(\sqrt{\alpha_i})=(-1)^{\bar
g_i}\sqrt{\alpha_i}$, identifies $\Gal(E/B)$ with $\Ftwo^8$. For
$e\in\Ftwo^8$ let $\chi_e$ be the character $g\mapsto(-1)^{\langle\bar
g,e\rangle}$ of $\Gal(E/B)$. Since $g(\sqrt{\alpha_e})=\chi_e(g)\sqrt{\alpha_e}$,
the character $\chi_e$ cuts out $B_e$, and the $256$ characters $\chi_e$ are
all the characters of $\Gal(E/B)$. For $e\ne0$ the extension $B_e/B$ is
quadratic, because the classes of $\alpha_0,\ldots,\alpha_7$ are independent
in $B^\times/B^{\times2}$ (Lemma~\ref{tw:base}(d)).

For a pro-$2$ group $H$, the Frattini subgroup $\Phi(H)=H^2[H,H]$ is $D_2H$.

\begin{lemma}\label{gf:in-tower}
The fixed field of $D_2G_B=\Phi(G_B)$ is $E$. Consequently $E$ is contained
in every field $K$ of Theorem~\ref{tw:field-family}.
\end{lemma}

\begin{proof}
Let $\mathcal G_B$ be the group of Subsection~\ref{tw:global-section}, so
that $\mathcal G_B/\Phi(\mathcal G_B)=\Gal(E/B)$, and let
$G_B=\mathcal G_B/N$ as in Definition~\ref{tw:GB}. As observed after that
definition, $N\subseteq\Phi(\mathcal G_B)$: the images of the Frattini
subgroups of the local groups above $3$ and $5$ and the fourth powers of
Frobenius elements lie in $\Phi(\mathcal G_B)$, and the kernel of the
surjection $\mathcal G_{\mathfrak p_j}\to D$ lies in $\Phi(\mathcal
G_{\mathfrak p_j})$ because both groups have generator rank three. Hence
$G_B/\Phi(G_B)=\mathcal G_B/\Phi(\mathcal G_B)=\Gal(E/B)$. A field $K$ of
Theorem~\ref{tw:field-family} contains the fixed field of
$D_4G_B\subseteq D_2G_B$, and hence contains $E$.
\end{proof}

\subsection{Local Data}
For a prime $\mathfrak p\notin S$ of $B$, the \emph{Frobenius vector}
$v_{\mathfrak p}\in\Ftwo^8$ is the elementary image of the Frobenius of
$\mathfrak p$ in $\Gal(E/B)$. By Euler's criterion, $(v_{\mathfrak p})_i=1$
exactly when $\alpha_i$ is not a square in $\mathcal O_B/\mathfrak p$. For a
place $\mathfrak p$ of $B$ we write $(\cdot,\cdot)_{\mathfrak p}$ for the
quadratic Hilbert symbol of $B_{\mathfrak p}$.

\begin{proposition}\label{gf:local-data}
Let $e\in\Ftwo^8\setminus\{0\}$. For a prime $\mathfrak p$ of $B$ let
$a_{\mathfrak p}(e)$ be the exponent of $\mathfrak p$ in the conductor of
$\chi_e$, and if $a_{\mathfrak p}(e)=0$, let $\chi_e(\mathfrak p)\in\{\pm1\}$
be the value of $\chi_e$ at the Frobenius of $\mathfrak p$.
\begin{enumerate}
\item If $\mathfrak p\notin S$, then $a_{\mathfrak p}(e)=0$ and
 $\chi_e(\mathfrak p)=(-1)^{\langle v_{\mathfrak p},e\rangle}$. Let $p$ be the
 rational prime below $\mathfrak p$. If $p$ splits in $B$, then
 $\mathfrak p=(p,\sqrt{241}-r)$ with $r^2\equiv241\pmod p$, and
 $(v_{\mathfrak p})_i=1$ exactly when the image of $\alpha_i$ under
 $\sqrt{241}\mapsto r$ is a quadratic nonresidue modulo $p$. If $p$ is
 inert, then $(v_{\mathfrak p})_i=1$ exactly when $N_{B/\Q}(\alpha_i)$ is a
 quadratic nonresidue modulo $p$. If $p=241$, the same holds with
 $\sqrt{241}\mapsto0$.
\item Let $\mathfrak p$ lie above $p\in\{3,5\}$, let $u=-1$ if $p=3$ and $u=2$
 if $p=5$, and let $\varpi=-\alpha_i$, where $\alpha_i$ generates
 $\mathfrak p$. If $(u,\alpha_e)_{\mathfrak p}=-1$, then
 $a_{\mathfrak p}(e)=1$. Otherwise $a_{\mathfrak p}(e)=0$ and
 $\chi_e(\mathfrak p)=(\varpi,\alpha_e)_{\mathfrak p}$.
\item Let $\mathfrak p=\mathfrak p_j$. If $(5,\alpha_e)_{\mathfrak p}=-1$, then
 $a_{\mathfrak p}(e)=3$. If $(5,\alpha_e)_{\mathfrak p}=1$ and
 $(-1,\alpha_e)_{\mathfrak p}=-1$, then $a_{\mathfrak p}(e)=2$. Otherwise
 $a_{\mathfrak p}(e)=0$ and $\chi_e(\mathfrak p)=(-2,\alpha_e)_{\mathfrak p}$.
\item The real place $v_k$ becomes complex in $B_e$, that is, $\chi_e$ is
 nontrivial at $v_k$, exactly when $\alpha_e<0$ at $v_k$.
\end{enumerate}
\end{proposition}

\begin{proof}
For (1), $\mathfrak p$ is odd and $\alpha_e$ is a unit at $\mathfrak p$, so
$\mathfrak p$ is unramified in $B_e$. By Euler's criterion the Frobenius
multiplies $\sqrt{\alpha_e}$ by the quadratic residue symbol of $\alpha_e$
modulo $\mathfrak p$, which is $\prod_i(\alpha_i/\mathfrak p)^{e_i}$. For a
split prime, $\mathcal O_B/\mathfrak p\cong\mathbb F_p$ via $\sqrt{241}\mapsto
r$. For an inert prime, $\mathcal O_B/\mathfrak p\cong\mathbb F_{p^2}$ and
the conjugation of $B$ induces the Frobenius $x\mapsto x^p$ of the residue
field, so the residue of $N_{B/\Q}(\alpha)$ is $\bar\alpha^{p+1}$. An element
of $\mathbb F_{p^2}^\times$ is a square exactly when
$\bar\alpha^{(p^2-1)/2}=(\bar\alpha^{p+1})^{(p-1)/2}$ equals $1$, that is,
when its norm is a square in $\mathbb F_p$. The prime above $241$ has
residue field $\mathbb F_{241}$ and $\sqrt{241}\mapsto0$.

For (2) and (3) we use local class field theory for $B_{\mathfrak p}=\Q_p$.
The local Artin map sends $b\in\Q_p^\times$ to the automorphism multiplying
$\sqrt{\alpha_e}$ by $(b,\alpha_e)_{\mathfrak p}$. Units map onto inertia, a
uniformizer maps to a Frobenius, and the conductor exponent is the least
$n\ge0$ such that $b\mapsto(b,\alpha_e)_{\mathfrak p}$ is trivial on
$U^{(n)}$, where $U^{(0)}=\mathbb Z_p^\times$ and $U^{(n)}=1+p^n\mathbb Z_p$
\cite{MilneCFT}.

For $p=3,5$, the unit group modulo squares is $\{1,u\}$, and $U^{(1)}$
consists of squares by Hensel's lemma. Hence $a_{\mathfrak p}(e)\le1$, with
equality exactly when $(u,\alpha_e)_{\mathfrak p}=-1$. In the unramified
case the Frobenius is the image of any uniformizer, such as $\varpi$.

For $p=2$, the squares of $\mathbb Z_2^\times$ form $U^{(3)}$. Moreover
$U^{(2)}=U^{(3)}\cup5U^{(3)}$ and
$U^{(0)}=U^{(1)}=\{\pm1,\pm5\}U^{(3)}$. The character is thus trivial on
$U^{(3)}$. It is trivial on $U^{(2)}$ exactly when $(5,\alpha_e)=1$, and
on $U^{(0)}$ exactly when also $(-1,\alpha_e)=1$. This gives the stated
exponents; the exponent $1$ cannot occur since $U^{(0)}=U^{(1)}$. In the
unramified case the Frobenius is the image of the uniformizer $-2$.

For (4), the real place $v_k$ splits in $B_e$ exactly when $\alpha_e>0$ at
$v_k$.
\end{proof}

Each condition in the proposition is linear in $e$. Indeed,
$(b,\alpha_e)_{\mathfrak p}=(-1)^{\langle w,e\rangle}$, where $w$ is the
elementary image of the local generator that acts on square roots by
$(b,\cdot)_{\mathfrak p}$ in Table~\ref{tw:vector-table}: at $\mathfrak p_j$
the vectors $\bar x_j,\bar y_j,\bar z_j$ belong to $b=5,-1,-2$, and above
$3$ and $5$ the images of the inertia and Frobenius generators belong to
$b=u$ and $b=\varpi$. Likewise $\alpha_e<0$ at $v_k$ exactly when
$\langle\bar\iota_k,e\rangle=1$, where $\bar\iota_k$ is the elementary image of
the complex conjugation $\iota_k$. We put
$\nu_k(e)=\langle\bar\iota_k,e\rangle\in\{0,1\}$. The conductor of $\chi_e$
is $\mathfrak f_e=\prod_{\mathfrak p\in S}\mathfrak p^{a_{\mathfrak p}(e)}$.

\begin{corollary}\label{gf:types}
The ramification index and residue degree of a prime $\mathfrak p$ of $B$ in
$E/B$ are $(4,2)$ at $\mathfrak p_1,\mathfrak p_2$; $(2,2)$ at
$\mathfrak q_1,\mathfrak q_2,\mathfrak r_1,\mathfrak r_2$; and $(1,1)$ or
$(1,2)$ at $\mathfrak p\notin S$, according as $v_{\mathfrak p}=0$ or not.
In particular $\mathfrak t_0,\mathfrak t_1,\mathfrak t_2$ have type $(1,2)$.
\end{corollary}

\begin{proof}
For $\mathfrak p\in S$, local class field theory identifies the
decomposition group of $\mathfrak p$ in $\Gal(E/B)=\Ftwo^8$ with the image
of $\Q_p^\times$ under the local Artin map, and the inertia group with the
image of the units. By the preceding paragraph, at $\mathfrak p_j$ these are
spanned by $\bar x_j,\bar y_j,\bar z_j$ and by $\bar x_j,\bar y_j$; the three
vectors are independent (Table~\ref{tw:vector-table}), so the decomposition
group has order $8$ and inertia has order $4$. Above $3$ and $5$ the two
elementary images are independent, which gives $(2,2)$. For
$\mathfrak p\notin S$ the decomposition group is generated by
$v_{\mathfrak p}$. The Frobenius vectors of $\mathfrak t_0,\mathfrak
t_1,\mathfrak t_2$ are nonzero (Table~\ref{tw:vector-table}).
\end{proof}

\subsection{The Zeta Function of the Kummer Field}
Let $\chi_{241}$ be the Dirichlet character $n\mapsto(n/241)$.

\begin{proposition}\label{gf:factorization}
For $\operatorname{Re}s>1$,
\[
 \zeta_E(s)=\zeta_B(s)\prod_{e\in\Ftwo^8\setminus\{0\}}L(s,\chi_e),
 \qquad \zeta_B(s)=\zeta(s)L(s,\chi_{241}),
\]
where $L(s,\chi_e)=\zeta_{B_e}(s)/\zeta_B(s)$ is the Hecke $L$-function of
$\chi_e$.
\end{proposition}

\begin{proof}
Compare Euler factors at a prime $\mathfrak p$ of $B$. Let $T_{\mathfrak
p}\subseteq Z_{\mathfrak p}\subseteq G=\Gal(E/B)$ be its inertia and
decomposition groups, and put $f_{\mathfrak p}=|Z_{\mathfrak p}/T_{\mathfrak
p}|$ and $g_{\mathfrak p}=|G/Z_{\mathfrak p}|$. The Euler factor of $\zeta_E$
at $\mathfrak p$ is $(1-X^{f_{\mathfrak p}})^{-g_{\mathfrak p}}$ with
$X=\mathrm N\mathfrak p^{-s}$. A character nontrivial on $T_{\mathfrak p}$
contributes the factor $1$. The characters trivial on $T_{\mathfrak p}$ are
those of $G/T_{\mathfrak p}$; they restrict onto the $f_{\mathfrak p}$
characters of the cyclic group $Z_{\mathfrak p}/T_{\mathfrak p}$, each with
$g_{\mathfrak p}$ extensions. Hence their factors multiply to
$\prod_{\zeta^{f_{\mathfrak p}}=1}(1-\zeta X)^{-g_{\mathfrak p}}
=(1-X^{f_{\mathfrak p}})^{-g_{\mathfrak p}}$. The same identity for $B_e/B$
identifies $L(s,\chi_e)$ with $\zeta_{B_e}/\zeta_B$. The second identity is
the case of $B/\Q$. Since $241\equiv1\pmod8$, the prime $2$ splits in $B$,
and quadratic reciprocity shows that the character of $B/\Q$ is
$\chi_{241}$.
\end{proof}

Put $\Gamma_\R(s)=\pi^{-s/2}\Gamma(s/2)$ and
$\Gamma_\C(s)=2(2\pi)^{-s}\Gamma(s)$, so that
$\Gamma_\C(s)=\Gamma_\R(s)\Gamma_\R(s+1)$ by the duplication formula. For a
number field $M$ with discriminant $\Delta_M$ and signature
$(r_1(M),r_2(M))$, the completed Dedekind zeta function is
\[
 \Lambda_M(s)=|\Delta_M|^{s/2}\Gamma_\R(s)^{r_1(M)}\Gamma_\C(s)^{r_2(M)}
 \zeta_M(s).
\]

\begin{proposition}\label{gf:degree-two}
Let $e\ne0$, write $L(s,\chi_e)=\sum_{n\ge1}a_n(e)n^{-s}$, and put
\[
 Q_e=241\,\mathrm N\mathfrak f_e,\qquad
 \Lambda(s,\chi_e)=Q_e^{s/2}\,\Gamma_\R(s+\nu_1(e))\,\Gamma_\R(s+\nu_2(e))\,
 L(s,\chi_e).
\]
\begin{enumerate}
\item The coefficients $a_n(e)$ are integers with $|a_n(e)|\le\tau(n)$. The
 Euler factor at a rational prime $p$ is
 $\prod_{\mathfrak p\mid p}(1-\chi_e(\mathfrak p)\mathrm N\mathfrak p^{-s})^{-1}$,
 with $\chi_e(\mathfrak p)=0$ when $a_{\mathfrak p}(e)>0$.
\item The function $\Lambda(s,\chi_e)$ equals $\Lambda_{B_e}(s)/\Lambda_B(s)$.
 It is entire of order at most one and satisfies
 $\Lambda(s,\chi_e)=\Lambda(1-s,\chi_e)$.
\item Exactly $63$ of the characters have $\nu_1(e)=\nu_2(e)=0$, exactly
 $64$ have $\nu_1(e)=\nu_2(e)=1$, and $128$ have $\nu_1(e)\ne\nu_2(e)$.
\end{enumerate}
\end{proposition}

\begin{proof}
The Euler factor in (1) is Proposition~\ref{gf:factorization} applied prime
by prime. Above a rational prime there are at most two primes of $B$. So
the local factor is either a product of at most two factors
$(1-cp^{-s})^{-1}$, or a single factor $(1-cp^{-2s})^{-1}$, with
$c\in\{0,\pm1\}$. In both cases the coefficient of $p^{-ks}$ has modulus at
most $k+1=\tau(p^k)$. Multiplicativity gives (1).

For (2), the completed Dedekind zeta functions satisfy
$\Lambda_M(1-s)=\Lambda_M(s)$, and $L(s,\chi_e)$ extends to an entire
function because $\chi_e$ is nontrivial \cite{Hecke1920,Tate1967}. The
discriminants satisfy
$|\Delta_{B_e}|=|\Delta_B|^2\,\mathrm N\mathfrak d_{B_e/B}$, where
$\mathfrak d_{B_e/B}$ is the relative discriminant, and the
conductor--discriminant formula gives $\mathfrak d_{B_e/B}=\mathfrak f_e$
for the quadratic extension $B_e/B$. Hence
$|\Delta_{B_e}|/|\Delta_B|=241\,\mathrm N\mathfrak f_e=Q_e$. If $\nu_k(e)=0$,
the place $v_k$ splits into two real places of $B_e$ and contributes
$\Gamma_\R(s)^2/\Gamma_\R(s)=\Gamma_\R(s)$. If $\nu_k(e)=1$, it becomes one
complex place and contributes $\Gamma_\C(s)/\Gamma_\R(s)=\Gamma_\R(s+1)$.
Thus $\Lambda_{B_e}/\Lambda_B=\Lambda(s,\chi_e)$. It is the completed Hecke
$L$-function of $\chi_e$, which is entire of order at most one
\cite{Hecke1920,Tate1967}.

For (3), the vectors $\bar\iota_1,\bar\iota_2$ are linearly independent. So
$e\mapsto(\nu_1(e),\nu_2(e))$ maps $\Ftwo^8$ onto $\Ftwo^2$ with fibers of
size $64$. Removing $e=0$ from the fiber over $(0,0)$ gives the counts.
\end{proof}

The characters with $\nu_1=\nu_2$ have \emph{pure} gamma factors
$\Gamma_\R(s+\nu)^2$, and the others have \emph{mixed} gamma factors
$\Gamma_\R(s)\Gamma_\R(s+1)=\Gamma_\C(s)$. Only the primes of $S$ divide the
conductors. Since $\mathrm N\mathfrak f_e=2^{a_{\mathfrak p_1}(e)+a_{\mathfrak
p_2}(e)}3^{a_{\mathfrak q_1}(e)+a_{\mathfrak q_2}(e)}5^{a_{\mathfrak
r_1}(e)+a_{\mathfrak r_2}(e)}$, the conductors are $Q_e=241\cdot2^i\cdot m$
with $2^i\in\{1,4,8,16,32,64\}$ and $m\mid225$. The number of characters with
given $2^i$ and $m$ is as follows.
\begin{equation}\label{gf:conductor-table}
\begin{array}{c|ccccccccc|c}
 2^i\backslash m&1&3&5&9&15&25&45&75&225&\text{total}\\\hline
 1&0&2&2&1&4&1&2&2&1&15\\
 4&2&4&4&2&8&2&4&4&2&32\\
 8&4&8&8&4&16&4&8&8&4&64\\
 16&1&2&2&1&4&1&2&2&1&16\\
 32&4&8&8&4&16&4&8&8&4&64\\
 64&4&8&8&4&16&4&8&8&4&64\\\hline
 \text{total}&15&32&32&16&64&16&32&32&16&255
\end{array}
\end{equation}
No nontrivial character has $\mathrm N\mathfrak f_e=1$, since the narrow
class number of $B$ is one. The conductors range from $Q_e=723$ to
$Q_e=241\cdot14400=3470400$.

\begin{remark}\label{gf:checks}
The supplementary program \texttt{lfun241.py} computes, for each $e\ne0$,
the conductor $Q_e$, the gamma type and the coefficients $a_n(e)$ for
$n\le N_e=4\lfloor\sqrt{Q_e}\rfloor+1$ from Propositions~\ref{gf:local-data}
and~\ref{gf:degree-two}(1); these are the data used in
Section~\ref{an:section}. The supplementary program \texttt{check255.gp}
compares all $255$ of them with PARI/GP's own number fields and
$L$-functions \cite{PARI}. For each $e$ it computes the maximal order of
$B_e$, checks $|\Delta_{B_e}|=241\,Q_e$ and the gamma type against the
signature of $B_e$, and checks every coefficient $a_n(e)$ with $n\le N_e$
against the Dirichlet series of $\zeta_{B_e}/\zeta_B$. All agree. It also
prints the table \eqref{gf:conductor-table}, the counts of
Proposition~\ref{gf:degree-two}(3) and the range of $Q_e$. These comparisons
check the implementation; the identification of the $L$-functions rests on
Propositions~\ref{gf:local-data} and~\ref{gf:degree-two}.
\end{remark}

\section{The Relative Zeta Value and an Explicit Upper Bound}\label{an:section}

We bound the relative zeta value $L_F(1)$ needed in the geometric
construction. Logarithms of Dedekind zeta functions are normalized by the
degree: for a number field $A$ we work with $[A:\Q]^{-1}\log\zeta_A(s)$.

\begin{definition}\label{an:family}
Let $\mathcal K$ be the set of all finite Galois extensions $K/B$ contained
in $B_G$, of degree a power of two at least $2^{49}$, such that
$G_B\to\overline G_B$ factors through $\Gal(K/B)$ (notation of
Definition~\ref{tw:GB}). These are the fields $K$ of
Theorem~\ref{tw:field-family}.
\end{definition}

For $K\in\mathcal K$ we use the notation of
Theorem~\ref{tw:field-family}: $\iota_1$ is a complex conjugation at a place of $K$ above
$v_1$, $F=K^{\langle\iota_1\rangle}$, $d=[F:\Q]$, so that $[K:\Q]=2d$, $(b,c)$
is the signature of $F$, $\theta=c/d$, and $\lambda=2^{9/4}\sqrt{3615}$ is the
root discriminant of both $K$ and $F$. Put
\[
 \ell=\log\lambda=\tfrac94\log2+\tfrac12\log3615=5.65600472355\ldots
\]
Different choices of $\iota_1$ are conjugate in $\Gal(K/B)$ and give
conjugate fields $F$; hence $b$, $c$, $\theta$ and the function
$L_F(s)=\zeta_K(s)/\zeta_F(s)$ depend only on $K$. The extension $K/F$ is
quadratic and unramified at every finite place, $K$ is totally imaginary,
and $b\ge1$. Let $\gamma$ be Euler's constant. For a
prime $\mathfrak p$ of $B$, let $e_K(\mathfrak p)$ and $f_K(\mathfrak p)$ be
its ramification index and residue degree in the Galois extension $K/B$, so
that $(e_K(\mathfrak p),f_K(\mathfrak p))$ is the type of $\mathfrak p$ in
$K/B$ (Definition~\ref{tw:type}).

\begin{theorem}\label{an:ceiling}
There is an $m_0$ such that every $K\in\mathcal K$ with $[K:B]\ge m_0$
satisfies
\[
 \frac1d\log L_F(1)<C:=\ceiling.
\]
\end{theorem}

The proof assumes neither the generalized Riemann hypothesis nor a
Brauer--Siegel type asymptotic formula, and it does not give an effective
value of $m_0$. It has three steps. Monotonicity of the completion of $L_F$
reduces the value at $s=1$ to $\zeta_K$ at a real point $\sigma>1$
(Subsection~\ref{an:reduction}). The Tsfasman--Vl\u{a}du\c{t} inequality
bounds the error of this shift (Subsection~\ref{an:budget-section},
Proposition~\ref{an:prime-budget}). Finally, $\zeta_K(\sigma)$ is bounded
prime by prime by $\zeta_E(\sigma)$, for the Kummer field $E$ of
Section~\ref{gf:section}, minus corrections at the primes above $2$, $7$
and $29$, whose types in $K/B$ are known, and at the primes of
$\mathcal P_4$ (Definition~\ref{an:P4}), whose residue degree in $K/B$ is
at least $4$ (Proposition~\ref{an:Ystar}); $\zeta_E(\sigma)$ is evaluated
rigorously in
Subsections~\ref{an:afe-section}--\ref{an:value-section}
(Proposition~\ref{an:genus-value}). Subsection~\ref{an:census} computes
these corrections and the lower bounds for the types in $K/B$ that
Proposition~\ref{an:prime-budget} uses, and the last subsection combines
the steps. Throughout, $\sigma$ denotes a real number greater than $1$ and
$\varepsilon=\sigma-1$; in this section $\varepsilon$ is this positive
number, not the unit $\alpha_1$ of \eqref{tw:alpha}.

\subsection{\texorpdfstring{Reduction to $s>1$}{Reduction to s > 1}}
\label{an:reduction}
For a number field $A$, write $h_A$, $\operatorname{Reg}_A$, $w_A$ and
$\Delta_A$ for its class number, regulator, number of roots of unity and
discriminant. The residue of $\zeta_A$ at $s=1$ is
$2^{r_1(A)}(2\pi)^{r_2(A)}h_A\operatorname{Reg}_A/(w_A|\Delta_A|^{1/2})$.
Since $K/F$ is unramified at the finite places, $|\Delta_K|=|\Delta_F|^2$.
Hence
\begin{equation}\label{an:class-formula}
 L_F(1)=\lim_{s\to1}\frac{\zeta_K(s)}{\zeta_F(s)}
 =2^c\pi^{d-c}\lambda^{-d/2}\,
 \frac{h_K\operatorname{Reg}_Kw_F}{h_F\operatorname{Reg}_Fw_K}>0.
\end{equation}

Let $s>1$ be real. At a prime of $F$ of norm $q$, the Euler factor of $L_F$
is $(1-q^{-s})^{-1}$ if the prime splits in $K$ and $(1+q^{-s})^{-1}$ if it
is inert. Its square is at most the Euler factor of $\zeta_K$, which is
$(1-q^{-s})^{-2}$, respectively $(1-q^{-2s})^{-1}$. Consequently
\begin{equation}\label{an:relative-euler}
 \frac{\log L_F(s)}d\le\frac{\log\zeta_K(s)}{2d}
 =\frac{\log\zeta_K(s)}{[K:\Q]}.
\end{equation}

The function $L_F$ is the Hecke $L$-function of the quadratic character of
$K/F$. This character is unramified at every finite place and nontrivial at
every real place of $F$, since $K$ is totally imaginary. With $\Gamma_\R$,
$\Gamma_\C$ and $\Lambda_A$ as in Section~\ref{gf:section}, put
\begin{equation}\label{an:completion}
 \mathcal L_F(s)=\frac{\Lambda_K(s)}{\Lambda_F(s)}
 =\lambda^{ds/2}\,\Gamma_\R(s+1)^b\,\Gamma_\C(s)^c\,L_F(s).
\end{equation}
By Hecke's theorem \cite{Hecke1920,Tate1967}, $\mathcal L_F$ is entire of
order at most one, and it satisfies $\mathcal L_F(s)=\mathcal L_F(1-s)$
because both completed Dedekind zeta functions do.

\begin{lemma}\label{an:monotonicity}
The function $\mathcal L_F(s)$ is positive and nondecreasing on $[1,\infty)$.
\end{lemma}

\begin{proof}
The Euler product shows that $\mathcal L_F$ has no zeros with
$\operatorname{Re}s>1$, and the functional equation excludes
$\operatorname{Re}s<0$. So every zero $\rho$ satisfies $0\le\operatorname{Re}
\rho\le1$. There is no zero on $[1,\infty)$, by the Euler product and by
\eqref{an:class-formula}. As $\mathcal L_F(s)>0$ for large real $s$,
it is positive on $[1,\infty)$.

Put $\Xi(z)=\mathcal L_F(\frac12+z)$. It is entire of order at most one,
even, and real on the real axis. Its zeros $z=\rho-\frac12$ satisfy
$|\operatorname{Re}z|\le\frac12$, and $\sum_{z\ne0}|z|^{-2}<\infty$. Group
the nonzero zeros into pairs $\{z_0,-z_0\}$, and let $m\ge0$ be the order of
$\Xi$ at $0$. By Hadamard's theorem for functions of order at most one,
\[
 \frac{\Xi'}{\Xi}(z)-\frac mz-\sum_{\{z_0,-z_0\}}\frac{2z}{z^2-z_0^2}
\]
is constant, where the series converges absolutely and locally uniformly
away from the zeros. Since $\Xi$ is even, $\Xi'/\Xi$ is odd, so the constant
is zero. For real $x\ge\frac12$ the logarithmic derivative is real, so it
equals the sum of the real parts. If $z_0=a+iy$, then
\[
 \operatorname{Re}\frac{2x}{x^2-z_0^2}
 =\frac{2x(x^2-a^2+y^2)}{(x^2-a^2+y^2)^2+4a^2y^2}\ge0,
\]
because $|a|\le\frac12\le x$. Since also $m/x\ge0$, the logarithmic
derivative is nonnegative, and $\log\Xi$ is nondecreasing on
$[\frac12,\infty)$.
\end{proof}

For $0<\varepsilon\le1$ and $0\le\theta\le\frac12$ put
\begin{equation}\label{an:gamma}
 \Upsilon(\varepsilon,\theta)
 =\frac\varepsilon2(\ell-\log\pi)-\varepsilon\theta\log2
 +(1-2\theta)\log\Gamma\Bigl(1+\frac\varepsilon2\Bigr)
 +\theta\log\Gamma(1+\varepsilon).
\end{equation}

\begin{lemma}\label{an:shift}
Let $K\in\mathcal K$ and $0<\varepsilon\le1$. Then
\begin{equation}\label{an:finite-gamma}
 \frac{\log L_F(1)}d
 \le\frac{\log\zeta_K(1+\varepsilon)}{[K:\Q]}+\Upsilon(\varepsilon,\theta).
\end{equation}
Moreover $\Upsilon(\varepsilon,\theta)\to0$ as $\varepsilon\downarrow0$,
uniformly for $0\le\theta\le\frac12$.
\end{lemma}

\begin{proof}
By Lemma~\ref{an:monotonicity},
$\mathcal L_F(1)\le\mathcal L_F(1+\varepsilon)$. In \eqref{an:completion}
the gamma factors change by
\[
 \frac{\Gamma_\R(2+\varepsilon)}{\Gamma_\R(2)}
 =\pi^{-\varepsilon/2}\,\Gamma\Bigl(1+\frac\varepsilon2\Bigr),\qquad
 \frac{\Gamma_\C(1+\varepsilon)}{\Gamma_\C(1)}
 =(2\pi)^{-\varepsilon}\,\Gamma(1+\varepsilon).
\]
Taking logarithms, dividing by $d$, and using $b/d=1-2\theta$, $c/d=\theta$
and \eqref{an:relative-euler}, we obtain \eqref{an:finite-gamma}. Since
$\Upsilon$ is affine in $\theta$, $\Upsilon(\varepsilon,\theta)\to0$ as
$\varepsilon\downarrow0$, uniformly for $0\le\theta\le\frac12$.
\end{proof}

The term $\Upsilon$ does not appear in the final estimate. With the value
$\varepsilon=1/300$ used in the proof of Theorem~\ref{an:ceiling},
\eqref{an:finite-gamma} would add $\Upsilon(1/300,\theta)\approx0.0054$ for
$\theta$ near $\frac12$. Instead, the proof of
Proposition~\ref{an:prime-budget} applies \eqref{an:finite-gamma} with
$s-1$ in place of $\varepsilon$ for each fixed $s\in(1,2]$, passes to the
limit along a sequence of fields, and then lets $s\downarrow1$, so that
$\Upsilon$ vanishes in the limit; the
Tsfasman--Vl\u{a}du\c{t} inequality then bounds the change from
$1+\varepsilon$ to $1$ by the term $\varepsilon(\kappa_\infty+\mathcal D)$
of that proposition, which is about $0.00215$ in the proof of
Theorem~\ref{an:ceiling}. This limit is also the reason why $m_0$ is not
effective.

\subsection{\texorpdfstring{The Tsfasman--Vl\u{a}du\c{t} Inequality over $B$}{The Tsfasman-Vladut Inequality over B}}
\label{an:budget-section}
For $q>1$ and $s\ge1$ put $g_s(q)=-\log(1-q^{-s})$.

\begin{lemma}\label{an:local-contribution}
Let $M$ be a finite Galois extension of $B$ and $s>1$. For a prime
$\mathfrak p$ of $B$ let $e_M(\mathfrak p)$ and $f_M(\mathfrak p)$ be its
ramification index and residue degree in $M/B$. Then
\begin{align}
 \frac{\log\zeta_M(s)}{[M:\Q]}
 &=\sum_{\mathfrak p}\frac{g_s(\mathrm N\mathfrak p^{f_M(\mathfrak p)})}
 {2e_M(\mathfrak p)f_M(\mathfrak p)},\label{an:local-sum}\\
 -\frac1{[M:\Q]}\frac{\zeta_M'}{\zeta_M}(2)
 &=\frac1{[M:\Q]}\sum_{\mathfrak P}\frac{\log\mathrm N\mathfrak P}
 {\mathrm N\mathfrak P^2-1}
 =\sum_{\mathfrak p}\frac{\log\mathrm N\mathfrak p}
 {2e_M(\mathfrak p)(\mathrm N\mathfrak p^{2f_M(\mathfrak p)}-1)},
 \label{an:local-debit}
\end{align}
where $\mathfrak P$ runs over the primes of $M$. Each summand on the right
does not increase when $e_M(\mathfrak p)$ or $f_M(\mathfrak p)$ is replaced
by a larger real number. If $M\subseteq M'$ are both Galois over $B$, then
$e_M(\mathfrak p)\mid e_{M'}(\mathfrak p)$ and
$f_M(\mathfrak p)\mid f_{M'}(\mathfrak p)$.
\end{lemma}

\begin{proof}
Above $\mathfrak p$ there are $[M:B]/(ef)$ primes of $M$, each of norm
$\mathrm N\mathfrak p^f$, and $[M:\Q]=2[M:B]$. This gives both formulas.
The first equality in \eqref{an:local-debit} follows from
$-\zeta_M'/\zeta_M(s)=\sum_{\mathfrak P}\log\mathrm N\mathfrak P/(\mathrm
N\mathfrak P^s-1)$. For monotonicity in $f$, write
$g_s(x^f)/f=\sum_{j\ge1}x^{-fjs}/(fj)$ for $x>1$; each term decreases in
$f$. The divisibilities hold because ramification indices and residue
degrees are multiplicative in towers.
\end{proof}

Put
\[
 \kappa_\infty=\frac{\ell-\gamma-\log(4\pi)}4,\qquad
 w(q)=2\log q\sum_{m\ge1}\frac1{q^m+1}.
\]
Here $2\kappa_\infty$ is the right side, and $w(q)$ the weight of the prime
power $q$, in the Tsfasman--Vl\u{a}du\c{t} inequality \eqref{an:tv-budget}
below.

\begin{proposition}\label{an:prime-budget}
Let $\sigma=1+\varepsilon$ with $0<\varepsilon\le1$, and let $Y_*$ be a real
number with $[K:\Q]^{-1}\log\zeta_K(\sigma)\le Y_*$ for every
$K\in\mathcal K$. For each prime $\mathfrak p$ of $B$, let
$e^0_{\mathfrak p},f^0_{\mathfrak p}\ge1$ be integers with
$e_K(\mathfrak p)\ge e^0_{\mathfrak p}$ and $f_K(\mathfrak p)\ge
f^0_{\mathfrak p}$ for every $K\in\mathcal K$, that is, lower bounds for
the type of $\mathfrak p$ in $K/B$, and let
\[
 \mathcal D\ge\sum_{\mathfrak p}\frac{\log\mathrm N\mathfrak p}
 {2e^0_{\mathfrak p}(\mathrm N\mathfrak p^{2f^0_{\mathfrak p}}-1)}.
\]
Then every $C'>Y_*+\varepsilon(\kappa_\infty+\mathcal D)$ has the following
property: there is an $m_0$ such that $d^{-1}\log L_F(1)<C'$ for every
$K\in\mathcal K$ with $[K:B]\ge m_0$.
\end{proposition}

\begin{proof}
\emph{A limiting sequence.} Suppose not. Then there are fields $K^{(j)}\in\mathcal K$ of strictly
increasing degree, hence pairwise non-isomorphic, with
$d_j^{-1}\log L_{F^{(j)}}(1)\ge C'$, where $[K^{(j)}:\Q]=2d_j$. For a prime power $q$
let $N_q(K)$ be the number of primes of $K$ of norm $q$. If $q=p^f$, then
$N_q(K)\le[K:\Q]/f$, since the local degrees above $p$ add up to $[K:\Q]$.
By a diagonal argument we may assume that
$\beta_q=\lim_jN_q(K^{(j)})/(2d_j)$ exists for every $q$.

\emph{The Tsfasman--Vl\u{a}du\c{t} inequality.} Tsfasman and
Vl\u{a}du\c{t} call a sequence $(A_i)$ of pairwise
non-isomorphic number fields asymptotically exact if its genus
$\mathfrak g(A_i)=\log|\Delta_{A_i}|^{1/2}$ tends to infinity and the limits
$\phi_q=\lim N_q(A_i)/\mathfrak g(A_i)$,
$\phi_\R=\lim r_1(A_i)/\mathfrak g(A_i)$ and
$\phi_\C=\lim r_2(A_i)/\mathfrak g(A_i)$ exist. Their unconditional Basic
Inequality$'$ \cite[Section~3.2, Proposition~3.1]{TsfasmanVladut2002} states
that every such sequence satisfies
\[
 2\sum_q\phi_q\log q\sum_{m\ge1}\frac1{q^m+1}
 +\phi_\R\Bigl(\frac\gamma2+\frac12+\log2\sqrt\pi\Bigr)
 +\phi_\C(\gamma+\log4\pi)\le1.
\]
For our sequence $\mathfrak g(K^{(j)})=d_j\ell$ and $K^{(j)}$ is totally imaginary, so
$\phi_q=2\beta_q/\ell$, $\phi_\R=0$ and $\phi_\C=1/\ell$, and the inequality
becomes
\begin{equation}\label{an:tv-budget}
 \sum_q\beta_qw(q)\le\frac{\ell-\gamma-\log(4\pi)}2=2\kappa_\infty.
\end{equation}

\emph{The value at $1$.} Put $Z(s)=\sum_q\beta_qg_s(q)$ for $s\ge1$. Since $g_1(q)\le1/(q-1)$ and
$w(q)\ge2\log q/(q+1)$, we have $g_1(q)\le3w(q)/(2\log2)$, so $Z(1)$ is
finite. For fixed real $s>1$,
$(2d_j)^{-1}\log\zeta_{K^{(j)}}(s)=\sum_q(N_q(K^{(j)})/2d_j)g_s(q)$ tends to $Z(s)$ by
dominated convergence: for $q=p^f$ the terms are at most
$g_s(q)/f\le2q^{-s}$, and $\sum_q q^{-s}<\infty$. By Lemma~\ref{an:shift}
with $s-1$ in place of $\varepsilon$, for $1<s\le2$,
\[
 C'\le\limsup_j\frac{\log L_{F^{(j)}}(1)}{d_j}
 \le Z(s)+\sup_{0\le\theta\le1/2}\Upsilon(s-1,\theta).
\]
As $s\downarrow1$, $Z(s)$ increases to $Z(1)$ and the supremum tends to
zero. Hence $C'\le Z(1)$.

\emph{From $1+\varepsilon$ to $1$.} For every prime power $q$,
\begin{align*}
 g_1(q)-g_{1+\varepsilon}(q)&=\int_1^{1+\varepsilon}\frac{\log q}{q^u-1}\,du
 \le\frac{\varepsilon\log q}{q-1},\\
 \frac{\log q}{q-1}-\frac{w(q)}2&=\log q\sum_{m\ge1}\frac1{q^m(q^m+1)}
 \le\frac{\log q}{q^2-1}.
\end{align*}
With \eqref{an:tv-budget} these give
\[
 Z(1)\le Z(1+\varepsilon)
 +\varepsilon\Bigl(\kappa_\infty+\sum_q\beta_q\frac{\log q}{q^2-1}\Bigr).
\]
By Fatou's lemma, \eqref{an:local-debit} for $M=K^{(j)}$, and the monotonicity
in Lemma~\ref{an:local-contribution},
\[
 \sum_q\beta_q\frac{\log q}{q^2-1}
 \le\liminf_j\sum_{\mathfrak p}\frac{\log\mathrm N\mathfrak p}
 {2e_{K^{(j)}}(\mathfrak p)(\mathrm N\mathfrak p^{2f_{K^{(j)}}(\mathfrak p)}-1)}
 \le\mathcal D.
\]
Finally
$Z(1+\varepsilon)=\lim_j(2d_j)^{-1}\log\zeta_{K^{(j)}}(1+\varepsilon)\le Y_*$.
Thus $C'\le Z(1)\le Y_*+\varepsilon(\kappa_\infty+\mathcal D)$, a
contradiction.
\end{proof}

\subsection{Comparison with the Kummer Field}\label{an:comparison}
We use the set $S$ of the six primes of $B$ above $2$, $3$ and $5$ and the
primes $\mathfrak t_0=7\mathcal O_B$ and $\mathfrak t_1,\mathfrak t_2$ above
$29$ (Lemma~\ref{tw:base}), the group $G_B$ of Definition~\ref{tw:GB} and
its Zassenhaus filtration $D_nG_B$, the relation space $R_2$ in degree two
(Definition~\ref{tw:relation-space}), which Lemma~\ref{tw:quadratic-layer}
computes, the restricted square $v^{[2]}$ of a vector $v\in\Ftwo^8$ as in
Section~\ref{tw:tower}, and the Frobenius vectors
$\overline{\Frob}_{\mathfrak p}$ of Definition~\ref{gf:frobenius-vector}.

\begin{definition}\label{an:P4}
Let $\Sigma_2=\{v\in\Ftwo^8:v^{[2]}\in R_2\}$, and let $\mathcal P_4$ be the
set of primes $\mathfrak p\notin S\cup\{\mathfrak t_0,\mathfrak
t_1,\mathfrak t_2\}$ of $B$ with $\mathrm N\mathfrak p\le10^6$ and
$\overline{\Frob}_{\mathfrak p}\notin\Sigma_2$.
\end{definition}

\begin{lemma}\label{an:census-degree}
Let $K\in\mathcal K$, and let $\mathfrak p\notin S$ be a prime of $B$ with
$\overline{\Frob}_{\mathfrak p}\notin\Sigma_2$. Then $e_K(\mathfrak p)=1$ and
$f_K(\mathfrak p)\ge4$.
\end{lemma}

\begin{proof}
The field $K$ is unramified over $B$ outside $S$, so $e_K(\mathfrak p)=1$,
and a prime $\mathfrak P$ of $K$ above $\mathfrak p$ has a Frobenius
$\Frob_{\mathfrak P}\in\Gal(K/B)$ of order $f_K(\mathfrak p)$. Choose
$g\in G_B$ mapping to $\Frob_{\mathfrak P}$. Its image in
$G_B/D_2G_B=\Gal(E/B)$ is the Frobenius of $\mathfrak p$ in $E$
(Lemma~\ref{gf:in-tower}), whose elementary image is $\overline{\Frob}_{\mathfrak p}$; it is
nonzero, since $0\in\Sigma_2$. By Lemma~\ref{tw:quadratic-layer}, which uses
the complete list of generators of $N_B$ in
Lemma~\ref{tw:GB-presentation}, and since $\overline{\Frob}_{\mathfrak p}^{[2]}\notin R_2$,
$g$ has order at least $4$ in $G_B/D_3G_B$. The projection $G_B\to G_B/D_3G_B$
factors through $\Gal(K/B)$, because $K$ contains the fixed field of
$D_4G_B\subseteq D_3G_B$. So $\Frob_{\mathfrak P}$ has order at least $4$.
\end{proof}

The next proposition compares $\zeta_K(\sigma)$ with $\zeta_E(\sigma)$ and
subtracts two corrections. The first concerns the primes
$\mathfrak p_1,\mathfrak p_2$ above $2$ and $\mathfrak t_0,\mathfrak
t_1,\mathfrak t_2$ above $7$ and $29$, whose types in $K/B$ are known; the
second concerns the primes of $\mathcal P_4$, which are found in
Subsection~\ref{an:census}. Put
\begin{align}
 \Delta_{2,7,29}(\sigma)&=2\Bigl(\frac{g_\sigma(2^2)}{16}
 -\frac{g_\sigma(2^4)}{64}\Bigr)
 +2\Bigl(\frac{g_\sigma(29^2)}4-\frac{g_\sigma(29^4)}8\Bigr)
 +\frac{g_\sigma(7^4)}4-\frac{g_\sigma(7^8)}8,
 \label{an:sel}\\
 \Delta_{\mathcal P_4}(\sigma)&=\sum_{\mathfrak p\in\mathcal P_4}
 \Bigl(\frac{g_\sigma(\mathrm N\mathfrak p^2)}4
 -\frac{g_\sigma(\mathrm N\mathfrak p^4)}8\Bigr).
 \label{an:cen}
\end{align}
Every term of both sums is nonnegative, since $g_\sigma(x^2)\le g_\sigma(x)$.

\begin{proposition}\label{an:Ystar}
Every $K\in\mathcal K$ satisfies
\[
 \frac{\log\zeta_K(\sigma)}{[K:\Q]}
 \le\frac{\log\zeta_E(\sigma)}{512}-\Delta_{2,7,29}(\sigma)
 -\Delta_{\mathcal P_4}(\sigma).
\]
\end{proposition}

\begin{proof}
By Lemma~\ref{gf:in-tower}, $E\subseteq K$, and both are Galois over $B$.
Apply \eqref{an:local-sum} to $M=E$ and $M=K$. By
Lemma~\ref{an:local-contribution}, each term for $K$ is at most the
corresponding term for $E$. For the primes listed below we bound the term
for $K$ from above using what is known about the type in $K/B$, and
subtract the difference between the term for $E$ and this upper bound.

The types in $E/B$ are given by Corollary~\ref{gf:types}: $(4,2)$ at
$\mathfrak p_1,\mathfrak p_2$, and $(1,2)$ at $\mathfrak t_0,\mathfrak
t_1,\mathfrak t_2$ and at the primes of $\mathcal P_4$. By
Theorem~\ref{tw:field-family}, the primes of $K$ above $2$, $29$ and $7$ have
absolute types $(8,4)$, $(1,4)$ and $(1,8)$. The primes $\mathfrak p_j$,
$\mathfrak t_1$ and $\mathfrak t_2$ have residue degree one over $\Q$, and
$\mathfrak t_0$ has residue degree two. So the types in $K/B$ are $(8,4)$ at
$\mathfrak p_j$ and $(1,4)$ at $\mathfrak t_0,\mathfrak t_1,\mathfrak t_2$.
By Lemma~\ref{an:census-degree}, the primes of $\mathcal P_4$ have $e_K=1$
and $f_K\ge4$. The differences of the terms of $E$ and the upper bounds for
the terms of $K$ are exactly the summands of \eqref{an:sel} and
\eqref{an:cen}. For example, at $\mathfrak p_j$ the term of $E$ is
$g_\sigma(2^2)/(2\cdot4\cdot2)$ and that of $K$ is
$g_\sigma(2^4)/(2\cdot8\cdot4)$.
\end{proof}

\subsection{An Approximate Functional Equation for Two Gamma Factors}
\label{an:afe-section}
We evaluate $\zeta_E(\sigma)$ through Proposition~\ref{gf:factorization}. The
factors $L(\sigma,\chi_e)$ are computed from an approximate functional
equation. For real $\mu$ and $x,y>0$ define the incomplete gamma and Bessel
integrals
\[
 H_\mu(x)=\int_1^\infty u^{\mu-1}e^{-xu}\,du,\qquad
 I_\mu(y)=\int_1^\infty u^\mu K_0(yu)\,du,
\]
where $K_0$ is the modified Bessel function of the second kind of order
zero, and their envelopes (upper bounds, by Lemma~\ref{an:kernels}(2))
\[
 \mathcal V_1(x)=\frac{e^{-x}}x\Bigl(1+\frac1x\Bigr),\qquad
 \mathcal V_2(y)=\sqrt{\frac\pi2}\,\frac{e^{-y}}{y^{3/2}}\Bigl(1+\frac1y\Bigr).
\]
Numerical evaluation of $L$-functions from their functional equations is
developed in \cite{Dokchitser2004}. Part (1) of the next lemma treats mixed
gamma factors, and part (2), for pure gamma factors, uses the Bessel
function $K_0$. We prove both.

\begin{lemma}\label{an:balanced-afe}
Let $a_n$ be real numbers with $|a_n|\le\tau(n)$, let $Q>0$ and
$\nu_1,\nu_2\in\{0,1\}$ (integers, not places), and let
$L(s)=\sum_{n\ge1}a_nn^{-s}$. Suppose that
$\Lambda(s)=Q^{s/2}\Gamma_\R(s+\nu_1)\Gamma_\R(s+\nu_2)L(s)$ extends to an
entire function of order at most one with $\Lambda(1-s)=\Lambda(s)$. Put
$t=2\pi/\sqrt Q$ and let $\sigma>1$ be real.
\begin{enumerate}
\item If $\nu_1\ne\nu_2$, then
\[
 L(\sigma)=\frac{t^\sigma}{\Gamma(\sigma)}\sum_{n\ge1}a_n
 \bigl(H_\sigma(tn)+H_{1-\sigma}(tn)\bigr).
\]
\item If $\nu_1=\nu_2=\nu$, then
\[
 L(\sigma)=\frac{4(t/2)^{\sigma+\nu}}{\Gamma((\sigma+\nu)/2)^2}
 \sum_{n\ge1}a_nn^\nu
 \bigl(I_{\sigma+\nu-1}(tn)+I_{\nu-\sigma}(tn)\bigr).
\]
\end{enumerate}
The series converge absolutely.
\end{lemma}

\begin{proof}
In case (1) put $k(x)=e^{-x}$ and $\nu=0$, and in case (2) put
$k(x)=K_0(x)$. Their Mellin transforms are $\Gamma(w)$ and
$2^{w-2}\Gamma(w/2)^2$ for $\operatorname{Re}w>0$
\cite[10.43.19]{DLMF}. Let $\Theta(u)=\sum_na_nn^\nu k(tnu)$ for $u>0$. For
$\operatorname{Re}w>1+\nu$, termwise integration gives
\[
 \widetilde\Theta(w)=\int_0^\infty\Theta(u)u^{w-1}\,du=
 \begin{cases}
 t^{-w}\Gamma(w)L(w)=\tfrac12\Lambda(w)&\text{in case (1)},\\
 2^{w-2}t^{-w}\Gamma(w/2)^2L(w-\nu)=\tfrac14Q^{\nu/2}\Lambda(w-\nu)
 &\text{in case (2)}.
 \end{cases}
\]
In case (1) we used $\Gamma_\R(w)\Gamma_\R(w+1)=\Gamma_\C(w)$ and
$Q^{w/2}(2\pi)^{-w}=t^{-w}$. In both cases $\widetilde\Theta$ is entire and
$\widetilde\Theta(w)=\widetilde\Theta(1+2\nu-w)$.

The function $L=\Lambda/(Q^{s/2}\Gamma_\R(s+\nu_1)\Gamma_\R(s+\nu_2))$ of the
complex variable $s$ is entire of finite order. Fix $\eta>0$. The function
$L$ is bounded on $\operatorname{Re}s=1+\eta$, and, by the functional
equation and Stirling's formula, of polynomial growth on
$\operatorname{Re}s=-\eta$. By the Phragm\'en--Lindel\"of principle it has
polynomial growth in $|\operatorname{Im}s|$ on every vertical strip. With
Stirling's formula this shows that $\widetilde\Theta(w)$ decays exponentially in
$|\operatorname{Im}w|$, uniformly on vertical strips. Mellin inversion on a
line $\operatorname{Re}w=c>1+\nu$ and a shift to
$\operatorname{Re}w=1+2\nu-c$ are therefore justified, and no residues
occur. With the functional equation of $\widetilde\Theta$ they give
\[
 \Theta(u)=u^{-1-2\nu}\,\Theta(1/u)\qquad(u>0).
\]
In particular $\Theta(u)$ decays rapidly as $u\downarrow0$. Splitting the
Mellin integral at $u=1$ and substituting $u\mapsto1/u$ below $1$ gives
\[
 \widetilde\Theta(\sigma+\nu)=\int_1^\infty\Theta(u)
 \bigl(u^{\sigma+\nu-1}+u^{\nu-\sigma}\bigr)\,du.
\]
Termwise integration, justified by the exponential decay of $k$, turns the
right side into $\sum_na_nn^\nu$ times $H_\sigma(tn)+H_{1-\sigma}(tn)$ in case
(1), and times $I_{\sigma+\nu-1}(tn)+I_{\nu-\sigma}(tn)$ in case (2). On the
left, $\widetilde\Theta(\sigma)=t^{-\sigma}\Gamma(\sigma)L(\sigma)$ in case (1), and
$\widetilde\Theta(\sigma+\nu)=2^{\sigma+\nu-2}t^{-\sigma-\nu}
\Gamma((\sigma+\nu)/2)^2L(\sigma)$ in case (2).
\end{proof}

\begin{lemma}\label{an:kernels}
Let $\mu\in\R$.
\begin{enumerate}
\item The functions $H_\mu$ and $I_\mu$ are Laplace transforms of positive
 measures on $[1,\infty)$. They are completely monotone on $(0,\infty)$,
 extend holomorphically to $\operatorname{Re}z>0$, and satisfy
 $|H(z)|\le H(\operatorname{Re}z)$ there, for $H=H_\mu$ or $H=I_\mu$.
\item If $\mu\le2$, then $H_\mu\le\mathcal V_1$. If $\mu\le3/2$, then
 $I_\mu\le\mathcal V_2$.
\item They satisfy
 \[
 xH_\mu'(x)+\mu H_\mu(x)=-e^{-x},\qquad
 yI_\mu'(y)+(\mu+1)I_\mu(y)=-K_0(y),
 \]
 and $K_0$ satisfies $y^2K_0''+yK_0'-y^2K_0=0$ and $K_0'=-K_1$, where
 $K_1$ is the modified Bessel function of the second kind of order one.
\end{enumerate}
\end{lemma}

\begin{proof}
The definition exhibits $H_\mu$ as the Laplace transform of
$u^{\mu-1}\,du$ on $[1,\infty)$. Taking the order $0$ in
\cite[10.32.9]{DLMF} and substituting $v=\cosh r$ in its integral gives
$K_0(x)=\int_1^\infty e^{-xv}(v^2-1)^{-1/2}\,dv$. Substituting this into
$I_\mu$, and then $v=r/u$ in the inner integral, shows
\[
 I_\mu(y)=\int_1^\infty e^{-yr}\Bigl(\int_1^ru^\mu(r^2-u^2)^{-1/2}\,du\Bigr)
 \,dr,
\]
with a nonnegative inner integral. Differentiation under the integral
gives complete monotonicity. Since the measures are positive and
$|e^{-zr}|=e^{-r\operatorname{Re}z}$, we get $|H(z)|\le H(\operatorname{Re}z)$
for $\operatorname{Re}z>0$, which proves (1).

For (2), $u^{\mu-1}\le u$ for $u\ge1$ and $\mu\le2$, and
$\int_1^\infty ue^{-xu}\,du=\mathcal V_1(x)$. For $I_\mu$ we use
$K_0(x)\le K_{1/2}(x)=\sqrt{\pi/(2x)}\,e^{-x}$ \cite[10.39.2]{DLMF}. It
holds because the modified Bessel function of order $\eta\ge0$ is
$\int_0^\infty e^{-x\cosh r}\cosh(\eta r)\,dr$ \cite[10.32.9]{DLMF}, which
increases with $\eta$. Since $u^{\mu-1/2}\le u$ for $\mu\le3/2$,
\[
 I_\mu(y)\le\sqrt{\frac\pi{2y}}\int_1^\infty u^{\mu-1/2}e^{-yu}\,du
 \le\mathcal V_2(y).
\]

For (3), differentiate under the integral and integrate by parts:
$xH_\mu'(x)=-\int_1^\infty u^\mu xe^{-xu}\,du=-e^{-x}-\mu H_\mu(x)$, and
$yI_\mu'(y)=\int_1^\infty u^{\mu+1}\frac{d}{du}K_0(yu)\,du
=-K_0(y)-(\mu+1)I_\mu(y)$. The last two identities are Bessel's equation
and a standard derivative \cite[10.25.1, 10.29.3]{DLMF}.
\end{proof}

In the application $\sigma=301/300$, so $1<\sigma\le3/2$. The four exponents $\sigma$,
$1-\sigma$, $\sigma+\nu-1$ and $\nu-\sigma$ in Lemma~\ref{an:balanced-afe}
are then covered by the envelopes: $\mathcal V_1$ bounds $H_\sigma$ and
$H_{1-\sigma}$, and $\mathcal V_2$ bounds $I_{\sigma+\nu-1}$ and
$I_{\nu-\sigma}$.

\begin{lemma}\label{an:tails}
Let $1<\sigma\le3/2$, let $|a_n|\le\tau(n)$, let $t>0$, $N\ge1$ and
$\omega>1$.
\begin{enumerate}
\item If $t(N+1)\ge\omega-1$, then for $\mu\in\{\sigma,1-\sigma\}$
\[
 \sum_{n>N}|a_n|H_\mu(tn)\le\zeta(\omega)^2(N+1)^\omega
 \mathcal V_1\bigl(t(N+1)\bigr).
\]
\item If $\nu\in\{0,1\}$ and $t(N+1)\ge\omega+\nu-\frac32$, then for
 $\mu\in\{\sigma+\nu-1,\nu-\sigma\}$
\[
 \sum_{n>N}|a_n|n^\nu I_\mu(tn)\le\zeta(\omega)^2(N+1)^{\omega+\nu}
 \mathcal V_2\bigl(t(N+1)\bigr).
\]
\end{enumerate}
\end{lemma}

\begin{proof}
The logarithmic derivative of
$x^\omega\mathcal V_1(x)=x^{\omega-1}e^{-x}(1+1/x)$ is
$(\omega-1)/x-1-1/(x(x+1))<0$ for $x\ge\omega-1$. So
$n^\omega\mathcal V_1(tn)\le(N+1)^\omega\mathcal V_1(t(N+1))$ for $n>N$. By
Lemma~\ref{an:kernels}(2),
\[
 \sum_{n>N}|a_n|H_\mu(tn)\le\sum_{n>N}\frac{\tau(n)}{n^\omega}\,
 n^\omega\mathcal V_1(tn)\le(N+1)^\omega\mathcal V_1\bigl(t(N+1)\bigr)
 \sum_{n\ge1}\frac{\tau(n)}{n^\omega},
\]
and the last sum is $\zeta(\omega)^2$. Part (2) is the same, since
$y^{\omega+\nu}\mathcal V_2(y)$ is a multiple of
$y^{\omega+\nu-3/2}e^{-y}(1+1/y)$, which decreases for
$y\ge\omega+\nu-\frac32$.
\end{proof}

\subsection{\texorpdfstring{Enclosures of $H_\mu$ and $I_\mu$, and Finite Sums}{Enclosures of H and I, and Finite Sums}}
\label{an:kernel-enclosures}
An \emph{enclosure} of a real number is an interval that contains it. We
enclose $H_\mu$ and $I_\mu$ on a geometric grid. Let $H$ be $H_\mu$ or $I_\mu$
with $\mu$ as in Lemma~\ref{an:tails}, and let $\mathcal V$ be its envelope,
$\mathcal V_1$ or $\mathcal V_2$. For $J\ge0$ and $\xi>0$ put
\[
 \rho_J(\xi)=\mathcal V(\xi/2)\,\frac{16^{-J-1}}{1-1/16}.
\]

\begin{lemma}\label{an:taylor}
Let $\xi>0$ and let $h_j=H^{(j)}(\xi)/j!$ be the Taylor coefficients of $H$
at $\xi$.
\begin{enumerate}
\item $(-1)^jh_j\ge0$ and $|h_j|\le\mathcal V(\xi/2)(2/\xi)^j$ for all
 $j\ge0$.
\item For $|x-\xi|\le\xi/32$,
 $\bigl|H(x)-\sum_{j=0}^Jh_j(x-\xi)^j\bigr|\le\rho_J(\xi)$. For
 $x=31\xi/32$ the difference lies in $[0,\rho_J(\xi)]$.
\item Let $k_j$ be the Taylor coefficients at $\xi$ of $e^{-x}$ if
 $H=H_\mu$, and of $K_0$ if $H=I_\mu$. If $H=H_\mu$, then
 $(j+1)\xi h_{j+1}=-k_j-(j+\mu)h_j$, and $k_j=(-1)^je^{-\xi}/j!$. If
 $H=I_\mu$, then $(j+1)\xi h_{j+1}=-k_j-(j+\mu+1)h_j$, where
 $k_0=K_0(\xi)$, $k_1=-K_1(\xi)$, and
 \[
 (j+1)(j+2)\xi^2k_{j+2}=-(j+1)(2j+1)\xi k_{j+1}-(j^2-\xi^2)k_j
 +2\xi k_{j-1}+k_{j-2},
 \]
 with $k_{-1}=k_{-2}=0$.
\end{enumerate}
\end{lemma}

\begin{proof}
Complete monotonicity gives the signs in (1). The disc $|z-\xi|\le\xi/2$ lies
in $\operatorname{Re}z\ge\xi/2$, where $|H(z)|\le H(\xi/2)\le\mathcal
V(\xi/2)$ by Lemma~\ref{an:kernels}. Cauchy's estimate gives the bound on
$h_j$. For $|x-\xi|\le\xi/32$, the omitted terms have modulus at most
$\mathcal V(\xi/2)16^{-j}$, and their sum over $j>J$ is at most
$\rho_J(\xi)$. For $x<\xi$ every term $h_j(x-\xi)^j$ is nonnegative by (1),
so the remainder is nonnegative. Part (3) follows by comparing coefficients
of $(x-\xi)^j$ in the differential equations of Lemma~\ref{an:kernels}(3),
written at $x=\xi+(x-\xi)$.
\end{proof}

The grid points are $\xi_i=64\,(31/32)^i$ for $0\le i\le i_1$, where
$i_1=1878$ is the least index with $\xi_{i_1}\le2^{-80}$. Taylor polynomials
are formed at $\xi_0,\ldots,\xi_{i_1-1}$. The initial values
$e^{-\xi_i}$, $K_0(\xi_i)$ and $K_1(\xi_i)$ of part~(3) of
Lemma~\ref{an:taylor} are enclosed with the rigorous exponential and Bessel
functions of Arb \cite{Johansson2017}, which is now part of FLINT
\cite{FLINT} and is called through its Python interface.

\begin{lemma}\label{an:grid}
Let $J\ge0$, and suppose that enclosures of $e^{-\xi_i}$, $K_0(\xi_i)$ and
$K_1(\xi_i)$ are given for $0\le i<i_1$. Start from the enclosure
$[0,\mathcal V(64)]$ of $H(\xi_0)$. For $i=0,1,\ldots,i_1-1$, compute from
the enclosure of $H(\xi_i)$ enclosures of $h_0,\ldots,h_J$ at $\xi_i$ by the
recurrences of Lemma~\ref{an:taylor}(3), and take
\[
 \sum_{j=0}^Jh_j(-\xi_i/32)^j+[0,\rho_J(\xi_i)],
\]
evaluated in interval arithmetic, as the enclosure of $H(\xi_{i+1})$. Then
every interval so obtained contains the corresponding true value.
\end{lemma}

\begin{proof}
By Lemma~\ref{an:kernels}, $H(\xi_0)\in[0,\mathcal V(64)]$. Suppose that
the enclosure of $H(\xi_i)$ contains $H(\xi_i)$. The true Taylor
coefficients at $\xi_i$ are obtained from the true value $H(\xi_i)$ by the
same recurrence, so their enclosures contain them. Since
$\xi_{i+1}=31\xi_i/32$, Lemma~\ref{an:taylor}(2) shows that the enclosure
of $H(\xi_{i+1})$ contains $H(\xi_{i+1})$. Induction on $i$ proves the
lemma.
\end{proof}

Now let $t>0$, $\nu\in\{0,1\}$, integers $a_1,\ldots,a_N$ and $J\ge0$ be
given, with $tN<\xi_0$ and $t>\xi_{i_1}$. For $0\le i<i_1$ let
$\mathcal B_i=\{n:1\le n\le N,\ \xi_{i+1}<tn\le\xi_i\}$. These sets
partition $\{1,\ldots,N\}$. For nonempty $\mathcal B_i$, let $\bar n_i$ be
the integer part of the average of its least and largest elements, and put
\[
 M_{i,j}=\sum_{n\in\mathcal B_i}a_nn^\nu(n-\bar n_i)^j,\qquad
 A_i=\sum_{n\in\mathcal B_i}|a_n|n^\nu,\qquad
 P_{i,j}=t^j\sum_{m=j}^J\binom mjh_{i,m}(t\bar n_i-\xi_i)^{m-j},
\]
where $h_{i,m}$ are the Taylor coefficients of $H$ at $\xi_i$. The moments
$M_{i,j}$ and the absolute sums $A_i$ are exact integers.

\begin{lemma}\label{an:binning}
With this notation,
\begin{equation}\label{an:bin-sum}
 \Bigl|\sum_{n=1}^Na_nn^\nu H(tn)-\sum_i\sum_{j=0}^JP_{i,j}M_{i,j}\Bigr|
 \le\sum_i\rho_J(\xi_i)A_i,
\end{equation}
where $i$ runs over the indices with $\mathcal B_i$ nonempty.
\end{lemma}

\begin{proof}
Let $n\in\mathcal B_i$. Then $|tn-\xi_i|<\xi_i/32$, so by
Lemma~\ref{an:taylor}(2) the value $H(tn)$ differs from
$\sum_{m=0}^Jh_{i,m}(tn-\xi_i)^m$ by at most $\rho_J(\xi_i)$. Since
$tn-\xi_i=t(n-\bar n_i)+(t\bar n_i-\xi_i)$, the binomial theorem gives
$\sum_{m=0}^Jh_{i,m}(tn-\xi_i)^m=\sum_{j=0}^JP_{i,j}(n-\bar n_i)^j$.
Multiply by $a_nn^\nu$ and sum over $n\in\mathcal B_i$ and over $i$.
\end{proof}

In the application below $t\ge2\pi/\sqrt{3470400}>0.0033$, because
$Q_e\le3470400$ for every $e$ (Corollary~\ref{gf:conductors}). So every argument
$tn$ exceeds $0.0033$, and the part of the grid below that point is not
used.

\subsection{\texorpdfstring{The Value of $\zeta_E$ at $\sigma=301/300$}{The Value of zeta(E) at sigma = 301/300}}
\label{an:value-section}
From now on $\sigma=301/300$. Let $e\in\Ftwo^8\setminus\{0\}$. By
Proposition~\ref{gf:degree-two}, $L(s,\chi_e)$ satisfies the hypotheses of
Lemma~\ref{an:balanced-afe} with $Q=Q_e$ and $\nu_k=\nu_k(e)$. Put
$t=2\pi/\sqrt{Q_e}$ and
\[
 N_e=4\lfloor\sqrt{Q_e}\rfloor+1,
\]
and let $\Pi_e$ be the prefactor of that lemma:
$\Pi_e=t^\sigma/\Gamma(\sigma)$ if $\nu_1(e)\ne\nu_2(e)$, and
$\Pi_e=4(t/2)^{\sigma+\nu}/\Gamma((\sigma+\nu)/2)^2$ if
$\nu_1(e)=\nu_2(e)=\nu$; in the first case put $\nu=0$. Let
$\mathcal S_1$ and $\mathcal S_2$ be the sums
$\sum_i\sum_{j=0}^{20}P_{i,j}M_{i,j}$ of Lemma~\ref{an:binning}, with
$N=N_e$, $a_n=a_n(e)$ and $J=20$, for the two functions of the series of
Lemma~\ref{an:balanced-afe}: $H=H_\sigma$ and $H=H_{1-\sigma}$ in the first
case, and $H=I_{\sigma+\nu-1}$ and $H=I_{\nu-\sigma}$ in the second; their
envelope is $\mathcal V=\mathcal V_1$ in the first case and
$\mathcal V=\mathcal V_2$ in the second. For every $e$, $N_e$ lies between
$105$ and $7449$, and $t(N_e+1)$ between $24.7$ and $25.5$ (the
supplementary program \texttt{check255.gp} prints both ranges). Hence
$tN_e<\xi_0$ and $t>\xi_{i_1}$, as Lemma~\ref{an:binning} requires, and the
hypotheses of Lemma~\ref{an:tails} hold for every $\omega$ in the list
$\frac{101}{100},\frac{21}{20},\frac{11}{10},\frac98,\frac65,\frac54,
\frac43,\frac32,2$. The value of $\omega$ in this list is chosen to minimize
the bound
\[
 \mathcal T_e=\sum_i\rho_{20}(\xi_i)A_i
 +\zeta(\omega)^2(N_e+1)^{\omega+\nu}\mathcal V\bigl(t(N_e+1)\bigr).
\]

\begin{lemma}\label{an:row}
For $\sigma=301/300$ and every $e\in\Ftwo^8\setminus\{0\}$,
\begin{equation}\label{an:row-enclosure}
 \bigl|L(\sigma,\chi_e)-\Pi_e(\mathcal S_1+\mathcal S_2)\bigr|
 \le2\Pi_e\mathcal T_e.
\end{equation}
\end{lemma}

\begin{proof}
Split each of the two series of Lemma~\ref{an:balanced-afe} at $N_e$. By
Lemma~\ref{an:binning}, each finite part differs from its approximation
$\mathcal S_1$ or $\mathcal S_2$ by at most $\sum_i\rho_{20}(\xi_i)A_i$; the
same bound serves both series, since their two functions have the same
envelope. By Lemma~\ref{an:tails}, each tail is at most
$\zeta(\omega)^2(N_e+1)^{\omega+\nu}\mathcal V(t(N_e+1))$. Multiplying by
$\Pi_e$ gives \eqref{an:row-enclosure}.
\end{proof}

\begin{proposition}\label{an:genus-value}
For $\sigma=301/300$,
\[
 0.08264460177456<\frac{\log\zeta_E(\sigma)}{512}<0.08264460807138.
\]
\end{proposition}

\begin{proof}[Proof (computer-assisted)]
By Proposition~\ref{gf:factorization},
\[
 \log\zeta_E(\sigma)=\log\zeta(\sigma)+\log L(\sigma,\chi_B)
 +\sum_{e\ne0}\log L(\sigma,\chi_e).
\]
The coefficients $a_n(e)$ for $n\le N_e$ are computed exactly from the Euler
factors of Proposition~\ref{gf:local-data} by the supplementary program
\texttt{lfun241.py}. The supplementary program \texttt{afe241.py} evaluates
\eqref{an:row-enclosure} in Arb midpoint--radius interval arithmetic
\cite{Johansson2017,FLINT} with $256$ bits of precision, rounding every
quantity outward. Its routines for
Lemmas~\ref{an:balanced-afe}--\ref{an:binning} are taken unchanged from an
earlier supplementary archive of the author (Subsection~\ref{intro:code}), and
they implement these lemmas as stated here. Every enclosure of
$L(\sigma,\chi_e)$ is positive. The factor
$\zeta_B(\sigma)=\zeta(\sigma)L(\sigma,\chi_B)$ is evaluated with
\[
 L(\sigma,\chi_B)=241^{-\sigma}\sum_{a=1}^{240}\Bigl(\frac
 a{241}\Bigr)\zeta\Bigl(\sigma,\frac a{241}\Bigr),
\]
where $\zeta(s,x)=\sum_{n\ge0}(n+x)^{-s}$ is the Hurwitz zeta function,
using Arb's rigorous implementation of it \cite{Johansson2015Hurwitz}. This
gives $\zeta_B(\sigma)=725.51409164864\ldots$ and
$L(\sigma,\chi_B)=2.41373420241\ldots$. Summing the logarithms gives the
enclosure
\[
 [0.08264460177456168\ldots,\ 0.08264460807137072\ldots]
\]
of $\log\zeta_E(\sigma)/512$, which the proposition states rounded outward.
The logarithms are summed in Arb before any output conversion. The JSON
endpoints for the individual factors are rounded outward to double
precision, and their serialized decimal values are checked against the
Arb enclosures. The bounds for the normalized logarithm are retained as
Arb midpoint--radius strings, including their radii.
So the inequalities follow from
Lemmas~\ref{an:balanced-afe}--\ref{an:row}, the exact coefficients, and
outward-rounded interval arithmetic.
\end{proof}

As a check
independent of the approximate functional equation, the supplementary
programs \texttt{yecheck.gp} and \texttt{lvalues255.gp} evaluate the same
values numerically with PARI/GP \cite{PARI}: each of the $255$ values
$L(\sigma,\chi_e)$ lies in its enclosure, and PARI's value
$0.08264460492293\ldots$ of $\log\zeta_E(\sigma)/512$ lies in the enclosure
above.

\subsection{Frobenius Vectors and Lower Bounds for the Types}
\label{an:census}

\begin{lemma}\label{an:sigma2}
The set $\Sigma_2$ of Definition~\ref{an:P4} has exactly $21$ elements: zero
and the following twenty vectors, written as strings $c_0c_1\cdots c_7$ of
their coordinates, as in Section~\ref{tw:tower}:
\[
\begin{array}{ccccc}
 00100000&11110010&11010010&00010000&10000001\\
 10010001&00001000&10100111&10101111&00000100\\
 11101011&11101111&00000010&01100101&01100111\\
 00000001&01011010&01011011&10111010&11000101.
\end{array}
\]
In the notation of Table~\ref{tw:vector-table}, these twenty vectors are the
elementary images of the local involutions $x_j$, $y_j$ and $x_jy_j$ for
$j=1,2$, of the nonidentity elements of the four local groups
$C_2\times C_2$ above $3$ and $5$, and of $\iota_1$ and $\iota_2$.
\end{lemma}

\begin{proof}
Each of these images lies in $\Sigma_2$ by Lemma~\ref{tw:quadratic-layer},
since an involution has trivial square. The supplementary program
\texttt{census241.py} computes $\Sigma_2$ by linear algebra over $\Ftwo$
from the relations of Lemma~\ref{tw:quadratic-layer}, and finds no other
vectors.
\end{proof}

\begin{proposition}[Frobenius vectors of the primes of norm at most $10^6$]\label{an:census-prop}
Among the $78616$ primes
$\mathfrak p\notin S\cup\{\mathfrak t_0,\mathfrak t_1,\mathfrak t_2\}$ of
$B$ with $\mathrm N\mathfrak p\le10^6$, exactly $246$ have
$\overline{\Frob}_{\mathfrak p}=0$, exactly $5994$ have
$\overline{\Frob}_{\mathfrak p}\in\Sigma_2\setminus\{0\}$, and the remaining $72376$ form
$\mathcal P_4$.
\end{proposition}

\begin{proof}
The supplementary program \texttt{census241.py} computes the Frobenius
vector of each of these primes from Proposition~\ref{gf:local-data}(1) with
exact modular arithmetic, and compares it with the elements of $\Sigma_2$
(Lemma~\ref{an:sigma2}).
\end{proof}

Below we use this classification of the primes, not the three counts.

\begin{lemma}\label{an:corrections}
For $\sigma=301/300$,
\begin{equation}\label{an:census-values}
 \Delta_{2,7,29}(\sigma)>0.0344534299,\qquad
 \Delta_{\mathcal P_4}(\sigma)>0.0016298411.
\end{equation}
\end{lemma}

\begin{proof}
The supplementary program \texttt{ceiling241.py} uses the classification of
Proposition~\ref{an:census-prop} and evaluates \eqref{an:sel} and
\eqref{an:cen} in Arb interval arithmetic. It gives
$\Delta_{2,7,29}(\sigma)=0.03445342991613\ldots$ and
$\Delta_{\mathcal P_4}(\sigma)=0.00162984111205\ldots$, and the lemma states
these values rounded down.
\end{proof}

For every prime $\mathfrak p$ of $B$ put
\begin{equation}\label{an:floors}
 (e^0_{\mathfrak p},f^0_{\mathfrak p})=
 \begin{cases}
 (8,4)&\mathfrak p=\mathfrak p_1,\mathfrak p_2,\\
 (2,2)&\mathfrak p=\mathfrak q_1,\mathfrak q_2,\mathfrak r_1,\mathfrak r_2,\\
 (1,4)&\mathfrak p\in\{\mathfrak t_0,\mathfrak t_1,\mathfrak t_2\}\cup\mathcal P_4,\\
 (1,2)&\mathrm N\mathfrak p\le10^6,\ \overline{\Frob}_{\mathfrak p}\in\Sigma_2\setminus\{0\},\\
 (1,1)&\text{otherwise},
 \end{cases}
\end{equation}
where $\mathfrak q_1,\mathfrak q_2$ and $\mathfrak r_1,\mathfrak r_2$ are the
primes above $3$ and $5$ (Lemma~\ref{tw:base}).

\begin{lemma}\label{an:floors-valid}
For every $K\in\mathcal K$ and every prime $\mathfrak p$ of $B$,
$e_K(\mathfrak p)\ge e^0_{\mathfrak p}$ and
$f_K(\mathfrak p)\ge f^0_{\mathfrak p}$. Thus \eqref{an:floors} gives lower
bounds for the types in $K/B$, as Proposition~\ref{an:prime-budget}
requires.
\end{lemma}

\begin{proof}
At $\mathfrak p_j$ and $\mathfrak t_k$ the types in $K/B$ are $(8,4)$ and
$(1,4)$ by Theorem~\ref{tw:field-family}, as in the proof of
Proposition~\ref{an:Ystar}. At $\mathfrak q_j$ and $\mathfrak r_j$ the type
in $E/B$ is $(2,2)$ by Corollary~\ref{gf:types}, and if
$\overline{\Frob}_{\mathfrak p}\ne0$, then $f_E(\mathfrak p)=2$ by the same corollary.
Since $E\subseteq K$, the ramification indices and residue degrees in $E/B$
divide those in $K/B$ (Lemma~\ref{an:local-contribution}). On
$\mathcal P_4$ use Lemma~\ref{an:census-degree}. The bounds $(1,1)$ hold
trivially.
\end{proof}

\begin{lemma}\label{an:D}
For the lower bounds \eqref{an:floors},
\begin{equation}\label{an:R}
 \sum_{\mathfrak p}\frac{\log\mathrm N\mathfrak p}
 {2e^0_{\mathfrak p}(\mathrm N\mathfrak p^{2f^0_{\mathfrak p}}-1)}
 <0.0085105513.
\end{equation}
Hence $\mathcal D=0.0085105513$ satisfies the hypothesis of
Proposition~\ref{an:prime-budget}.
\end{lemma}

\begin{proof}
The primes of norm greater than $X=10^6$ receive $(1,1)$. At most two primes
of $B$ have a given norm, and $x\mapsto\log x/x^2$ decreases on
$[2,\infty)$. Their contribution to the sum is therefore at most
\[
 \sum_{m>X}\frac{\log m}{m^2-1}
 \le\frac{1}{1-X^{-2}}\sum_{m>X}\frac{\log m}{m^2}
 \le\frac{1}{1-X^{-2}}\int_X^\infty\frac{\log x}{x^2}\,dx
 =\frac{\log X+1}{X(1-X^{-2})}.
\]
The supplementary program \texttt{ceiling241.py} sums the primes of norm at
most $X$ in Arb interval arithmetic and evaluates this tail bound entirely
in Arb, including the logarithm. The resulting upper endpoint is less
than $0.0085105513$.
\end{proof}

\subsection{\texorpdfstring{Proof of Theorem~\ref{an:ceiling}}{Proof of the Upper Bound}}

\begin{corollary}\label{an:Ystar-value}
Let $\sigma=301/300$. Every $K\in\mathcal K$ satisfies
$[K:\Q]^{-1}\log\zeta_K(\sigma)\le Y_*$, where
\[
 Y_*=0.0826446081-0.0344534299-0.0016298411=0.0465613371.
\]
\end{corollary}

\begin{proof}
Combine Proposition~\ref{an:Ystar}, the upper bound of
Proposition~\ref{an:genus-value} rounded up, and the lower bounds of
Lemma~\ref{an:corrections}, which are rounded down.
\end{proof}

\begin{proof}[Proof of Theorem~\ref{an:ceiling}]
Let $\sigma=301/300$ and $\varepsilon=1/300$. We apply
Proposition~\ref{an:prime-budget} with $Y_*=0.0465613371$
(Corollary~\ref{an:Ystar-value}), with the lower bounds \eqref{an:floors}
for the types (Lemma~\ref{an:floors-valid}), and with
$\mathcal D=0.0085105513$ (Lemma~\ref{an:D}). Numerically
$\kappa_\infty=(\ell-\gamma-\log(4\pi))/4=0.63694120292\ldots<0.6369412030$,
so
\[
 Y_*+\varepsilon(\kappa_\infty+\mathcal D)
 <0.0465613371+\frac{0.6369412030+0.0085105513}{300}
 <0.0487128430<C.
\]
The interval evaluation of the same expression from the unrounded
enclosures, by \texttt{ceiling241.py}, gives the upper bound
$0.0487128429$. The replay compares the complete serialized Arb enclosure,
including its printed radius, with $C$. Proposition~\ref{an:prime-budget} with $C'=C$ proves
the theorem.
\end{proof}

\begin{remark}\label{an:limit}
The primes above $2$, $3$, $5$, $7$ and $29$ alone limit what
Proposition~\ref{an:prime-budget} can give. Let $(K^{(j)})$, $\beta_q$ and
$Z(s)=\sum_q\beta_qg_s(q)$ be as in the proof of that proposition. The
absolute types of the primes of $K^{(j)}$ above $2,3,5,7,29$ are known exactly
(Theorem~\ref{tw:field-family}), so these primes contribute exactly
\[
 \frac{2g_1(2^4)}{64}+\frac{2g_1(3^2)}8+\frac{2g_1(5^2)}8
 +\frac{2g_1(29^4)}8+\frac{g_1(7^8)}8=0.0416684615\ldots
\]
to $Z(1)$, a value printed by \texttt{ceiling241.py}; all other
contributions are nonnegative. The proof of
Proposition~\ref{an:prime-budget} shows that
$Z(1)\le Y_*+\varepsilon(\kappa_\infty+\mathcal D)$ for every choice of
$\sigma$, $Y_*$, $\mathcal D$ and lower bounds for the types that satisfies
its hypotheses. So no such choice gives an upper bound below this number.
The bound $C=\ceiling$ exceeds it by $0.0070444$.
\end{remark}

\begin{remark}[The hypothesis of the Lean formalization]\label{an:lean-hypothesis}
The Lean formalization \cite{NaslundLean2026} proves, under one explicit
hypothesis that it does not prove, that there are finite planar sets $U_j$
with $|U_j|\to\infty$ and $u(U_j)/|U_j|^{1.0427}\to\infty$. The hypothesis is
\begin{equation}\label{an:h241}
 \frac{\log\zeta_E(1+\frac1{300})}{512}
 +\frac1{300}\Bigl(\frac{\ell-\gamma-\log(4\pi)}4
 -\frac1{512}\frac{\zeta_E'}{\zeta_E}(2)\Bigr)<0.0852.
\end{equation}
In the formalization, $E$ is the subfield of an algebraic closure of $\Q$
generated by $\sqrt{241}$ and square roots of $\alpha_0,\ldots,\alpha_7$; it
is isomorphic to the Kummer field $E$, and the formalization calls it the
genus field.

The left side of \eqref{an:h241} is the bound of this section with all
local data taken from the fixed subfield $E$. By \eqref{an:local-debit} for
$M=E$, $-(\zeta_E'/\zeta_E)(2)/512$ is the sum defining $\mathcal D$ in
Proposition~\ref{an:prime-budget} when the lower bounds are the types in
$E/B$. These lower bounds are valid since $E\subseteq K$. The choice
$Y_*=\log\zeta_E(\sigma)/512$ is valid by Lemma~\ref{an:local-contribution}.
So the left side of \eqref{an:h241} is the bound given by
Proposition~\ref{an:prime-budget} without the refinements of
Proposition~\ref{an:Ystar} and with the types in $E/B$ in place of those in
$K/B$. The supplementary program \texttt{h241\_receipt.py} bounds it by
$0.0848335193<0.0852$. It uses the upper endpoint in
Proposition~\ref{an:genus-value}, the value of $\kappa_\infty$, and
$-(\zeta_E'/\zeta_E)(2)/512<0.0197321539$, which it obtains by summing
\eqref{an:local-debit} in interval arithmetic over the primes of norm at
most $10^6$ and bounding the rest as in the proof of Lemma~\ref{an:D}.

From \eqref{an:h241}, the formalization subtracts corrections at the primes
above $2$, $7$ and $29$ and at the primes above $41$, $47$, $53$, $59$,
$61$, $67$, $79$, $83$ and $97$, which lie in $\mathcal P_4$. It obtains the
weaker bound $0.0495$ for $d^{-1}\log L_F(1)$ for the fields of the tower
in the formalization. This suffices for the exponent $1.0427$ but not for
$1.04273$. The bound $C=\ceiling$ of Theorem~\ref{an:ceiling} uses
Proposition~\ref{an:census-prop} and the lower bounds for the types in
$K/B$; this refinement is not formalized.
\end{remark}

\section{Units of Relative Norm One and Planar Point Sets}
\label{geo:section}

This section and Sections~\ref{fw:section} and~\ref{pf:section} turn
arithmetic data of a quadratic extension into planar point sets. Let $K/F$ be a quadratic extension of
number fields, let $\iota$ be the nontrivial automorphism of $K$ over
$F$, and let $(b,c)$ be the signature of $F$. Throughout these three
sections we assume:
\begin{enumerate}
\item[(G1)] $K$ is totally imaginary, and $b\ge1$;
\item[(G2)] $K/F$ is unramified at every finite place;
\item[(G3)] $\rd(K)=\rd(F)=\lambda$.
\end{enumerate}
Here the root discriminant $\rd$ and the discriminants $\Delta_M$ are as in
Section~\ref{tw:tower}. By (G2), $|\Delta_K|=|\Delta_F|^2$, so (G3) amounts
to $\rd(F)=\lambda$.
Put $d=[F:\Q]=b+2c$, $\theta=c/d$ and $\ell=\log\lambda$. Then
$[K:\Q]=2d$, $|\Delta_K|=\lambda^{2d}$, $|\Delta_F|=\lambda^d$ and
$0\le\theta<1/2$, where $\Delta_K$ and $\Delta_F$ are the absolute
discriminants. Theorem~\ref{tw:field-family} provides such extensions,
with $\lambda=2^{9/4}\sqrt{3615}$ and $\theta\ge\theta_*$; there $\iota$
is the complex conjugation $\iota_1$. The arguments of these three
sections use no Galois structure over $B$ or over $\Q$. For a nonzero
fractional ideal $\mathfrak a$ of $K$ or of $F$, $N\mathfrak a$ denotes
its absolute norm, so that $N(x\mathcal O_K)=|N_{K/\Q}x|$
\cite[Propositions~4.1 and~4.2]{MilneANT}. We take the unit theorem and
the covolumes of ideal lattices from \cite{MilneANT}, and the Herbrand
quotient, Hilbert's Theorem~90 and the analytic class-number formula
from \cite{MilneCFT}.

The argument has five steps. Proposition~\ref{geo:mass} is the
class-number formula for the units of relative norm one: it expresses
their regulator, together with the relative class number and the
capitulation kernel, through $L_F(1)$. Lemma~\ref{geo:selection} uses the
split primes and one ideal class to produce many elements
$\beta_1,\dots,\beta_t$ of relative norm one in a single fractional ideal,
with pairwise disjoint cosets $\beta_i\mathcal O_K^1$.
Lemmas~\ref{geo:unfold-lemma} and~\ref{geo:joint-average-lemma} average
over the units of relative norm one and over translates of the ideal
lattice; this counts exactly, on average, the pairs of lattice points in a
window (Definition~\ref{geo:window}) that differ by the image of an element
of $\bigcup_i\beta_i\mathcal O_K^1$. Lemmas~\ref{geo:dual}
and~\ref{geo:poisson} bound the number of lattice points in a window
uniformly, by Poisson summation. Finally, Proposition~\ref{geo:transfer},
the geometric transfer, takes the window to be a sublevel set of the total
energy of a pair of profiles (Definition~\ref{geo:profiles} and
Lemma~\ref{geo:concentration}) and combines these steps.

\subsection{Coordinates}
A real place of $F$ extends to $K$ either as two real places or as one
complex place, and by (G1) the second case occurs. A complex place of
$F$ has two extensions to $K$, since $\C$ has no quadratic extension.
Thus $K$ has exactly $d$ complex places. For each real place $v$ of
$F$, let $\phi_v\colon K\to\C$ be one of the two embeddings inducing
the place of $K$ above $v$. Number the complex places of $F$ from $1$
to $c$. For the $j$-th one fix an embedding $\phi_j^+\colon K\to\C$
whose restriction $\tau_j$ to $F$ induces it, and put
$\phi_j^-=\phi_j^+\circ\iota$. The embeddings $\phi_j^\pm$ induce
the two places of $K$ above this place: they are distinct, and they are
not complex conjugate, because $\tau_j$ is not real. We identify
$K\otimes_\Q\R$ with $\C^d$ through the map $x\mapsto x_\infty$, where
\[
 x_\infty=\bigl((\phi_v(x))_v,\ (\phi_j^+(x),\phi_j^-(x))_{j}\bigr),
\]
write $\mathfrak a_\infty=\{x_\infty:x\in\mathfrak a\}$ for
$\mathfrak a\subset K$, and identify $\C^d$ with $\R^{2d}$ through real
and imaginary parts. Volumes are Lebesgue measure in these coordinates,
and $\langle x,\xi\rangle=\operatorname{Re}\sum_w x_w\overline{\xi_w}$
is the real inner product, where $w$ runs over the $d$ coordinates. For a
Borel set $\Omega$, $|\Omega|$ denotes its volume; for a finite set, $|\cdot|$
denotes its cardinality.

At a real place $v$ of $F$, the automorphism $\iota$ induces complex
conjugation on the completion $\C$ of $K$, so
$\phi_v(\iota x)=\overline{\phi_v(x)}$. Writing $v(y)\in\R$ for the
image of $y\in F$ under the real embedding $v$, we obtain for $x\in K$
\begin{equation}\label{geo:norm-modulus}
 |\phi_v(x)|^2=v(N_{K/F}x),\qquad
 \phi_j^+(x)\,\phi_j^-(x)=\tau_j(N_{K/F}x).
\end{equation}
For $x\in K^\times$ put
\[
 l_j(x)=\tfrac12\bigl(\log|\phi_j^+(x)|-\log|\phi_j^-(x)|\bigr),
 \qquad l(x)=(l_1(x),\dots,l_c(x))\in\R^c .
\]
If $N_{K/F}\beta=1$, then by \eqref{geo:norm-modulus}
$|\phi_v(\beta)|=1$ for every real place $v$, and
$|\phi_j^\pm(\beta)|=e^{\pm l_j(\beta)}$ for every $j$. Accordingly the
$b$ coordinates $\phi_v$ are called the \emph{compact coordinates}:
there the elements of relative norm one lie on the unit circle. The $c$
pairs $(\phi_j^+,\phi_j^-)$ are the \emph{pair coordinates}: there
they have reciprocal moduli. A compact coordinate, or a pair of
coordinates, is called a \emph{block}. The planar sets will be images
under one compact coordinate, which is injective on $K$ and sends every
element of relative norm one to a unit vector. We count ordered pairs of
points at distance one, and divide by two at the end.

\subsection{The Class-Number Formula for Units of Relative Norm One}

Let $\mathcal O_K^\times$ and $\mathcal O_F^\times$ be the unit groups,
and put
\[
 \mathcal O_K^1=\ker\bigl(N_{K/F}\colon\mathcal O_K^\times\to
 \mathcal O_F^\times\bigr),
 \qquad I_N=[\mathcal O_F^\times:N_{K/F}(\mathcal O_K^\times)].
\]
Let $w_K$ be the number of roots of unity in $K$, let $h_K$ and $h_F$ be
the class numbers of $K$ and $F$, put $h_{\rm rel}=h_K/h_F$, let $\kappa$
be the order of the \emph{capitulation kernel}
$\ker\bigl(\Cl(F)\to\Cl(K)\bigr)$, the map being induced by extension of
ideals, and let
$L_F(s)=\zeta_K(s)/\zeta_F(s)$. At $s=1$, $L_F(1)$ denotes the quotient
of the residues of $\zeta_K$ and $\zeta_F$. The regulators
$\operatorname{Reg}_K$ and $\operatorname{Reg}_F$ are the ordinary ones,
which use twice the logarithm of the absolute value at a complex place
\cite[Chapter~5]{MilneANT}.

\begin{proposition}\label{geo:mass}
The roots of unity of $K$ lie in $\mathcal O_K^1$, and the restriction of
$l$ to $\mathcal O_K^1$ has kernel of order $w_K$ and image a full
lattice in $\R^c$. Let $\operatorname{Reg}^1$ be the covolume of this
lattice, with $\operatorname{Reg}^1=1$ if $c=0$. Then
\[
 \kappa=I_N/2^{b-1},\qquad
 \operatorname{Reg}_K/\operatorname{Reg}_F=2^{c-1}I_N\operatorname{Reg}^1,
\]
and consequently
\begin{equation}\label{geo:mass-density}
 \frac{w_K}{h_{\rm rel}\kappa\operatorname{Reg}^1}
 =\frac{2^{d-1}\pi^{b+c}}{\lambda^{d/2}L_F(1)}.
\end{equation}
\end{proposition}

\begin{proof}
\emph{Signs.} Fix a real place $v_0$ of $F$. By
\eqref{geo:norm-modulus}, $v_0(N_{K/F}x)=|\phi_{v_0}(x)|^2>0$ for
$x\in K^\times$, so $-1$ is not a norm from $K^\times$. If $\zeta\in K$
is a root of unity, then $N_{K/F}\zeta$ is a root of unity of $F$,
hence $\pm1$ because $F$ has a real place, and it is positive at $v_0$.
Thus $N_{K/F}\zeta=1$.

\emph{Logarithmic embeddings.} For $x\in\mathcal O_K^\times$ and
$y\in\mathcal O_F^\times$ let
\[
 \operatorname{Log}_K(x)=\bigl((2\log|\phi_vx|)_v,\,(2\log|\phi_j^+x|,
 2\log|\phi_j^-x|)_j\bigr),\qquad
 \operatorname{Log}_F(y)=\bigl((\log|v(y)|)_v,\,(2\log|\tau_j(y)|)_j\bigr).
\]
Their images are lattices in the hyperplanes of coordinate sum zero,
the kernel of $\operatorname{Log}_F$ is $\{\pm1\}$
\cite[Proposition~5.8]{MilneANT}, and $\operatorname{Reg}_K$ and
$\operatorname{Reg}_F$ are the covolumes of the images after one
coordinate is deleted; the choice of the deleted coordinate does not
matter.

\emph{The lattice $l(\mathcal O_K^1)$.} Let $\beta\in\mathcal O_K^1$. Then
$|\phi_v(\beta)|=1$ and $|\phi_j^\pm(\beta)|=e^{\pm l_j(\beta)}$, so
\[
 \operatorname{Log}_K(\beta)=\bigl(0,\ (2l_j(\beta),-2l_j(\beta))_j\bigr).
\]
Hence $l(\beta)=0$ exactly when every conjugate of $\beta$ has absolute
value one, that is, when $\beta$ is a root of unity
\cite[Corollary~5.6]{MilneANT}. By the first step the kernel of $l$ on
$\mathcal O_K^1$ has order $w_K$. The norm of $y\in\mathcal O_F^\times$
is $y^2$, so $N_{K/F}(\mathcal O_K^\times)$ has finite index in
$\mathcal O_F^\times$ and rank $b+c-1$. By the unit theorem
\cite[Theorem~5.1]{MilneANT}, $\mathcal O_K^\times$ has rank $d-1$, so
$\mathcal O_K^1$ has rank $c$. As $\operatorname{Log}_K(\mathcal O_K^1)$
is discrete, so is $l(\mathcal O_K^1)$, and it is a full lattice in
$\R^c$.

\emph{The capitulation kernel.} The Herbrand quotient
$h(M)=|\widehat H^0(\langle\iota\rangle,M)|/|H^1(\langle\iota\rangle,M)|$
\cite[Chapter~II, \S3]{MilneCFT} of the units satisfies
$2\,h(\mathcal O_K^\times)=\prod_v[K_w:F_v]$, where $v$ runs over the
archimedean places of $F$ and $w$ lies above $v$
\cite[Chapter~VII, Proposition~3.1]{MilneCFT}. Each real place
contributes $2$ and each complex place $1$, so
$h(\mathcal O_K^\times)=2^{b-1}$. Since
$\widehat H^0(\langle\iota\rangle,\mathcal O_K^\times)=\mathcal
O_F^\times/N_{K/F}(\mathcal O_K^\times)$ has order $I_N$, we get
$|H^1(\langle\iota\rangle,\mathcal O_K^\times)|=I_N/2^{b-1}$. This group
is $\mathcal O_K^1$ modulo the units $\gamma/\iota\gamma$ with
$\gamma\in\mathcal O_K^\times$. If $\mathfrak a$ is a fractional ideal
of $F$ with $\mathfrak a\mathcal O_K=x\mathcal O_K$, the class of the
unit $x/\iota x$ depends neither on the choice of $x$ nor on the choice
of $\mathfrak a$ in its ideal class. This defines a homomorphism from
$\ker(\Cl(F)\to\Cl(K))$ to $H^1(\langle\iota\rangle,\mathcal O_K^\times)$.
It is injective: if $x/\iota x=\gamma/\iota\gamma$, then $y=x/\gamma$
lies in $F$ and $\mathfrak a\mathcal O_K=y\mathcal O_K$, so
$\mathfrak a=y\mathcal O_F$ by unique factorization. It is surjective:
by Hilbert's Theorem~90 \cite[Chapter~II, Corollary~1.23]{MilneCFT}, an
element of $\mathcal O_K^1$ has the form $x/\iota x$ with
$x\in K^\times$, and then the fractional ideal
$x\mathcal O_K=(\iota x)\mathcal O_K$ is invariant under $\iota$. By
(G2), every $\iota$-invariant fractional ideal of $K$ is extended from
$F$: at a split prime $\mathfrak P\,\iota\mathfrak P$ invariance forces
equal exponents, an inert prime of $F$ stays prime, and no prime
ramifies. Hence $\kappa=I_N/2^{b-1}$.

\emph{Regulators.} Delete the coordinate of $v_0$ from both logarithmic
embeddings, and let $\mathcal U_K\subset\R^{d-1}$ and
$\mathcal U_F\subset\R^{b+c-1}$ be the resulting lattices, of covolumes
$\operatorname{Reg}_K$ and $\operatorname{Reg}_F$. By
\eqref{geo:norm-modulus}, $\operatorname{Log}_F(N_{K/F}x)$ is obtained
from $\operatorname{Log}_K(x)$ by keeping the coordinates of the real
places and adding the two coordinates of each pair. This induces a
linear map $T\colon\R^{d-1}\to\R^{b+c-1}$ such that
$T(\mathcal U_K)$ is the image of $N_{K/F}(\mathcal O_K^\times)$ in
$\mathcal U_F$. The substitution $(x_j,y_j)\mapsto(x_j,x_j+y_j)$ on each
pair has determinant one and turns $T$ into a coordinate projection
whose kernel consists of the $c$ coordinates $x_j$. The covolume of a
lattice in a product of two coordinate spaces is the covolume of its
intersection with the first factor times the covolume of its projection
to the second, whenever both are lattices.

A unit $x\in\mathcal O_K^\times$ lies in the kernel of $T$ exactly
when $\operatorname{Log}_F(N_{K/F}x)=0$, because the deleted coordinate
is minus the sum of the others. Then $N_{K/F}x=\pm1$, and $N_{K/F}x=1$
by the first step. So the intersection of $\mathcal U_K$ with the kernel is
the image of $\mathcal O_K^1$, which after the substitution is
$2l(\mathcal O_K^1)$, of covolume $2^c\operatorname{Reg}^1$. Since
$-1\notin N_{K/F}(\mathcal O_K^\times)$, the projection
$T(\mathcal U_K)$ has index
$[\mathcal O_F^\times:\{\pm1\}N_{K/F}(\mathcal O_K^\times)]=I_N/2$ in
$\mathcal U_F$. Hence
$\operatorname{Reg}_K=2^c\operatorname{Reg}^1\cdot\operatorname{Reg}_FI_N/2$.

\emph{The class-number formula.} The analytic class-number formula
\cite[Chapter~V, Theorem~2.4]{MilneCFT} (compare
\eqref{an:class-formula}), applied to $K$ (with $r_1=0$
and $r_2=d$) and to $F$ (with $r_1=b$, $r_2=c$ and $w_F=2$), gives
\[
 L_F(1)=\frac{(2\pi)^dh_K\operatorname{Reg}_K}{w_K\lambda^d}\cdot
 \frac{2\lambda^{d/2}}{2^b(2\pi)^ch_F\operatorname{Reg}_F}
 =\frac{2^{c+1}\pi^{b+c}h_{\rm rel}}{w_K\lambda^{d/2}}\,
 \frac{\operatorname{Reg}_K}{\operatorname{Reg}_F}.
\]
Substituting
$\operatorname{Reg}_K/\operatorname{Reg}_F=2^{c-1}I_N\operatorname{Reg}^1
=2^{b+c-2}\kappa\operatorname{Reg}^1$ gives \eqref{geo:mass-density}.
\end{proof}

\subsection{Norm-One Elements from Split Primes}

\begin{definition}\label{geo:uniform-types}
Let $\mathcal R$ be a finite set of rational primes. The extension $K/F$
has \emph{uniform local types} $(e_r,f_r)_{r\in\mathcal R}$ if, for every
$r\in\mathcal R$, every prime of $F$ above $r$ splits in $K$ and has
absolute type $(e_r,f_r)$ (Definition~\ref{tw:type}).
\end{definition}

If $K/F$ has uniform local types $(e_r,f_r)_{r\in\mathcal R}$, then $F$
has exactly $d/(e_rf_r)$ primes above $r$, each of norm $r^{f_r}$, since
the products $e_{\mathfrak p}f_{\mathfrak p}$ of the absolute ramification
indices and residue degrees of the primes $\mathfrak p$ of $F$ above $r$
add up to $[F:\Q]=d$. Given integers $k_r\ge0$, define
\begin{equation}\label{geo:HJ}
 H=\sum_{r\in\mathcal R}\frac{k_r\log r}{e_r},\qquad
 J=\sum_{r\in\mathcal R}\frac{\log(k_r+1)}{e_rf_r}.
\end{equation}

\begin{lemma}\label{geo:selection}
Assume that $K/F$ has uniform local types $(e_r,f_r)_{r\in\mathcal R}$, and
fix integers $k_r\ge0$. There are a
fractional ideal $I$ of $K$ with $N(I)=e^{-dH}$, an integer
$t\ge e^{dJ}/(\kappa h_{\rm rel})$, and elements
$\beta_1,\dots,\beta_t\in I$ of relative norm one whose cosets
$\beta_i\mathcal O_K^1$ are pairwise disjoint.
\end{lemma}

\begin{proof}
For each prime $\mathfrak p$ of $F$ above a prime $r\in\mathcal R$, write
$\mathfrak p\mathcal O_K=\mathfrak P_{\mathfrak p}\,\iota\mathfrak
P_{\mathfrak p}$ and $k_{\mathfrak p}=k_r$. For integer vectors
$\boldsymbol j=(j_{\mathfrak p})$ with $0\le j_{\mathfrak p}\le
k_{\mathfrak p}$ put
\[
 \mathfrak A_{\boldsymbol j}=\prod_{\mathfrak p}
 \mathfrak P_{\mathfrak p}^{j_{\mathfrak p}},\qquad
 \mathfrak J_{\boldsymbol j}=\prod_{\mathfrak p}
 \mathfrak P_{\mathfrak p}^{j_{\mathfrak p}}
 (\iota\mathfrak P_{\mathfrak p})^{k_{\mathfrak p}-j_{\mathfrak p}}.
\]
There are $\prod_{\mathfrak p}(k_{\mathfrak p}+1)=e^{dJ}$ such vectors.
The quotient of $\Cl(K)$ by the image of $\Cl(F)$ has order
$h_K\kappa/h_F=\kappa h_{\rm rel}$, so one class of this quotient
contains the classes of $\mathfrak A_{\boldsymbol j}$ for a set of
$t\ge e^{dJ}/(\kappa h_{\rm rel})$ vectors $\boldsymbol j$. Fix one of
them, $\boldsymbol j_0$, and put $I=\mathfrak J_{\boldsymbol j_0}^{-1}$.
For each $\boldsymbol j$ in the set, write
$\mathfrak A_{\boldsymbol j}\mathfrak A_{\boldsymbol j_0}^{-1}
=x\,\mathfrak a\mathcal O_K$ with $x\in K^\times$ and $\mathfrak a$ a
fractional ideal of $F$, and put $\beta_{\boldsymbol j}=x/\iota x$.
Then $N_{K/F}\beta_{\boldsymbol j}=1$, and since $\iota$ fixes
$\mathfrak a\mathcal O_K$,
\[
 \beta_{\boldsymbol j}\mathcal O_K
 =\mathfrak A_{\boldsymbol j}\mathfrak A_{\boldsymbol j_0}^{-1}
 \bigl(\iota\mathfrak A_{\boldsymbol j}\bigr)^{-1}
 \iota\mathfrak A_{\boldsymbol j_0}
 =\mathfrak J_{\boldsymbol j}I.
\]
As $\mathfrak J_{\boldsymbol j}$ is integral,
$\beta_{\boldsymbol j}\in I$. The ideals $\mathfrak J_{\boldsymbol j}I$
are distinct for distinct $\boldsymbol j$, and every element of
$\beta_{\boldsymbol j}\mathcal O_K^1$ generates
$\mathfrak J_{\boldsymbol j}I$, so the cosets are disjoint. Finally,
since each $\mathfrak p$ splits in $K$,
$N(I)^{-1}=\prod_{\mathfrak p}(N\mathfrak p)^{k_{\mathfrak p}}
=\prod_r r^{f_rk_rd/(e_rf_r)}=e^{dH}$.
\end{proof}

\subsection{Unit Averaging and Translation}
In this subsection and the next, $I$ is a fractional ideal of $K$ and
$\beta_1,\dots,\beta_t\in I$ are elements of relative norm one whose
cosets $\beta_i\mathcal O_K^1$ are pairwise disjoint;
Lemma~\ref{geo:selection} provides such data. For $h\in\R^c$ let
$W_h\colon\C^d\to\C^d$ multiply the coordinate $\phi_j^+$ by
$e^{-h_j}$ and $\phi_j^-$ by $e^{h_j}$, and fix the compact
coordinates. It is a real diagonal map of determinant one. Put
$\Lambda_h=W_hI_\infty$. For a lattice $\Lambda\subset\C^d$,
$\covol(\Lambda)$ denotes the volume of a fundamental domain. By
\cite[Proposition~4.26]{MilneANT}, extended to fractional ideals by
scaling and \cite[Proposition~4.2]{MilneANT},
\begin{equation}\label{geo:covolume}
 \covol(\mathfrak a_\infty)=2^{-d}|\Delta_K|^{1/2}N\mathfrak a
\end{equation}
for every fractional ideal $\mathfrak a$ of $K$. Hence
$\covol(\Lambda_h)=2^{-d}\lambda^dN(I)$, independently of $h$; for the
ideal of Lemma~\ref{geo:selection},
\begin{equation}\label{geo:determinant}
 \covol(\Lambda_h)=2^{-d}\lambda^de^{-dH}.
\end{equation}
If $N_{K/F}\beta=1$, then $W_h\beta_\infty$ has compact coordinates of
modulus one and pair coordinates of moduli $e^{\pm(l_j(\beta)-h_j)}$.
Let $\Pi^1$ be a fundamental parallelepiped of the lattice
$l(\mathcal O_K^1)$; its volume is $\operatorname{Reg}^1$. When $c=0$,
integrals over $\R^0$ and over $\Pi^1$ are evaluations at the single
point.

\begin{lemma}\label{geo:unfold-lemma}
For every Borel function $\Phi\colon\R^c\to[0,\infty]$,
\begin{equation}\label{geo:unfold}
 \frac1{\operatorname{Reg}^1}\int_{\Pi^1}\sum_{i=1}^t
 \sum_{\beta\in\beta_i\mathcal O_K^1}\Phi(l(\beta)-h)\,dh
 =\frac{tw_K}{\operatorname{Reg}^1}\int_{\R^c}\Phi(u)\,du .
\end{equation}
\end{lemma}

\begin{proof}
The map $l$ is a homomorphism on $K^\times$, so
$l(\beta_i\mathcal O_K^1)=l(\beta_i)+l(\mathcal O_K^1)$, and by
Proposition~\ref{geo:mass} each point of this translate of
$l(\mathcal O_K^1)$ is the image of exactly $w_K$ elements of
$\beta_i\mathcal O_K^1$. The translates $\Pi^1-x$, with
$x\in l(\mathcal O_K^1)$, tile $\R^c$, so by Tonelli's theorem
$\int_{\Pi^1}\sum_{x\in l(\mathcal
O_K^1)}\Phi(l(\beta_i)+x-h)\,dh=\int_{\R^c}\Phi(u)\,du$ for each
$i$.
\end{proof}

\begin{definition}\label{geo:window}
A \emph{window} is a bounded Borel set $\Omega\subset\C^d$ of positive
volume that is invariant under rotation of each coordinate separately.
For $u\in\R^c$ let $\nu(u)\in\C^d$ have compact coordinates $1$ and
pair coordinates $(e^{u_j},e^{-u_j})$, and for a Borel set
$\Omega\subset\C^d$ put
\[
 \Phi_\Omega(u)=|\Omega\cap(\Omega-\nu(u))|.
\]
\end{definition}

\begin{lemma}\label{geo:phi-omega}
Let $\Omega$ be a window, let $h\in\R^c$, and let $\beta\in K$ satisfy
$N_{K/F}\beta=1$. Then
$|\Omega\cap(\Omega-W_h\beta_\infty)|=\Phi_\Omega(l(\beta)-h)$.
\end{lemma}

\begin{proof}
If $x\in\C^d$ has compact coordinates of modulus one and pair
coordinates of moduli $e^{\pm u_j}$, then a rotation of the coordinates
preserves $\Omega$ and maps $\nu(u)$ to $x$, so
$|\Omega\cap(\Omega-x)|=\Phi_\Omega(u)$. By the remark before
Lemma~\ref{geo:unfold-lemma}, $x=W_h\beta_\infty$ is such a point, with
$u=l(\beta)-h$.
\end{proof}

For a window $\Omega$, the function $\Phi_\Omega$ is continuous, bounded
by $|\Omega|$, and vanishes when some $e^{|u_j|}$ exceeds the diameter of
$\Omega$: if $x$ and $x+\nu(u)$ both lie in $\Omega$, then $e^{|u_j|}$,
the modulus of a coordinate of their difference $\nu(u)$, is at most the
diameter of $\Omega$.

Let $\Omega$ be a window. Fix a $\Z$-basis of $I_\infty$, and for
$y\in[0,1)^{2d}$ let $y_I\in\C^d$ be the point with coordinate vector $y$
in this basis. For $h\in\R^c$ let
$\mathcal X(h,y)=(\Lambda_h+W_hy_I)\cap\Omega$ and
$n(h,y)=|\mathcal X(h,y)|$, and let $\mathcal E(h,y)$ be the number of
\emph{ordered edges} of $\mathcal X(h,y)$, that is, of pairs $(x,\beta)$
with $x\in\mathcal X(h,y)$, $\beta\in\bigcup_i\beta_i\mathcal O_K^1$ and
$x+W_h\beta_\infty\in\Omega$. Since $\beta\in I$, the point
$x+W_h\beta_\infty$ then lies in $\mathcal X(h,y)$.

\begin{lemma}\label{geo:joint-average-lemma}
Let $\Omega$ be a window and suppose $n(h,y)\le n_{\max}<\infty$ for
all $(h,y)$. Then there is a pair $(h,y)$ with $n(h,y)>0$ and
\begin{equation}\label{geo:joint-average}
 \frac{\mathcal E(h,y)}{n(h,y)}\ge
 \frac{tw_K}{\operatorname{Reg}^1}\,
 \frac{\int_{\R^c}\Phi_\Omega(u)\,du}{|\Omega|}.
\end{equation}
\end{lemma}

\begin{proof}
For fixed $x$, distinct $\beta$ give distinct points
$x+W_h\beta_\infty\ne x$, so $\mathcal E\le n^2\le n_{\max}^2$. Both $n$
and $\mathcal E$ are countable sums of indicator functions of Borel sets
in $(h,y)$. For fixed $h$, the map $y\mapsto W_hy_I$ has Jacobian
$\covol(\Lambda_h)$ and maps $[0,1)^{2d}$ onto a fundamental domain of
$\Lambda_h$, so tiling and Lemma~\ref{geo:phi-omega} give
\[
 \int_{[0,1)^{2d}}n(h,y)\,dy=\frac{|\Omega|}{\covol(\Lambda_h)},\qquad
 \int_{[0,1)^{2d}}\mathcal E(h,y)\,dy
 =\frac1{\covol(\Lambda_h)}\sum_{i}\sum_{\beta\in\beta_i\mathcal O_K^1}
 \Phi_\Omega(l(\beta)-h).
\]
Average over $h\in\Pi^1$ with respect to $dh/\operatorname{Reg}^1$ and
apply \eqref{geo:unfold}. For the probability measure
$dh\,dy/\operatorname{Reg}^1$ on $\Pi^1\times[0,1)^{2d}$ this gives
$\mathbb E\,\mathcal E=\rho\,\mathbb En$ and
$\mathbb En=|\Omega|/\covol(\Lambda_h)>0$, where $\rho$ is the right
side of \eqref{geo:joint-average}. If $\mathcal E<\rho n$ held wherever
$n>0$, then, since $\mathcal E=0$ where $n=0$ and $\{n>0\}$ has positive
measure, we would get $\mathbb E\,\mathcal E<\rho\,\mathbb En$.
\end{proof}

\begin{lemma}\label{geo:planar}
Let $\Omega$ be a window, let $(h,y)$ be a pair, and let $v_0$ be a real
place of $F$. For $x\in\C^d$ let $x_{v_0}$ be its compact coordinate at
$v_0$, and let $U=\{x_{v_0}:x\in\mathcal X(h,y)\}\subset\C=\R^2$. Then
$|U|=n(h,y)$ and $u(U)\ge\mathcal E(h,y)/2$.
\end{lemma}

\begin{proof}
Two distinct points of $\mathcal X(h,y)$ differ by $W_h\gamma_\infty$
with $0\ne\gamma\in I$, whose $v_0$-coordinate $\phi_{v_0}(\gamma)$ is
nonzero, so $|U|=n(h,y)$. An ordered edge $(x,\beta)$ gives two points of
$U$ at distance $|\phi_{v_0}(\beta)|=1$, and distinct ordered edges give
distinct ordered pairs of points. Hence $u(U)\ge\mathcal E(h,y)/2$.
\end{proof}

\subsection{Uniform Lattice Counting}
For a lattice $\Lambda\subset\C^d$ let
$\Lambda^*=\{\xi\in\C^d:\langle\xi,x\rangle\in\Z\text{ for all
}x\in\Lambda\}$ be its dual lattice, and let
$\|\xi\|_*=\sum_w|\xi_w|$, a norm on $\C^d$. We use the Fourier
transform $\widehat\Psi(\xi)=\int_{\C^d}\Psi(x)e^{-2\pi i\langle
x,\xi\rangle}\,dx$.

\begin{lemma}\label{geo:dual}
Let $\mathfrak a$ be a fractional ideal of $K$ and $h\in\R^c$. Every
nonzero $\xi$ in the dual lattice of $W_h\mathfrak a_\infty$ satisfies
\[
 \prod_w|\xi_w|\ge\frac{2^d}{\lambda^d(N\mathfrak a)^{1/2}},
 \qquad
 \|\xi\|_*\ge\frac{2d}{\lambda(N\mathfrak a)^{1/(2d)}}.
\]
In particular, if $N(I)=e^{-dH}$, as for the ideal of
Lemma~\ref{geo:selection}, every nonzero dual vector of $\Lambda_h$
satisfies
\begin{equation}\label{geo:dual-product}
 \prod_w|\xi_w|\ge\mu^d,\qquad \mu=2e^{H/2-\ell},\qquad
 \|\xi\|_*\ge d\mu.
\end{equation}
\end{lemma}

\begin{proof}
First let $h=0$. For $x,y\in K$ we have
$\langle x_\infty,\overline{y_\infty}\rangle
=\operatorname{Re}\sum_w\phi_w(xy)=\tfrac12\Tr_{K/\Q}(xy)$, where the
bar is coordinatewise complex conjugation and $\phi_w$ is the
embedding of the coordinate $w$, because these embeddings and their
complex conjugates are all the embeddings of $K$. Let
$\mathfrak a^\vee=\{y\in K:\Tr_{K/\Q}(y\mathfrak a)\subset\Z\}$. It is an
$\mathcal O_K$-module, and since the trace form is nondegenerate, the
dual basis of a $\Z$-basis of $\mathfrak a$ is a $\Z$-basis of
$\mathfrak a^\vee$. So $\mathfrak a^\vee$ is a fractional ideal, and
the dual lattice of $\mathfrak a_\infty$ is
$\overline{(2\mathfrak a^\vee)_\infty}$. Dual lattices have reciprocal
covolumes, and complex conjugation preserves volume. With
\eqref{geo:covolume} and $N(2\mathfrak a^\vee)=2^{2d}N(\mathfrak a^\vee)$, this
gives $N(\mathfrak a^\vee)=\lambda^{-2d}(N\mathfrak a)^{-1}$. If
$0\ne y\in\mathfrak a^\vee$, then $y\mathcal O_K\subset\mathfrak a^\vee$
and $|N_{K/\Q}y|=N(y\mathcal O_K)\ge N(\mathfrak a^\vee)$
\cite[Propositions~4.1 and~4.2]{MilneANT}. For
$\xi=\overline{(2y)_\infty}$,
\[
 \prod_w|\xi_w|=|N_{K/\Q}(2y)|^{1/2}
 \ge\bigl(2^{2d}\lambda^{-2d}(N\mathfrak a)^{-1}\bigr)^{1/2}.
\]
For general $h$, the dual lattice of $W_h\mathfrak a_\infty$ is the
image of the dual lattice of $\mathfrak a_\infty$ under $W_{-h}$,
because $W_h$ is real diagonal, and $W_{-h}$ preserves
$\prod_w|\xi_w|$. The bound for $\|\xi\|_*$ follows from the
arithmetic-geometric mean inequality. For \eqref{geo:dual-product} take
$\mathfrak a=I$.
\end{proof}

\begin{lemma}\label{geo:poisson}
Let $\Lambda\subset\C^d$ be a lattice such that
$\prod_w|\xi_w|\ge\mu^d$ for every nonzero $\xi\in\Lambda^*$, where
$\mu>0$. Let $\Psi\colon\C^d\to[0,\infty)$ be smooth, bounded and
integrable, with $\int\Psi>0$ and all partial derivatives bounded, and
suppose that
\[
 |\widehat\Psi(\xi)|\le M^de^{-\sigma\|\xi\|_*}\int\Psi
 \qquad(\xi\in\C^d)
\]
for some $M\ge1$ and $\sigma>0$. If the \emph{Fourier condition}
\begin{equation}\label{geo:poisson-condition}
 \sigma\mu\ge\log M+2\log5+2
\end{equation}
holds, then for every $y\in\C^d$,
\begin{equation}\label{geo:pointwise-poisson}
 \sum_{x\in\Lambda+y}\Psi(x)\le\frac2{\covol(\Lambda)}\int\Psi.
\end{equation}
\end{lemma}

\begin{proof}
A nonzero $\xi\in\Lambda^*$ has no zero coordinate, and the
arithmetic-geometric mean inequality gives $\|\xi\|_*\ge d\mu$. Hence
distinct points of $\Lambda^*$ are at $\|\cdot\|_*$-distance at least
$d\mu$, and the open $\|\cdot\|_*$-balls of radius $d\mu/2$ about them
are disjoint. For $j\ge1$, the balls about the points with
$jd\mu\le\|\xi\|_*<(j+1)d\mu$ lie in the ball of radius $(j+3/2)d\mu$
about $0$. Comparing volumes in $\R^{2d}$ bounds the number of these
points by $(2j+3)^{2d}\le5^{2dj}$. By the Fourier condition
\eqref{geo:poisson-condition} and
$M\ge1$,
\[
 M^d\sum_{\xi\in\Lambda^*\setminus\{0\}}e^{-\sigma\|\xi\|_*}
 \le\sum_{j\ge1}e^{d(\log M+2j\log5-j\sigma\mu)}
 \le\sum_{j\ge1}e^{-2dj}\le1 .
\]

For $\varepsilon>0$ put $\Psi_\varepsilon(x)=\Psi(x)e^{-\varepsilon|x|^2}$, a Schwartz
function. Its Fourier transform is $\widehat\Psi*G_\varepsilon$, where
$G_\varepsilon(\zeta)=(\pi/\varepsilon)^de^{-\pi^2|\zeta|^2/\varepsilon}$ is a probability
density on $\R^{2d}$. Hence
$|\widehat{\Psi_\varepsilon}(\xi)|\le m_\varepsilon M^de^{-\sigma\|\xi\|_*}\int\Psi$,
with $m_\varepsilon=\int e^{\sigma\|\zeta\|_*}G_\varepsilon(\zeta)\,d\zeta$. The function
$y\mapsto\sum_{x\in\Lambda}\Psi_\varepsilon(x+y)$ is smooth and
$\Lambda$-periodic, its Fourier coefficient at $\xi\in\Lambda^*$ is
$\widehat{\Psi_\varepsilon}(\xi)/\covol(\Lambda)$, and these coefficients are
absolutely summable. This gives the Poisson summation formula
\[
 \sum_{x\in\Lambda}\Psi_\varepsilon(x+y)
 =\frac1{\covol(\Lambda)}\sum_{\xi\in\Lambda^*}\widehat{\Psi_\varepsilon}(\xi)
 e^{2\pi i\langle\xi,y\rangle}.
\]
Using $\widehat{\Psi_\varepsilon}(0)=\int\Psi_\varepsilon$ and the bound above, the sum
is at most $(\int\Psi_\varepsilon+m_\varepsilon\int\Psi)/\covol(\Lambda)$. As
$\varepsilon\to0$, we have $m_\varepsilon\to1$ by dominated convergence, since
$G_\varepsilon$ is the image of $G_1$ under $\zeta\mapsto\sqrt\varepsilon\,\zeta$;
moreover $\Psi_\varepsilon$ increases to $\Psi$. Monotone convergence gives
\eqref{geo:pointwise-poisson}.
\end{proof}

The next bound is used in Section~\ref{pf:section} for the Gaussian
real-place profile.

\begin{lemma}\label{geo:gaussian}
If $0<a_\R\le1$ and $\Psi_\R(z)=e^{-a_\R|z|^2}$ on $\C$, then
$\widehat{\Psi_\R}(\xi)=(\pi/a_\R)e^{-\pi^2|\xi|^2/a_\R}$ and
$|\widehat{\Psi_\R}(\xi)|\le2e^{-|\xi|}\int\Psi_\R$.
\end{lemma}

\begin{proof}
The formula is the Gaussian integral, and $\int\Psi_\R=\pi/a_\R$. Since
$a_\R\le1$, it suffices that $\pi^2x^2-x+\log2\ge0$ for $x\ge0$, and the
minimum of the left side is $\log2-1/(4\pi^2)>0$.
\end{proof}

\subsection{Profiles at Real and Complex Places}

\begin{definition}\label{geo:profiles}
Fix $0<\delta<1$ and put $p=2/(1+\delta)$, the \emph{Lebesgue exponent}
(not a prime), so that $1<p<2$ and $p(1+\delta)=2$. A \emph{real-place
profile} is a positive radial function $f_\R$ on $\C$, used at the compact
coordinates, and a \emph{complex-place profile} is a positive function
$g_\C$ on $\C^2$ that depends only on $|z|$ and $|w|$, used at the pair
coordinates. Their \emph{masses} $A_\R$, $A_\C$, \emph{overlaps} $I_\R$,
$I_\C$ and \emph{functionals} $J_\R$, $J_\C$ are
\[
 A_\R=\int_\C f_\R^p,\qquad I_\R=\int_\C f_\R(z)f_\R(z+1)\,dz,
 \qquad A_\C=\int_{\C^2}g_\C^p,
\]
\[
 I_\C=\int_\R\int_{\C^2}g_\C(z,w)\,g_\C(z+e^u,w+e^{-u})\,dz\,dw\,du,
\]
\[
 J_\R=\log I_\R-(1+\delta)\log A_\R,\qquad
 J_\C=\log I_\C-(1+\delta)\log A_\C.
\]
When these integrals are finite and positive, the \emph{endpoint laws}
are the probability densities $f_\R(z)f_\R(z+1)/I_\R$ on $\C$ and
$g_\C(z,w)g_\C(z+e^u,w+e^{-u})/I_\C$ on $\C^2\times\R$. The functions
$V_\R=-\log f_\R$ and $V_\C=-\log g_\C$ are the \emph{energies} of the
profiles.
\end{definition}

The exponent $p$ is chosen so that the terms in the mean energy cancel at
the end of the proof of Proposition~\ref{geo:transfer}.

\begin{definition}\label{geo:admissible}
Let $M\ge1$ and $\sigma>0$. The pair $(f_\R,g_\C)$ is \emph{admissible
with Fourier constants $(M,\sigma)$} if:
\begin{enumerate}
\item[(P1)] $f_\R^p$ and $g_\C^p$ are smooth and bounded, and all their
partial derivatives are bounded;
\item[(P2)] $A_\R$, $I_\R$, $A_\C$ and $I_\C$ are finite and positive;
\item[(P3)] $V_\R$ and $V_\C$ are bounded below and coercive, that is,
their sublevel sets are bounded;
\item[(P4)] $V_\R(z)$ and $V_\C(z,w)$ have finite second moments under
the endpoint laws;
\item[(P5)] $|\widehat{f_\R^p}(\xi)|\le Me^{-\sigma|\xi|}A_\R$ and
$|\widehat{g_\C^p}(\xi_1,\xi_2)|\le M^2e^{-\sigma(|\xi_1|+|\xi_2|)}A_\C$
for all $\xi,\xi_1,\xi_2\in\C$.
\end{enumerate}
\end{definition}

Section~\ref{pf:section} proves admissibility for the profiles we use.

\begin{lemma}\label{geo:tensor}
Let $(f_\R,g_\C)$ be admissible with Fourier constants $(M,\sigma)$, and
let $\Psi=\prod_vf_\R^p\prod_jg_\C^p$ be the product function on $\C^d$,
with one factor for each block. Then $\Psi$ satisfies the hypotheses that
Lemma~\ref{geo:poisson} places on $\Psi$, other than the Fourier condition
\eqref{geo:poisson-condition}: it is smooth and bounded, with bounded
derivatives, $\int\Psi=A_\R^bA_\C^c$, and
\[
 |\widehat\Psi(\xi)|\le M^de^{-\sigma\|\xi\|_*}\int\Psi\qquad(\xi\in\C^d).
\]
\end{lemma}

\begin{proof}
The first properties follow from (P1) and (P2). The Fourier transform of
$\Psi$ is the product of the transforms of the factors, so (P5) gives
$|\widehat\Psi(\xi)|\le M^{b+2c}e^{-\sigma\|\xi\|_*}\int\Psi
=M^de^{-\sigma\|\xi\|_*}\int\Psi$.
\end{proof}

\begin{lemma}\label{geo:concentration}
Let $f_\R$ and $g_\C$ be profiles satisfying (P1), (P2) and (P4), and let
$\eta>0$. Let $m_\R$ and $m_\C$ be the means of $V_\R(z)$ and $V_\C(z,w)$
under the endpoint laws, and let $\mathrm{Var}_\R$ and $\mathrm{Var}_\C$ be
their variances. For $x\in\C^d$ with compact coordinates $x_v$ and pair
coordinates $(x_j^+,x_j^-)$, let
$V(x)=\sum_vV_\R(x_v)+\sum_jV_\C(x_j^+,x_j^-)$ be the total energy, put
$\bar m=(bm_\R+cm_\C)/d$, and let
\[
 \Omega=\bigl\{x\in\C^d:V(x)\le d(\bar m+\eta)\bigr\}.
\]
Then, for all $d$ larger than a bound that depends only on $\eta$,
$\mathrm{Var}_\R$ and $\mathrm{Var}_\C$,
\[
 \int_{\R^c}\Phi_\Omega(u)\,du\ge\tfrac12I_\R^bI_\C^ce^{2d(\bar m-\eta)}.
\]
\end{lemma}

\begin{proof}
By (P1), $f_\R$ and $g_\C$ are continuous, so $V$ is continuous and
$\Omega$ is closed, hence a Borel set. The function
$(x,u)\mapsto e^{-V(x)}e^{-V(x+\nu(u))}/(I_\R^bI_\C^c)$ is the
product of the endpoint densities of the blocks, hence a probability
density on $\C^d\times\R^c$. Under it, the energies $V(x)$ and
$V(x+\nu(u))$ of the two endpoints are sums of independent terms, one for
each block, and $V(x)$ has mean $d\bar m$. On a compact block the map
$z\mapsto-z-1$, and on a pair block the map
$(z,w,u_j)\mapsto(-z-e^{u_j},-w-e^{-u_j},u_j)$, preserves the endpoint
density and carries the first endpoint to minus the second. Since the
profiles are even, $V(x+\nu(u))$ has the same law as $V(x)$. By
Chebyshev's inequality, each of the two energies differs from $d\bar m$
by more than $\eta d$ with probability at most
$\max(\mathrm{Var}_\R,\mathrm{Var}_\C)/(\eta^2d)$. So the set
$\mathcal Y\subset\C^d\times\R^c$ where both energies lie in
$[d(\bar m-\eta),d(\bar m+\eta)]$ has probability at least $1/2$ for all
$d$ larger than a bound that depends only on $\eta$, $\mathrm{Var}_\R$ and
$\mathrm{Var}_\C$. On $\mathcal Y$ both endpoints lie in $\Omega$ and
$e^{-V(x)}e^{-V(x+\nu(u))}\le e^{-2d(\bar m-\eta)}$. Therefore, with
$|\mathcal Y|$ the Lebesgue measure of $\mathcal Y$,
\[
 \int_{\R^c}\Phi_\Omega(u)\,du\ge|\mathcal Y|
 \ge e^{2d(\bar m-\eta)}I_\R^bI_\C^c\,\mathbb P(\mathcal Y)
 \ge\tfrac12I_\R^bI_\C^ce^{2d(\bar m-\eta)}.
\]
\end{proof}

\begin{proposition}[Geometric transfer]\label{geo:transfer}
Let $0<\delta<1$ and $C\in\R$. Let $(K_i/F_i)_{i\ge1}$ be quadratic
extensions satisfying (G1)--(G3) with a common $\lambda$ and with
$d_i=[F_i:\Q]\to\infty$, such that $\log L_{F_i}(1)\le d_iC$ and such
that every $K_i/F_i$ has the same uniform local types
$(e_r,f_r)_{r\in\mathcal R}$. Fix integers $k_r\ge0$, let $H$ and $J$ be
as in \eqref{geo:HJ}, and let $\mu=2e^{H/2-\ell}$. Let $(f_\R,g_\C)$ be
admissible with Fourier constants $(M,\sigma)$ satisfying the Fourier
condition \eqref{geo:poisson-condition}. Put
\begin{equation}\label{geo:general-margin}
 \mathcal M(\theta)=J-\delta H-(1/2-\delta)\ell-C
 +(1-\delta)\log2+(1-\theta)\log\pi
 +(1-2\theta)J_\R+\theta J_\C.
\end{equation}
We call $\mathcal M(\theta)$ the \emph{margin}. If
$\inf_i\mathcal M(\theta_i)>0$, where $\theta_i$ is the value of
$\theta$ for $F_i$, then there are finite sets $U_i\subset\R^2$ with
$|U_i|\to\infty$ and $u(U_i)/|U_i|^{1+\delta}\to\infty$.
\end{proposition}

\begin{proof}
\emph{The window.} Let $m_\R$, $m_\C$, $\mathrm{Var}_\R$ and
$\mathrm{Var}_\C$ be as in Lemma~\ref{geo:concentration}; the variances
are finite by (P4). Fix $\eta>0$ with
$4\eta<\inf_i\mathcal M(\theta_i)$, and consider one extension $K/F$ of
the sequence, with $d$ large. Take $I$, $t$ and
$\beta_1,\dots,\beta_t$ from Lemma~\ref{geo:selection}, so that
$\Lambda_h$ is defined; the counts $n(h,y)$ and $\mathcal E(h,y)$ are those
of the window $\Omega$ chosen next. With the total energy $V$ and the mean
$\bar m$ of Lemma~\ref{geo:concentration}, let
\[
 \Omega=\bigl\{x\in\C^d:V(x)\le d(\bar m+\eta)\bigr\}.
\]
By (P1) and (P3) (see the proof of Lemma~\ref{geo:concentration}), this
is a bounded Borel set, and it is invariant under
coordinate rotations because the profiles are radial in each
coordinate. By Lemma~\ref{geo:concentration},
\[
 \int_{\R^c}\Phi_\Omega(u)\,du\ge\tfrac12I_\R^bI_\C^ce^{2d(\bar m-\eta)}
\]
for all large $d$, uniformly in the signature. In particular
$|\Omega|>0$ for large $d$, so $\Omega$ is a window.

\emph{Counting points.} With the product function
$\Psi=\prod_vf_\R^p\prod_jg_\C^p=e^{-pV}$ of Lemma~\ref{geo:tensor}, we
have $\mathbf 1_\Omega\le e^{pd(\bar m+\eta)}\Psi$, so
\[
 |\Omega|\le e^{pd(\bar m+\eta)}A_\R^bA_\C^c,
\]
and Lemmas~\ref{geo:dual}, \ref{geo:poisson} and~\ref{geo:tensor} give,
for all $(h,y)$,
\[
 n(h,y)\le\frac{2e^{pd(\bar m+\eta)}A_\R^bA_\C^c}{\covol(\Lambda_h)}
 =2e^{dY_\eta},\qquad
 Y_\eta=H+\log2-\ell+(1-2\theta)\log A_\R+\theta\log A_\C
 +p(\bar m+\eta),
\]
by \eqref{geo:determinant}, $b/d=1-2\theta$ and $c/d=\theta$.

\emph{Counting edges.} Now take a pair $(h,y)$ from
Lemma~\ref{geo:joint-average-lemma}, applied with
$n_{\max}=2e^{dY_\eta}$. By $t\ge e^{dJ}/(\kappa h_{\rm rel})$,
Proposition~\ref{geo:mass}, $\log L_F(1)\le dC$ and $b+c=(1-\theta)d$,
\[
 \frac{tw_K}{\operatorname{Reg}^1}\ge
 \frac{e^{dJ}2^{d-1}\pi^{b+c}}{\lambda^{d/2}L_F(1)}
 \ge\tfrac12\exp\bigl(d[J+\log2+(1-\theta)\log\pi-\ell/2-C]\bigr).
\]
Combined with the bounds for $|\Omega|$ and for the integral of
$\Phi_\Omega$, \eqref{geo:joint-average} shows that the set
$\mathcal X=\mathcal X(h,y)$ has $n=|\mathcal X|>0$ points and at least
$\tfrac14e^{dX_\eta}n$ ordered edges, where
\[
 \begin{aligned}
 X_\eta={}&J+\log2+(1-\theta)\log\pi-\ell/2-C
 +(1-2\theta)(\log I_\R-\log A_\R)\\
 &+\theta(\log I_\C-\log A_\C)+(2-p)\bar m-(2+p)\eta.
 \end{aligned}
\]
Since $n\le2e^{dY_\eta}$ and $p(1+\delta)=2$, the terms in $\bar m$ cancel,
and
\[
 \frac{\mathcal E(h,y)}{n^{1+\delta}}
 \ge\frac{e^{dX_\eta}}{4\cdot2^\delta e^{\delta dY_\eta}}
 =2^{-2-\delta}e^{d(\mathcal M(\theta)-4\eta)}.
\]

\emph{The planar set.} Finally let $v_0$ be a real place of $F$, and let
$U=\{x_{v_0}:x\in\mathcal X\}\subset\C=\R^2$. By
Lemma~\ref{geo:planar}, $|U|=n$ and $u(U)\ge\mathcal E(h,y)/2$, so
\[
 \frac{u(U)}{|U|^{1+\delta}}
 \ge2^{-3-\delta}e^{d(\mathcal M(\theta)-4\eta)},
\]
which tends to infinity along the sequence. Since $u(U)\le|U|^2/2$ and
$\delta<1$, also $|U|\to\infty$. These sets are the $U_i$. The
conclusion concerns a sequence of cardinalities, not every large
cardinality.
\end{proof}

\section{Weighted Windows at Split Finite Places}
\label{fw:section}

We now place locally constant weights at the selected finite places.
This refines the choice of ideals in Lemma~\ref{geo:selection} in the
same way that the archimedean profiles refine a ball. We keep the
hypotheses (G1)--(G3) and the notation of Section~\ref{geo:section},
and put
\[
 \covol_0=\covol\bigl((\mathcal O_K)_\infty\bigr)=(\lambda/2)^d .
\]

\subsection{The Local Functional}

Let $L$ be a nonarchimedean local field with valuation ring $\mathcal
O$, uniformizer $\varpi$ and residue field of cardinality $Q$. We use
the additive Haar measure on $L$ with $|\mathcal O|=1$, the product
measure on $L^2$, and the multiplicative Haar measure $d\omega$ on
$\mathcal O^\times$ of total mass one. For $n\in\Z$ and
$\omega\in\mathcal O^\times$ put
\[
 s(n,\omega)=(\varpi^n\omega,\ \varpi^{-n}\omega^{-1})\in L^2,\qquad
 Tg(z)=\sum_{n\in\Z}\int_{\mathcal O^\times}g(z+s(n,\omega))\,d\omega .
\]
The two coordinates of $s(n,\omega)$ have product one; they model the
finite coordinates of an element of relative norm one at a split place.
For a nonnegative, locally constant, compactly supported function $g$
on $L^2$ with $g\ne0$, and $0<\delta<1$, $p=2/(1+\delta)$, put
\begin{equation}\label{fw:functional}
 A(g)=\int_{L^2}g^p,\qquad Z(g)=\int_{L^2}g\,Tg,\qquad
 \mathcal F_\delta(g)=\frac{Z(g)}{A(g)^{1+\delta}} .
\end{equation}
Only finitely many $n$ contribute to $Z(g)$: if $g(z)g(z+s(n,\omega))\ne0$
then both coordinates of $s(n,\omega)$ lie in the compact set
$\operatorname{supp}g-\operatorname{supp}g$, which bounds $n$ from above
and from below.

For $i\in\Z$ let $\mathcal B_i=\varpi^{-i}\mathcal O$, a ball of measure
$Q^i$. Write $[x]_+=\max(x,0)$.

\begin{lemma}\label{fw:gram-lemma}
For all integers $i_1,j_1,i_2,j_2$,
\begin{equation}\label{fw:rectangle-gram}
 \int_{L^2}\mathbf 1_{\mathcal B_{i_1}\times\mathcal B_{j_1}}\,
 T\mathbf 1_{\mathcal B_{i_2}\times\mathcal B_{j_2}}
 =Q^{\min(i_1,i_2)+\min(j_1,j_2)}\,
 [\max(i_1,i_2)+\max(j_1,j_2)+1]_+.
\end{equation}
\end{lemma}

\begin{proof}
Two balls of an ultrametric space are disjoint or nested. Hence
$\mathcal B_{i_1}\cap(\mathcal B_{i_2}-x)$ has measure
$Q^{\min(i_1,i_2)}$ if $x\in\mathcal B_{\max(i_1,i_2)}$, and is empty
otherwise. The left side of \eqref{fw:rectangle-gram} is therefore
$Q^{\min(i_1,i_2)+\min(j_1,j_2)}$ times the measure of the set of
$(n,\omega)$ with $\varpi^n\omega\in\mathcal B_{\max(i_1,i_2)}$ and
$\varpi^{-n}\omega^{-1}\in\mathcal B_{\max(j_1,j_2)}$. These conditions
say $-\max(i_1,i_2)\le n\le\max(j_1,j_2)$ and do not involve $\omega$,
and the unit integral has mass one.
\end{proof}

Put $\mathcal S_0=\mathcal O$ and $\mathcal S_i=\mathcal B_i\setminus
\mathcal B_{i-1}$ for $i\ge1$, the \emph{shells}, of measures $m_0=1$
and $m_i=Q^i-Q^{i-1}$.

\begin{definition}\label{fw:shell-profile}
Let $k\ge0$ be an integer, and let $w=(w_{ij})_{i,j\ge0}$ be nonnegative
weights, finitely many of them nonzero, with $w_{00}>0$. The
\emph{shell profile} of type $(Q,k,w)$ on $L^2$ is
\[
 g=\sum_{i,j\ge0}w_{ij}\mathbf 1_{\mathcal S_i\times\varpi^{-k}\mathcal S_j}.
\]
\end{definition}

A shell profile is nonnegative, locally constant and compactly
supported. It is invariant under $(x,y)\mapsto(\omega x,\omega'y)$ for
all $\omega,\omega'\in\mathcal O^\times$, since multiplication by a unit
preserves valuations and hence each shell; in particular it is even. It
is also invariant under translation by the group
\begin{equation}\label{fw:hard-period}
 \mathcal C=\mathcal O\times\varpi^{-k}\mathcal O,\qquad |\mathcal C|=Q^k,
\end{equation}
whatever the number of shells: adding an element of $\mathcal O$ does
not change the shell of a point. We call $\mathcal C$ the \emph{period}
of $g$. Define the symmetric matrix $\Delta=(\Delta_{ii'})_{i,i'\ge0}$ by
\[
 \Delta_{ii'}=m_{\min(i,i')}\ (i\ne i'),\qquad \Delta_{00}=0,\qquad
 \Delta_{ii}=im_i-Q^{i-1}\ (i\ge1),
\]
and the kernel
\[
 \mathcal K_{(i,j),(i',j')}=(k+1)m_im_j\mathbf 1_{i=i',\,j=j'}
 +m_j\Delta_{ii'}\mathbf 1_{j=j'}+m_i\Delta_{jj'}\mathbf 1_{i=i'}.
\]

\begin{lemma}\label{fw:shell-lemma}
The shell profile $g$ of type $(Q,k,w)$ satisfies
\begin{equation}\label{fw:shell-kernel}
 A(g)=Q^k\sum_{i,j}m_im_jw_{ij}^p,\qquad
 Z(g)=Q^k\sum_{(i,j),(i',j')}w_{ij}\,\mathcal K_{(i,j),(i',j')}\,w_{i'j'} .
\end{equation}
In particular $\mathcal F_\delta(g)$ depends only on the type, and the
hard window $w_{00}=1$, $w_{ij}=0$ otherwise, has
$\mathcal F_\delta(g)=(k+1)Q^{-\delta k}$.
\end{lemma}

\begin{proof}
The formula for $A(g)$ holds because the sets
$\mathcal S_i\times\varpi^{-k}\mathcal S_j$ are disjoint of measure
$Q^km_im_j$. For $Z(g)$, write $\mathbf 1_{\mathcal S_0}=\mathbf
1_{\mathcal B_0}$ and $\mathbf 1_{\mathcal S_i}=\mathbf 1_{\mathcal
B_i}-\mathbf 1_{\mathcal B_{i-1}}$ for $i\ge1$, and note
$\varpi^{-k}\mathcal B_j=\mathcal B_{j+k}$. Only balls with nonnegative
indices occur, so by \eqref{fw:rectangle-gram} and $k\ge0$ the positive
part can be dropped, and the Gram value of $\mathcal B_{i_1}\times
\mathcal B_{j_1+k}$ and $\mathcal B_{i_2}\times\mathcal B_{j_2+k}$, for
$i_1,j_1,i_2,j_2\ge0$, is
\[
 Q^k\bigl[(k+1)Q^{\min(i_1,i_2)}Q^{\min(j_1,j_2)}
 +\max(i_1,i_2)Q^{\min(i_1,i_2)}Q^{\min(j_1,j_2)}
 +Q^{\min(i_1,i_2)}\max(j_1,j_2)Q^{\min(j_1,j_2)}\bigr].
\]
For a function $\varphi$ of two nonnegative integers put
$(\partial\varphi)(i,i')=\varphi(i,i')-\varphi(i-1,i')-\varphi(i,i'-1)
+\varphi(i-1,i'-1)$, where terms with an index $-1$ are omitted; this is
the expansion of the shells into balls. Apply $\partial$ in
$(i_1,i_2)$ and in $(j_1,j_2)$. A direct computation gives, for
$i,i'\ge0$, that $\partial Q^{\min(i_1,i_2)}$ at $(i,i')$ is
$m_i\mathbf 1_{i=i'}$, and that
$\partial\bigl(\max(i_1,i_2)Q^{\min(i_1,i_2)}\bigr)$ is $\Delta_{ii'}$.
For example, when $1\le i<i'$ the latter is
$i'Q^i-i'Q^{i-1}-(i'-1)Q^i+(i'-1)Q^{i-1}=m_i$, and when $i=i'\ge1$ it is
$iQ^i-2iQ^{i-1}+(i-1)Q^{i-1}=im_i-Q^{i-1}$. This gives the formula for
$Z(g)$. For the hard window only $\mathcal K_{(0,0),(0,0)}=k+1$
contributes.
\end{proof}

All entries of $\mathcal K$ are nonnegative, since
$\Delta_{ii}=Q^{i-1}(iQ-i-1)\ge0$ for $i\ge1$. Hence
$Z(g)\ge(k+1)Q^kw_{00}^2>0$ for every shell profile.

For product weights $w_{ij}=x_ix'_j$, with vectors $x,x'$ of
nonnegative numbers, put
\[
 P_x=\sum_im_ix_i^p,\qquad S_x=\sum_im_ix_i^2,\qquad
 R_x=\sum_{i,i'}x_i\Delta_{ii'}x_{i'},
\]
and define $P_{x'},S_{x'},R_{x'}$ in the same way. Then
\eqref{fw:shell-kernel} becomes
\begin{equation}\label{fw:product-shells}
 A(g)=Q^kP_xP_{x'},\qquad
 Z(g)=Q^k\bigl[(k+1)S_xS_{x'}+R_xS_{x'}+S_xR_{x'}\bigr].
\end{equation}
For example, the one-shell choice $x=x'=(1,t)$ has $S_x=1+(Q-1)t^2$ and
$R_x=2t+(Q-2)t^2$, and its gain over the hard window is
\begin{equation}\label{fw:one-shell-gain}
 \log\frac{(k+1)S_x^2+2S_xR_x}{k+1}-2(1+\delta)\log\bigl(1+(Q-1)t^p\bigr).
\end{equation}
Since $p>1$, its right derivative at $t=0$ is $4/(k+1)>0$, so a
sufficiently small first shell always improves on the hard window. The
certificate uses the exact finite functional \eqref{fw:product-shells}.

\subsection{Finite Periods and Lattice Counting}

Let $\mathcal P$ be a finite set of finite places of $F$ that split in
$K$. For $v\in\mathcal P$ write $v\mathcal O_K=\mathfrak P_v\,\iota\mathfrak
P_v$, where we also write $v$ for the prime ideal of $F$. The completion
of $K$ at $\mathfrak P_v$ is $F_v$; let $j_v\colon K\to F_v$ be the
completion map. The map $x\mapsto(j_v(x),j_v(\iota x))$ records the
completions at $\mathfrak P_v$ and at $\iota\mathfrak P_v$, and $\iota$
exchanges the two coordinates. If $N_{K/F}\beta=1$, then
$j_v(\beta)j_v(\iota\beta)=1$. Let $\mathcal O_v$, $\varpi_v$ and $Q_v$
be the valuation ring, a uniformizer and the residue cardinality of
$F_v$. Put
\[
 \mathbb A_{\mathcal P}=\prod_{v\in\mathcal P}(F_v\times F_v),
\]
with the Haar measure that gives $\prod_v(\mathcal O_v\times\mathcal
O_v)$ measure one, and let $\Gamma=\mathcal O_{K,\mathcal P}$ be the
ring of elements of $K$ that are integral at every prime other than the
$\mathfrak P_v$ and $\iota\mathfrak P_v$. We embed $\Gamma$ in
$\C^d\times\mathbb A_{\mathcal P}$ by
$\gamma\mapsto(\gamma_\infty,\gamma_f)$, where
$\gamma_f=(j_v(\gamma),j_v(\iota\gamma))_v$. A subgroup of $\mathbb
A_{\mathcal P}$ of the form
\[
 \mathcal C_f=\prod_{v\in\mathcal P}\bigl(\varpi_v^{a_v}\mathcal
 O_v\times\varpi_v^{b_v}\mathcal O_v\bigr),\qquad a_v,b_v\in\Z,
\]
is called \emph{rectangular}; its measure is
$|\mathcal C_f|=\prod_vQ_v^{-a_v-b_v}$.

\begin{lemma}\label{fw:lattice}
Let $\mathcal C_f$ be rectangular and put
$\mathfrak L=\prod_v\mathfrak P_v^{a_v}(\iota\mathfrak P_v)^{b_v}$.
\begin{enumerate}
\item[(a)] The elements $\gamma\in\Gamma$ with $\gamma_f\in\mathcal C_f$
form $\mathfrak L$, and $N\mathfrak L=|\mathcal C_f|^{-1}$, so
$\covol(\mathfrak L_\infty)=\covol_0/|\mathcal C_f|$.
\item[(b)] Every coset of $\mathcal C_f$ in $\mathbb A_{\mathcal P}$
contains $\gamma_f$ for some $\gamma\in\Gamma$.
\item[(c)] $\Gamma$ is a lattice in $\C^d\times\mathbb A_{\mathcal P}$.
If $P_{\mathfrak L}$ is a fundamental parallelepiped of
$\mathfrak L_\infty$, then $P_{\mathfrak L}\times\mathcal C_f$ is a
fundamental domain for $\Gamma$, of measure $\covol_0$.
\end{enumerate}
\end{lemma}

\begin{proof}
(a) The valuation of $j_v(\gamma)$ is the exponent of $\mathfrak P_v$
in $\gamma$, and that of $j_v(\iota\gamma)$ is the exponent of
$\iota\mathfrak P_v$. Together with integrality at the other primes,
$\gamma_f\in\mathcal C_f$ says exactly that $\gamma\in\mathfrak L$. As
$N\mathfrak P_v=N(\iota\mathfrak P_v)=Q_v$, we get
$N\mathfrak L=\prod_vQ_v^{a_v+b_v}=|\mathcal C_f|^{-1}$, and the
covolume follows from the formula of Section~\ref{geo:section}.

(b) Choose $a'_v\le a_v$ and $b'_v\le b_v$ such that the given coset
lies in $\mathcal C'_f=\prod_v(\varpi_v^{a'_v}\mathcal O_v\times
\varpi_v^{b'_v}\mathcal O_v)$, and let $\mathfrak L'\supset\mathfrak L$
be the corresponding ideal. The map $\gamma\mapsto\gamma_f$ induces a
homomorphism $\mathfrak L'/\mathfrak L\to\mathcal C'_f/\mathcal C_f$. It
is injective by (a), applied to $\mathcal C_f$. Both groups are finite
of the same order: $|\mathcal C'_f/\mathcal C_f|=|\mathcal
C'_f|/|\mathcal C_f|$, and
$|\mathfrak L'/\mathfrak L|=N\mathfrak L/N\mathfrak L'$
\cite[Proposition~4.2]{MilneANT}, which is the same number by (a). So
the map is bijective, and every coset of $\mathcal C_f$ in $\mathcal
C'_f$ contains $\gamma_f$ for some $\gamma\in\mathfrak L'\subset\Gamma$.
No principality is needed.

(c) Given $(z,c)\in\C^d\times\mathbb A_{\mathcal P}$, by (b) there is
$\gamma\in\Gamma$ with $c-\gamma_f\in\mathcal C_f$, and then a unique
$\gamma'\in\mathfrak L$ with $z-(\gamma+\gamma')_\infty\in P_{\mathfrak
L}$. If $(z,c)$ had two such representations, the difference of the
two elements of $\Gamma$ would have finite part in $\mathcal C_f$, hence
lie in $\mathfrak L$ by (a), and then it would vanish because
$P_{\mathfrak L}$ is a fundamental domain. Thus
$P_{\mathfrak L}\times\mathcal C_f$ is a fundamental domain, of measure
$(\covol_0/|\mathcal C_f|)\cdot|\mathcal C_f|=\covol_0$. Finally $\Gamma$
is discrete: for a bounded open set $Y\ni0$ in $\C^d$, its intersection
with the neighborhood $Y\times\mathcal C_f$ of $0$ consists of the
finitely many $\gamma\in\mathfrak L$ with $\gamma_\infty\in Y$.
\end{proof}

For $h\in\R^c$ let $\Gamma_h=\{(W_h\gamma_\infty,\gamma_f):\gamma\in
\Gamma\}$. It is again a lattice, with fundamental domains
$W_hP_{\mathfrak L}\times\mathcal C_f$ of measure $\covol_0$. For a
rectangular $\mathcal C_f$ put
\[
 H_f=d^{-1}\log|\mathcal C_f|.
\]
By Lemma~\ref{fw:lattice}(a), $N\mathfrak L=e^{-dH_f}$, and
Lemma~\ref{geo:dual} gives the dual product bound $\mu_f^d$ for every
$W_h\mathfrak L_\infty$, with $\mu_f=2e^{H_f/2-\ell}$. Accordingly, the
Fourier condition used in this section is
\begin{equation}\label{fw:fourier-condition}
 \sigma\cdot2e^{H_f/2-\ell}\ge\log M+2\log5+2 .
\end{equation}

\begin{lemma}\label{fw:counting}
Let $(f_\R,g_\C)$ be admissible with Fourier constants $(M,\sigma)$, and
let $\Psi_\infty=\prod_vf_\R^p\prod_jg_\C^p$ on $\C^d$. Let
$\Psi_f\ge0$ be a function on $\mathbb A_{\mathcal P}$ that is invariant
under translation by a rectangular group $\mathcal C_f$ and vanishes
outside a finite union of its cosets. If \eqref{fw:fourier-condition}
holds, with $H_f$ defined by this $\mathcal C_f$, then for every
$h\in\R^c$ and every $z_0\in\C^d\times\mathbb A_{\mathcal P}$,
\begin{equation}\label{fw:weighted-count}
 \sum_{\gamma\in\Gamma_h+z_0}\Psi_\infty(\gamma_\infty)\Psi_f(\gamma_f)
 \le\frac2{\covol_0}\int\Psi_\infty\int\Psi_f ,
\end{equation}
where $\gamma_\infty$ and $\gamma_f$ are the two components of $\gamma$.
\end{lemma}

\begin{proof}
Group the points by the coset of $\mathcal C_f$ containing their finite
component. If two points of $\Gamma_h+z_0$ have finite components in
the same coset, their difference comes from an element of $\mathfrak L$,
by Lemma~\ref{fw:lattice}(a). So the archimedean components of the
points in one coset form a translate of $W_h\mathfrak L_\infty$, or the
empty set. By Lemma~\ref{geo:poisson}, with the dual bound $\mu_f$ and
the Fourier bound for $\Psi_\infty$ from Section~\ref{geo:section}, the
sum of $\Psi_\infty$ over such a translate is at most
$2|\mathcal C_f|\int\Psi_\infty/\covol_0$. The values of $\Psi_f$ on its
finitely many nonzero cosets sum to $\int\Psi_f/|\mathcal C_f|$.
\end{proof}

If $\alpha\in\mathbb A_{\mathcal P}$ has coordinates
$(\varpi_v^{n_v},\varpi_v^{-n_v})$, multiplication by $\alpha$ preserves
the Haar measure and maps rectangular groups to rectangular groups of
the same measure. So \eqref{fw:fourier-condition} is unchanged, and
Lemma~\ref{fw:counting} applies after such a rescaling of $\Psi_f$,
even though the rescaling does not preserve $\Gamma$.

\subsection{Class Selection and Unit Averaging}

For $u\in\R^c$ and $n=(n_v)\in\Z^{\mathcal P}$ let
$\nu(u,n)\in\C^d\times\mathbb A_{\mathcal P}$ have archimedean part
$\nu(u)$ and finite part $(\varpi_v^{n_v},\varpi_v^{-n_v})_v$.

\begin{lemma}\label{fw:selection}
Let $\mathcal N\subset\Z^{\mathcal P}$ be finite, and let
$\mathcal W\colon\mathcal N\to[0,\infty)$ have positive sum. There are
$n_0\in\mathcal N$ with $\mathcal W(n_0)>0$, a subset
$\mathcal N'\subset\mathcal N$ with
\[
 \sum_{n\in\mathcal N'}\mathcal W(n)\ge\frac1{\kappa h_{\rm rel}}
 \sum_{n\in\mathcal N}\mathcal W(n),
\]
and elements $\beta_n\in\Gamma$, for $n\in\mathcal N'$, with
$N_{K/F}\beta_n=1$ and
\[
 \beta_n\mathcal O_K=\prod_v\mathfrak P_v^{\,n_v-n_{0,v}}
 (\iota\mathfrak P_v)^{-(n_v-n_{0,v})}.
\]
The cosets $\beta_n\mathcal O_K^1$ are pairwise disjoint, and
$\beta_n\mathcal O_K^1$ is the set of all elements of relative norm one
that generate this ideal.
\end{lemma}

\begin{proof}
Map $n$ to the class of $\mathfrak A(n)=\prod_v\mathfrak P_v^{n_v}$ in
$\Cl(K)$ modulo the image of $\Cl(F)$, a group of order $\kappa h_{\rm
rel}$. Let $\mathcal N'$ be a fiber of largest total weight, and choose
$n_0\in\mathcal N'$ with $\mathcal W(n_0)>0$. For $n\in\mathcal N'$ write
$\mathfrak A(n)\mathfrak A(n_0)^{-1}=x_n\mathfrak a_n\mathcal O_K$ with
$x_n\in K^\times$ and $\mathfrak a_n$ a fractional ideal of $F$, and put
$\beta_n=x_n/\iota x_n$. As in the proof of Lemma~\ref{geo:selection},
$N_{K/F}\beta_n=1$ and $\beta_n$ generates the displayed ideal, which
is supported on the primes $\mathfrak P_v,\iota\mathfrak P_v$; hence
$\beta_n\in\Gamma$. Two elements of relative norm one generating the
same ideal differ by an element of $\mathcal O_K^1$, and distinct labels
give distinct ideals.
\end{proof}

If $\beta\in\beta_n\mathcal O_K^1$, then $j_v(\beta)$ has valuation
$n_v-n_{0,v}$ and $j_v(\iota\beta)=j_v(\beta)^{-1}$. Let
$\varpi_0\in\mathbb A_{\mathcal P}$ have coordinates
$(\varpi_v^{n_{0,v}},\varpi_v^{-n_{0,v}})$. Then $\varpi_0\beta_f$ has
coordinates $(\varpi_v^{n_v}\omega_v,\varpi_v^{-n_v}\omega_v^{-1})$ with
units $\omega_v\in\mathcal O_v^\times$.

\subsection{The Finite-Place Transfer}

\begin{proposition}\label{fw:transfer}
Let $0<\delta<1$ and $C\in\R$. Let $(K_j/F_j)_{j\ge1}$ be quadratic
extensions satisfying (G1)--(G3) with a common $\lambda$, with
$d_j=[F_j:\Q]\to\infty$ and $\log L_{F_j}(1)\le d_jC$. For each $j$ let
$\mathcal P_j$ be a finite set of finite places of $F_j$ that split in
$K_j$, and for $v\in\mathcal P_j$ let $g_v$ be a shell profile on
$F_v^2$ of type $(Q_v,k_v,w^{(v)})$, where the types range over a fixed
finite set independent of $j$. Let $\mathcal F_{\delta,v}=\mathcal
F_\delta(g_v)$, let $(f_\R,g_\C)$ be admissible with Fourier constants
$(M,\sigma)$, and suppose that \eqref{fw:fourier-condition} holds for
every $j$, with $\mathcal C_f=\prod_{v\in\mathcal P_j}(\mathcal
O_v\times\varpi_v^{-k_v}\mathcal O_v)$, that is, with
$H_f=d_j^{-1}\sum_{v\in\mathcal P_j}k_v\log Q_v$. Put
\begin{equation}\label{fw:margin}
 \begin{split}
 \mathcal M_{\rm fin}(\theta)={}&
 \frac1d\sum_{v\in\mathcal P}\log\mathcal F_{\delta,v}
 -(1/2-\delta)\ell-C+(1-\delta)\log2+(1-\theta)\log\pi\\
 &+(1-2\theta)J_\R+\theta J_\C ,
 \end{split}
\end{equation}
evaluated for $K_j/F_j$. If these margins have a positive lower bound
over $j$, then there are finite sets $U_j\subset\R^2$ with
$|U_j|\to\infty$ and $u(U_j)/|U_j|^{1+\delta}\to\infty$.
\end{proposition}

\begin{proof}
Fix $\eta>0$ with $4\eta$ less than the lower bound of the margins, and
consider one extension $K/F$ of the sequence, with $d$ large.

\emph{The profile and its window.} On $\C^d\times\mathbb A_{\mathcal
P}$ let $f=\prod_vf_\R\prod_jg_\C\prod_{v\in\mathcal P}g_v$, with one
factor for each archimedean block and each finite place, let
$V=-\log f$ on the set where $f>0$, and put
\[
 A=\int f^p=A_\R^bA_\C^c\prod_{v\in\mathcal P}A(g_v),\qquad
 Z=I_\R^bI_\C^c\prod_{v\in\mathcal P}Z(g_v).
\]
At a finite place, the \emph{endpoint law} is the probability density
$g_v(z)g_v(z+s(n,\omega))/Z(g_v)$ on $F_v^2\times\Z\times\mathcal
O_v^\times$; the energy $-\log g_v$ of its first endpoint takes finitely
many values. The map $(z,n,\omega)\mapsto(-z-s(n,\omega),n,\omega)$
preserves this law and carries the first endpoint to minus the second,
and $g_v$ is even. Together with the archimedean endpoint laws of
Section~\ref{geo:section}, the product of all these laws is a
probability law under which the energies $V$ of the two endpoints have
the same distribution and are sums of independent terms. Let $d\bar m$
be their mean. The residue characteristics of the places in $\mathcal
P$ lie in a fixed finite set, so $|\mathcal P|=O(d)$, and since the
types range over a finite set, the variance of either energy is $O(d)$.
As in the proof of Proposition~\ref{geo:transfer}, with probability at
least $1/2$ for large $d$, both energies lie in
$[d(\bar m-\eta),d(\bar m+\eta)]$.

Let $\Omega$ be the set of points of $\C^d\times\mathbb A_{\mathcal P}$
where $f>0$ and $V\le d(\bar m+\eta)$. It is a bounded Borel set: the
archimedean energies are coercive and bounded below, and the finite
part lies in the compact support of $\prod_vg_v$, on which the finite
energies are bounded. It is invariant under rotations of the
archimedean coordinates, under $(x_v,y_v)\mapsto(\omega x_v,\omega'y_v)$
for units at each $v\in\mathcal P$, and under translation by
$\mathcal C_f$. For $u\in\R^c$ and $n\in\Z^{\mathcal P}$ let
$\Phi(u,n)=|\Omega\cap(\Omega-\nu(u,n))|$, and let
$\mathcal W(n)=\int_{\R^c}\Phi(u,n)\,du$. By the unit invariance, the
overlap of $\Omega$ with its translate by a vector whose archimedean
part has the moduli of $\nu(u)$ and whose finite part is
$(\varpi_v^{n_v}\omega_v,\varpi_v^{-n_v}\omega_v^{-1})_v$ equals
$\Phi(u,n)$. Only finitely many $n$ have $\mathcal W(n)\ne0$. Undoing the
endpoint density on the event of the previous paragraph, as in the
proof of Proposition~\ref{geo:transfer}, gives
\[
 \sum_{n}\mathcal W(n)\ge\tfrac12Ze^{2d(\bar m-\eta)},
 \qquad
 |\Omega|\le Ae^{pd(\bar m+\eta)} .
\]

\emph{Class selection and rescaling.} Apply Lemma~\ref{fw:selection} to
the finitely many $n$ with $\mathcal W(n)>0$, obtaining $n_0$,
$\mathcal N'$ and $\beta_n$, and let $\varpi_0$ be as above. Put
$\Omega'=\varpi_0^{-1}\Omega$, where $\varpi_0$ acts on the finite part.
Then $|\Omega'|=|\Omega|$, and $\Omega'$ is invariant under the
rectangular group $\mathcal C'_f=\varpi_0^{-1}\mathcal C_f$, of measure
$|\mathcal C_f|$. For $n\in\mathcal N'$, $\beta\in\beta_n\mathcal O_K^1$
and $h\in\R^c$, the element $\beta_h=(W_h\beta_\infty,\beta_f)\in
\Gamma_h$ satisfies
\[
 |\Omega'\cap(\Omega'-\beta_h)|=|\Omega\cap(\Omega-(W_h\beta_\infty,
 \varpi_0\beta_f))|=\Phi(l(\beta)-h,n),
\]
by the last paragraph of the previous subsection and the unit
invariance.

\emph{Points and edges.} For $h\in\R^c$ and $z_0$ in the fundamental
domain $W_hP_{\mathfrak L'}\times\mathcal C'_f$ of $\Gamma_h$ given by
Lemma~\ref{fw:lattice}, where $\mathfrak L'$ corresponds to
$\mathcal C'_f$, let $\mathcal X(h,z_0)=(\Gamma_h+z_0)\cap\Omega'$ and
$N(h,z_0)=|\mathcal X(h,z_0)|$, and let $\mathcal E(h,z_0)$ count the
pairs $(x,\beta)$ with $x\in\mathcal X(h,z_0)$,
$\beta\in\bigcup_{n\in\mathcal N'}\beta_n\mathcal O_K^1$ and
$x+\beta_h\in\Omega'$. Parametrize $z_0$ by $y\in[0,1)^{2d}$ through a
basis of $\mathfrak L'_\infty$, as in Section~\ref{geo:section}, and by
the Haar probability on $\mathcal C'_f$. Since
$\mathbf 1_{\Omega'}\le e^{pd(\bar m+\eta)}(f\circ\varpi_0)^p$ and
$(f\circ\varpi_0)^p$ is the product of an archimedean factor and a
function invariant under $\mathcal C'_f$, Lemma~\ref{fw:counting} gives
\[
 N(h,z_0)\le N_{\max}=\frac{2Ae^{pd(\bar m+\eta)}}{\covol_0}
\]
for all $(h,z_0)$. By tiling, the average of $N(h,z_0)$ over $z_0$ is
$|\Omega'|/\covol_0$, and the average of $\mathcal E(h,z_0)$ is
$\covol_0^{-1}\sum_{n\in\mathcal N'}\sum_{\beta\in\beta_n\mathcal O_K^1}
\Phi(l(\beta)-h,n)$. Averaging over $h\in\Pi^1$ as in
Lemma~\ref{geo:unfold-lemma}, the average of $\mathcal E$ becomes
\[
 \frac{w_K}{\operatorname{Reg}^1\covol_0}\sum_{n\in\mathcal N'}\mathcal W(n)
 \ge\frac{w_K}{\kappa h_{\rm rel}\operatorname{Reg}^1\covol_0}
 \sum_n\mathcal W(n).
\]
Measurability and integrability hold as in
Lemma~\ref{geo:joint-average-lemma}, since $\mathcal E\le N^2\le
N_{\max}^2$. The argument of that lemma yields a pair $(h,z_0)$ with
$N=N(h,z_0)>0$ and
\[
 \frac{\mathcal E(h,z_0)}N\ge\frac{w_K}{\kappa h_{\rm rel}
 \operatorname{Reg}^1}\,\frac{\sum_n\mathcal W(n)}{|\Omega|}
 \ge\frac{w_K}{\kappa h_{\rm rel}\operatorname{Reg}^1}\,
 \frac{Ze^{2d(\bar m-\eta)}}{2Ae^{pd(\bar m+\eta)}} .
\]
With $N\le N_{\max}$ and $p(1+\delta)=2$,
\[
 \frac{\mathcal E(h,z_0)}{N^{1+\delta}}\ge
 \frac{w_K\covol_0^\delta}{2^{1+\delta}\kappa h_{\rm rel}
 \operatorname{Reg}^1}\,\frac{Z}{A^{1+\delta}}\,e^{-4\eta d}.
\]
By \eqref{geo:mass-density} and $\log L_F(1)\le dC$, the first factor
is at least
$2^{-2-\delta}2^{d(1-\delta)}\pi^{(1-\theta)d}\lambda^{-(1/2-\delta)d}e^{-dC}$,
and $Z/A^{1+\delta}=\exp\bigl(bJ_\R+cJ_\C+\sum_v\log\mathcal
F_{\delta,v}\bigr)$. Hence
$\mathcal E/N^{1+\delta}\ge2^{-2-\delta}e^{d(\mathcal M_{\rm
fin}(\theta)-4\eta)}$.

\emph{The planar set.} Let $v_0$ be a real place of $F$ and let $U$ be
the set of $v_0$-coordinates of the points of $\mathcal X(h,z_0)$.
Distinct points differ by $(W_h\gamma_\infty,\gamma_f)$ with
$0\ne\gamma\in\Gamma$, whose $v_0$-coordinate is nonzero, so $|U|=N$.
Each counted pair $(x,\beta)$ gives two points at distance
$|\sigma_{v_0}(\beta)|=1$, and distinct pairs give distinct ordered
pairs of points. As in the proof of Proposition~\ref{geo:transfer},
$u(U)\ge\mathcal E/2$, the ratio $u(U)/|U|^{1+\delta}$ tends to
infinity, and the bound $u(U)\le|U|^2/2$ forces $|U|\to\infty$.
\end{proof}

\subsection{Uniform Local Types}
Suppose that $K/F$ has uniform local types $(e_r,f_r)_{r\in\mathcal R}$
in the sense of Section~\ref{geo:section}, that $\mathcal P$ is the set
of all places of $F$ above $\mathcal R$, and that at every place above
$r$ we use a shell profile of the same type $(r^{f_r},k_r,w^{(r)})$.
Write $\mathcal F_{\delta,r}$ for its functional. Since there are
$d/(e_rf_r)$ places above $r$,
\begin{equation}\label{fw:uniform}
 \frac1d\sum_{v\in\mathcal P}\log\mathcal F_{\delta,v}
 =\sum_{r\in\mathcal R}\frac{\log\mathcal F_{\delta,r}}{e_rf_r},
 \qquad
 H_f=\sum_{r\in\mathcal R}\frac{k_r\log r}{e_r}=H .
\end{equation}
In particular \eqref{fw:fourier-condition} coincides with
\eqref{geo:poisson-condition}, whatever the shell weights. For hard
windows, Lemma~\ref{fw:shell-lemma} gives
$\log\mathcal F_{\delta,r}=\log(k_r+1)-\delta k_rf_r\log r$, and the sum
in \eqref{fw:uniform} is $J-\delta H$. Thus
Proposition~\ref{fw:transfer} contains Proposition~\ref{geo:transfer},
and the shell weights replace the term $J-\delta H$ of
\eqref{geo:general-margin} by the finite functional without changing the
Fourier condition.

\section{Archimedean Profiles}
\label{pf:section}

We use a Gaussian at the compact coordinates, and a positive polynomial modification of a
Student profile at the pair coordinates. For $(z,w)\in\C^2$ put
\[
 t=(1+a|z|^2)^{-1},\qquad t'=(1+a|w|^2)^{-1},\qquad
 g_\C(z,w)=t^s\,t'^{\,s}P(t,t'),
\]
where
\[
 s=1.14600585,\qquad a=0.0000128633241758,
\]
and
\[
 P(t,t')=\sum_{i,j=0}^3\beta_{ij}\,b_i^3(t)\,b_j^3(t'),\qquad
 b_i^m(t)=\binom mi t^i(1-t)^{m-i},
\]
with the symmetric matrix of Bernstein coefficients
\begin{equation}\label{pf:bernstein-witness}
 (\beta_{ij})=\begin{pmatrix}
 1&2.4129438977&2.7721799210&6.6857798061\\
 2.4129438977&3.8532983888&4.6371750119&8.1307168183\\
 2.7721799210&4.6371750119&5.4628959084&9.5512647588\\
 6.6857798061&8.1307168183&9.5512647588&13.8659430610
 \end{pmatrix}.
\end{equation}
All finite decimals here are exact rational numbers. The Bernstein
polynomials $b_i^m$ are nonnegative on $[0,1]$ and sum to one, so
$\min\beta\le P\le\max\beta$ on $[0,1]^2$, where $\min\beta$ and
$\max\beta$ are the smallest and largest coefficients.

The shape of $g_\C$ is adapted to the displacements at a pair of
coordinates. By \eqref{geo:norm-modulus} they have moduli
$(e^u,e^{-u})$, so as $u$ varies they run along a hyperbola. A profile
of width $a^{-1/2}$ in each coordinate keeps a large overlap with its
translate for all $u$ with $e^{|u|}\lesssim a^{-1/2}$, a range of length
about $\log(1/a)$; this produces the factor $\log(1/a)$ in
Proposition~\ref{pf:overlap}, against a mass penalty of order
$\delta\log(1/a)$. For Student profiles the overlaps along this
hyperbola reduce to beta integrals and digamma values
(Lemmas~\ref{pf:laplace}--\ref{pf:pair-reduction}), which makes
certified bounds possible, and the polynomial $P$ couples the two radii.

In this section we prove that these profiles are admissible in the sense
of Definition~\ref{geo:admissible}, and we derive bounds for $J_\R$ and
$J_\C$ that reduce to finite computations with rational data.
Section~\ref{sec:certificate} carries out these computations at
$\delta=\deltastar$. Throughout, $0<\delta<1$ and $p=2/(1+\delta)$. We
write $\psi=\Gamma'/\Gamma$, $\gamma$ for Euler's constant, $(x)_n$ for
the rising factorial, $H_n=\sum_{k=1}^n1/k$ with $H_0=0$, and
$\mathrm{Beta}(\alpha,\alpha')$ for the beta distribution on $(0,1)$
with density proportional to $t^{\alpha-1}(1-t)^{\alpha'-1}$.

\subsection{The Gaussian at the Compact Coordinates}

\begin{lemma}\label{pf:gaussian}
Let $a_\R=2p\delta$ and $f_\R(z)=e^{-a_\R|z|^2/p}$, so that
$f_\R^p=e^{-a_\R|z|^2}$. Then
\begin{equation}\label{pf:gaussian-compact}
 J_\R=\log(p/2)+\delta\log(a_\R/\pi)-a_\R/(2p)
 =\log(p/2)+\delta\bigl(\log(a_\R/\pi)-1\bigr),
\end{equation}
and $a_\R$ maximizes $J_\R$ over the Gaussians $e^{-a'|z|^2/p}$, $a'>0$.
\end{lemma}

\begin{proof}
For $f=e^{-a'|z|^2/p}$ we have $A_\R=\pi/a'$ and, since
$|z|^2+|z+1|^2=2|z+1/2|^2+1/2$, $I_\R=(\pi p/(2a'))e^{-a'/(2p)}$. Hence
$J_\R=\log(p/2)+\delta\log(a'/\pi)-a'/(2p)$, a concave function of
$\log a'$ with its maximum at $a'=2p\delta$.
\end{proof}

By \eqref{geo:general-margin}, the margin is affine in $\theta$ with
slope $-\log\pi-2J_\R+J_\C$. This slope depends on the profiles; for
ours it is positive (Section~\ref{sec:certificate}).

\subsection{The Pair Convolution}

\begin{lemma}\label{pf:laplace}
Let $\alpha,\alpha'>0$ with $A=\alpha+\alpha'-1>0$, and let $h\in\C$.
Then
\begin{equation}\label{pf:laplace-convolution}
 \int_{\C}(1+|z|^2)^{-\alpha}(1+|z+h|^2)^{-\alpha'}\,dz
 =\frac\pi A\,\mathbb E\bigl(1+T(1-T)|h|^2\bigr)^{-A},
 \qquad T\sim\mathrm{Beta}(\alpha,\alpha').
\end{equation}
\end{lemma}

\begin{proof}
Insert
$(1+|z|^2)^{-\alpha}=\Gamma(\alpha)^{-1}\int_0^\infty
r^{\alpha-1}e^{-r(1+|z|^2)}\,dr$ and the analogous formula with a
variable $r'$. Since
$r|z|^2+r'|z+h|^2=(r+r')|z+r'h/(r+r')|^2+rr'|h|^2/(r+r')$, the
integral of $e^{-r|z|^2-r'|z+h|^2}$ over $z\in\C$ is
$\pi(r+r')^{-1}e^{-rr'|h|^2/(r+r')}$; the factors $e^{-r}$ and
$e^{-r'}$ remain in the Laplace integrals. Substitute $r=RT$,
$r'=R(1-T)$, with Jacobian $R$, and integrate over $R$. This gives
\[
 \frac{\pi\,\Gamma(A)}{\Gamma(\alpha)\Gamma(\alpha')}
 \int_0^1T^{\alpha-1}(1-T)^{\alpha'-1}\bigl(1+T(1-T)|h|^2\bigr)^{-A}dT,
\]
which equals the right side of \eqref{pf:laplace-convolution} because
$\Gamma(\alpha+\alpha')=A\,\Gamma(A)$. All integrands are nonnegative,
so Tonelli's theorem justifies the interchanges.
\end{proof}

For $A,B>0$ and $y>0$ define
\[
 Q_{AB}(y)=\int_\R(1+ye^{2v})^{-A}(1+ye^{-2v})^{-B}\,dv .
\]

\begin{lemma}\label{pf:hyperbola}
Let $A,B>0$.
\begin{enumerate}
\item[(a)] For $y>0$,
\begin{equation}\label{pf:elementary-hyperbola-error}
 \Bigl|Q_{AB}(y)-\log(1/y)+\frac{\psi(A)+\psi(B)+2\gamma}{2}\Bigr|
 \le\frac{(A+B)y}2 .
\end{equation}
\item[(b)] For $0<y<1$, with absolutely convergent series,
\[
 Q_{AB}(y)=\sum_{n\ge0}\frac{(A)_n(B)_n}{(n!)^2}\,y^{2n}
 \Bigl[\log(1/y)+\psi(n+1)-\frac{\psi(A+n)+\psi(B+n)}2\Bigr].
\]
\item[(c)] Suppose $A,B\ge1$, $0<y<1$, $ABy^2<1$ and
$2\log(1/y)\ge H_{\lceil A\rceil-1}+H_{\lceil B\rceil-1}$. Then every
term of the series in (b) is nonnegative, and for every integer $m\ge0$
its remainder $\mathcal R_m(y)$ after the terms $n\le m$ satisfies
\[
 0\le\mathcal R_m(y)\le\log(1/y)\frac{(ABy^2)^{m+1}}{1-ABy^2}.
\]
\end{enumerate}
\end{lemma}

\begin{proof}
(a) Split the integral at $v=0$. On $v\ge0$, first replace
$(1+ye^{-2v})^{-B}$ by $1$. Since
$0\le1-(1+ye^{-2v})^{-B}\le Bye^{-2v}$, this changes the half-integral
by an amount in $[0,By/2]$. The substitution $x=ye^{2v}$ gives
$\int_0^\infty(1+ye^{2v})^{-A}dv=\frac12\int_y^\infty(1+x)^{-A}dx/x$,
and
\[
 \int_y^\infty\frac{dx}{x(1+x)^A}
 =\log\frac1y+c_A+\int_0^y\bigl(1-(1+x)^{-A}\bigr)\frac{dx}x,\qquad
 c_A=\int_0^\infty\bigl((1+x)^{-A}-\mathbf 1_{x<1}\bigr)\frac{dx}x .
\]
The last integral lies in $[0,Ay]$, because $0\le1-(1+x)^{-A}\le Ax$;
in particular it tends to $0$ with $y$. On the other hand, the
substitution $t=1/(1+x)$ gives
\[
 \int_y^\infty\frac{dx}{x(1+x)^A}
 =\int_0^{1/(1+y)}\frac{t^{A-1}-1}{1-t}\,dt+\log\frac{1+y}y .
\]
Letting $y\to0$ in the two expressions yields
$c_A=\int_0^1(t^{A-1}-1)(1-t)^{-1}dt=-\psi(A)-\gamma$
\cite[Equation~(5.9.16)]{DLMF}. Thus the half $v\ge0$ equals
$\frac12(\log(1/y)-\psi(A)-\gamma)$ plus a gain in $[0,Ay/2]$ minus a
loss in $[0,By/2]$. The half $v\le0$ is the same with $A$ and $B$
exchanged. Adding the halves proves (a).

(b) The substitutions $x=ye^{2v}$ and then $t=x/(x+y^2)$ give
\[
 2Q_{AB}(y)=\int_0^1t^{B-1}(1-t)^{A-1}\bigl(1-(1-y^2)t\bigr)^{-A}dt .
\]
By Euler's integral \cite[Equation~(15.6.1)]{DLMF}, the right side is
$\Gamma(A)\Gamma(B)\,\mathbf F(A,B;A+B;1-y^2)$, where $\mathbf F$ is
Olver's regularized hypergeometric function. The expansion
\cite[Equation~(15.8.10)]{DLMF}, in the case $c=a+b$, gives
\[
 \mathbf F(A,B;A+B;1-y^2)=\frac1{\Gamma(A)\Gamma(B)}\sum_{n\ge0}
 \frac{(A)_n(B)_n}{(n!)^2}\,y^{2n}\bigl[2\psi(n+1)-\psi(A+n)-\psi(B+n)
 -\log y^2\bigr]
\]
for $0<y<1$. Dividing by two proves (b). The coefficients grow at most
polynomially in $n$, so the series converges absolutely.

(c) For $A\ge1$ and $n\ge0$, monotonicity of $\psi$ and the recurrence
$\psi(x+1)=\psi(x)+1/x$ give
$\psi(n+1)\le\psi(A+n)\le\psi(\lceil A\rceil+n)\le\psi(n+1)+H_{\lceil
A\rceil-1}$. Hence the bracket in (b) lies between
$\log(1/y)-\frac12(H_{\lceil A\rceil-1}+H_{\lceil B\rceil-1})\ge0$ and
$\log(1/y)$. Moreover $(A)_n/n!=\prod_{k<n}(A+k)/(k+1)\le A^n$ for
$A\ge1$, and similarly for $B$. Summing the geometric series proves
(c).
\end{proof}

Let $d_{ij}$ be the monomial coefficients of $P$, so that
$P(t,t')=\sum_{i,j}d_{ij}t^it'^{\,j}$, and put $\eta_0=2s-1$,
\[
 A_{ik}=\eta_0+i+k,\qquad B_{jl}=\eta_0+j+l,\qquad
 w_{ijkl}=\frac{d_{ij}d_{kl}}{A_{ik}B_{jl}} .
\]
We call a quadruple $(i,j,k,l)$ with $d_{ij}d_{kl}\ne0$ a \emph{cross}.
In the pair overlap we rename the variable of integration $v$, so that
$I_\C=\int_\R\int_{\C^2}g_\C(z,w)g_\C(z+e^v,w+e^{-v})\,dz\,dw\,dv$.

\begin{lemma}\label{pf:pair-reduction}
Let $s\ge1$ and $a>0$. Then
\[
 I_\C=\frac{\pi^2}{a^2}\sum_{(i,j,k,l)}w_{ijkl}\,
 \mathbb E\,Q_{A_{ik}B_{jl}}(y),\qquad
 y=a\sqrt{T(1-T)T'(1-T')}\le a/4,
\]
where $T\sim\mathrm{Beta}(s+i,s+k)$ and $T'\sim\mathrm{Beta}(s+j,s+l)$
are independent. Each expectation is finite.
\end{lemma}

\begin{proof}
Expand $g_\C(z,w)=\sum_{i,j}d_{ij}(1+a|z|^2)^{-s-i}(1+a|w|^2)^{-s-j}$ and
multiply out $g_\C(z,w)g_\C(z+e^v,w+e^{-v})$. For one cross, the scaling
$z\mapsto z/\sqrt a$ and Lemma~\ref{pf:laplace} give
\[
 \int_\C(1+a|z|^2)^{-s-i}(1+a|z+e^v|^2)^{-s-k}\,dz
 =\frac\pi{aA_{ik}}\,\mathbb E\bigl(1+aT(1-T)e^{2v}\bigr)^{-A_{ik}},
\]
and similarly in $w$ with $e^{-2v}$, $B_{jl}$ and $T'$. By Tonelli's
theorem, the $v$-integral of the product is
$\mathbb E\int_\R(1+Ye^{2v})^{-A_{ik}}(1+Y'e^{-2v})^{-B_{jl}}dv$ with
$Y=aT(1-T)$ and $Y'=aT'(1-T')$, and the translation
$v\mapsto v+\frac14\log(Y'/Y)$ turns the inner integral into
$Q_{A_{ik}B_{jl}}(\sqrt{YY'})$. By
\eqref{pf:elementary-hyperbola-error}, $Q_{AB}(y)$ is at most
$\log(1/y)$ plus a constant, and $\log(1/(T(1-T)))$ is integrable under
every beta law; so each term is finite and the finite signed sum is
legitimate.
\end{proof}

\subsection{Admissibility}

For the Fourier bound let $m\ge1$ and write
$P(t,t')=\sum_{i,j=0}^m\beta_{ij}b_i^m(t)b_j^m(t')$ with all
$\beta_{ij}>0$. Let $\Delta_1$ and $\Delta_2$ be the forward difference
operators in the two indices. Fix $0<\rho<1$, the width of the complex
tube in the proof below, and put $\omega=(\rho+\rho^2)/(1-\rho^2)$,
\[
 E_{rr'}=\binom mr\binom m{r'}\max_{i,j}\bigl|\Delta_1^r\Delta_2^{r'}
 \beta_{ij}\bigr|,\qquad
 q_0=\frac1{\min\beta}\sum_{r+r'>0}E_{rr'}\,\omega^{r+r'} .
\]

\begin{lemma}\label{pf:tube}
Let $s>0$, $sp>1$ and $a>0$, and suppose $q_0<1$. Then for all
$\xi_1,\xi_2\in\C$,
\[
 \bigl|\widehat{g_\C^p}(\xi_1,\xi_2)\bigr|
 \le K_\C e^{-\sigma_\C(|\xi_1|+|\xi_2|)}A_\C,\qquad
 K_\C=(1-\rho^2)^{-2sp}(1+q_0)^p,\qquad
 \sigma_\C=2\pi\rho/\sqrt a .
\]
\end{lemma}

\begin{proof}
Write $x=\sqrt a\,z\in\R^2$ and consider complex vectors $x+iy$ with
$x,y\in\R^2$ and $|y|\le\rho$. Put
$\mathcal D(x+iy)=1+(x_1+iy_1)^2+(x_2+iy_2)^2$. Then
\[
 \operatorname{Re}\mathcal D=1+|x|^2-|y|^2\ge(1-\rho^2)(1+|x|^2),\qquad
 \bigl|\mathcal D^{-1}-(1+|x|^2)^{-1}\bigr|
 \le\frac{2|x|\rho+\rho^2}{(1-\rho^2)(1+|x|^2)^2}\le\omega .
\]
So the complexified $t=\mathcal D^{-1}$ lies within $\omega$ of the real
$t_0=(1+|x|^2)^{-1}\in(0,1]$, and likewise for $t'$. By the Bernstein
derivative formula,
$\partial_t^r\partial_{t'}^{r'}P/(r!\,r'!)=\binom mr\binom m{r'}
\sum_{i,j}(\Delta_1^r\Delta_2^{r'}\beta)_{ij}\,b_i^{m-r}(t)b_j^{m-r'}(t')$,
which is at most $E_{rr'}$ in absolute value on $[0,1]^2$. Taylor
expansion of the polynomial $P$ about $(t_0,t'_0)$ therefore gives
$|P(t,t')-P(t_0,t'_0)|\le q_0\min\beta\le q_0P(t_0,t'_0)$. Thus
$P(t,t')$ lies in the right half-plane, and with principal branches,
\[
 \bigl|\mathcal D(x+iy)^{-sp}\mathcal D(x'+iy')^{-sp}P(t,t')^p\bigr|
 \le K_\C\,t_0^{sp}\,t_0'^{\,sp}P(t_0,t'_0)^p ,
\]
where $x'+iy'$ is the variable attached to $w$. The same bounds hold
for $|y|,|y'|<\rho'$ with some $\rho'>\rho$, since $q_0$ depends
continuously on $\rho$, so the left side is holomorphic in each complex
coordinate on a neighborhood of the closed tube.

The function $g_\C^p$ is invariant under rotations of $z$ and of $w$.
Rotate so that $\xi_1$ and $\xi_2$ are positive multiples of the first
basis vector, and shift the first real coordinate of $x$ by $-i\rho$,
then that of $x'$ by $-i\rho$. Each shift is justified by Cauchy's
theorem in one variable: the integrand is holomorphic in the strip and
bounded there by a constant times $(1+|x|^2)^{-sp}(1+|x'|^2)^{-sp}$,
which tends to zero at infinity and is integrable because $sp>1$. After
the shifts the Fourier factor gains the factor
$e^{-2\pi\rho(|\xi_1|+|\xi_2|)/\sqrt a}$, and the displayed bound gives
the lemma.
\end{proof}

\begin{proposition}\label{pf:admissible}
Let $s\ge1$, $a>0$, and let $P$ have positive Bernstein coefficients
$\beta_{ij}$ of some degree $m$ in each variable. Let $0<\delta<1$ with
$a_\R=2p\delta\le1$, let $f_\R$ be the Gaussian of
Lemma~\ref{pf:gaussian}, and let $0<\rho<1$ with $q_0<1$. Then
$(f_\R,g_\C)$ is admissible with Fourier constants
\begin{equation}\label{pf:polynomial-fourier-contract}
 M=\max\{2,\sqrt{K_\C}\},\qquad \sigma=\min\{1,\sigma_\C\}.
\end{equation}
\end{proposition}

\begin{proof}
(P1) The functions $f_\R^p$ and
$g_\C^p=(1+a|z|^2)^{-sp}(1+a|w|^2)^{-sp}P(t,t')^p$ are smooth, because
$P\ge\min\beta>0$ on $[0,1]^2$, and they and all their derivatives are
bounded.

(P2) The integrals at the compact coordinates are computed in
Lemma~\ref{pf:gaussian}. The mass $A_\C$ is finite because $sp>1$
(Proposition~\ref{pf:mass} below), and $0<I_\C<\infty$ because
$g_\C\le\max\beta\,t^st'^{\,s}$ and Lemma~\ref{pf:pair-reduction}
applies to the pure Student profile $t^st'^{\,s}$.

(P3) $V_\R=a_\R|z|^2/p$, and
$V_\C=s\log(1+a|z|^2)+s\log(1+a|w|^2)-\log P(t,t')$ is at least
$-\log\max\beta$ and tends to infinity with $|z|$ or $|w|$.

(P4) The endpoint law at a compact coordinate is Gaussian. For the pair,
$-\log P$ is bounded, so it suffices to bound the second moments of
$\log(1+a|z|^2)$ and $\log(1+a|w|^2)$. Choose $0<\tau<s$; then
$\tau<2s-1$ as well, since $s\ge1$. We have
$\log^2(1+x)\le c_\tau(1+x)^\tau$ for $x\ge0$ and some constant
$c_\tau$, and $g_\C\le\max\beta\,t^st'^{\,s}$. Hence the second moment
of $\log(1+a|z|^2)$ is at most a constant times the integral of
$(1+a|z|^2)^{-(s-\tau)}(1+a|z+e^v|^2)^{-s}(1+a|w|^2)^{-s}
(1+a|w+e^{-v}|^2)^{-s}$. By the argument of
Lemma~\ref{pf:pair-reduction}, with exponents $s-\tau>0$ and $s$ and
$A=2s-\tau-1>0$, this integral is a multiple of
$\mathbb E\,Q_{A,2s-1}(y)$, which is finite by
\eqref{pf:elementary-hyperbola-error}. The same holds for $w$.

(P5) By Lemma~\ref{geo:gaussian} the Gaussian satisfies
$|\widehat{f_\R^p}(\xi)|\le2e^{-|\xi|}A_\R\le Me^{-\sigma|\xi|}A_\R$, and by
Lemma~\ref{pf:tube} the pair profile satisfies the bound with
$K_\C e^{-\sigma_\C(|\xi_1|+|\xi_2|)}\le
M^2e^{-\sigma(|\xi_1|+|\xi_2|)}$.
\end{proof}

\subsection{A Lower Bound for the Pair Overlap}
For $n\ge0$ and integers $i,k\ge0$ define
\[
 K_n(i,k)=\frac{(s+i)_n(s+k)_n}{n!\,(\eta_0+i+k+n)_{n+1}},\qquad
 S_x(j)=\sum_{h=0}^{j-1}\frac1{x+h}\quad(S_x(0)=0),
\]
\[
 r_n(i,k)=-\frac1{\eta_0}+\frac{H_n}2-\frac{S_{\eta_0}(i+k+n)}2
 +S_{\eta_0}(i+k+2n+1)-\frac{S_s(i+n)+S_s(k+n)}2,
\]
$R_n(i,k)=K_n(i,k)r_n(i,k)$, and
\[
 G_n=\sum d_{ij}d_{kl}K_n(i,k)K_n(j,l),\qquad
 \widetilde G_n=\sum d_{ij}d_{kl}\bigl[R_n(i,k)K_n(j,l)+K_n(i,k)R_n(j,l)\bigr],
\]
where the sums run over all crosses. For rational $s$ and $d_{ij}$ these
are rational numbers. Put
\begin{equation}\label{pf:student-constant}
 C_s=-2\psi(s)+\psi(2s)+(2s-1)^{-1}-\gamma ,\qquad z_0=a^2/16 .
\end{equation}

\begin{proposition}\label{pf:overlap}
Let $s\ge1$ and $a>0$, and suppose that for every cross
\begin{equation}\label{pf:cross-threshold}
 A_{ik}B_{jl}z_0<1,\qquad
 2\log(4/a)\ge H_{\lceil A_{ik}\rceil-1}+H_{\lceil B_{jl}\rceil-1},
 \qquad \log(4/a)>1/2 .
\end{equation}
Then for every integer $m\ge0$,
\begin{equation}\label{pf:collected-overlap}
 \begin{split}
 \frac{I_\C}{\pi^2a^{-2}}\ge{}&
 \sum_{n=0}^ma^{2n}\bigl\{G_n[\log(1/a)+C_s]+\widetilde G_n\bigr\}-E_{-,m},\\
 E_{-,m}={}&\log(4/a)\sum_{w_{ijkl}<0}(-w_{ijkl})
 \frac{(A_{ik}B_{jl}z_0)^{m+1}}{1-A_{ik}B_{jl}z_0}.
 \end{split}
\end{equation}
\end{proposition}

\begin{proof}
Start from Lemma~\ref{pf:pair-reduction}. Since $y\le a/4<1$, the
hypotheses of Lemma~\ref{pf:hyperbola}(c) hold at every value of $y$,
with $A=A_{ik}\ge1$ and $B=B_{jl}\ge1$. Split each
$Q_{A_{ik}B_{jl}}(y)$ into the terms $n\le m$ of
Lemma~\ref{pf:hyperbola}(b) and the remainder $\mathcal R_m(y)$.

Consider the term of index $n$ for one cross, with $\alpha=s+i$,
$\alpha'=s+k$, $A=A_{ik}=\alpha+\alpha'-1$, and similarly for $T'$. Its
factor $y^{2n}$ is $a^{2n}(T(1-T))^n(T'(1-T'))^n$, and
$\log(1/y)=\log(1/a)-\frac12\log(T(1-T))-\frac12\log(T'(1-T'))$. The
beta integral gives $\mathbb E(T(1-T))^n=(\alpha)_n(\alpha')_n/(A+1)_{2n}$,
so
\[
 \frac{(A)_n}{A\,n!}\,\mathbb E(T(1-T))^n
 =\frac{(\alpha)_n(\alpha')_n}{n!\,(A+n)_{n+1}}=K_n(i,k).
\]
Weighting by $(T(1-T))^n$ turns $\mathrm{Beta}(\alpha,\alpha')$ into
$\mathrm{Beta}(\alpha+n,\alpha'+n)$, under which the mean of
$\log(T(1-T))$ is $\psi(\alpha+n)+\psi(\alpha'+n)-2\psi(A+2n+1)$. The
recurrence $\psi(x+j)=\psi(x)+S_x(j)$ and $\psi(1)=-\gamma$ give
\[
 \begin{gathered}
 \psi(n+1)=-\gamma+H_n,\qquad
 \psi(\alpha+n)=\psi(s)+S_s(i+n),\qquad
 \psi(\alpha'+n)=\psi(s)+S_s(k+n),\\
 \psi(A+n)=\psi(\eta_0)+S_{\eta_0}(i+k+n),\qquad
 \psi(A+2n+1)=\psi(\eta_0)+S_{\eta_0}(i+k+2n+1),\\
 \tfrac12C_s=-\psi(s)+\tfrac12\psi(\eta_0)+\eta_0^{-1}-\tfrac12\gamma,
 \end{gathered}
\]
the last one because $\psi(2s)=\psi(\eta_0)+1/\eta_0$. Substituting,
\[
 \tfrac12\psi(n+1)-\tfrac12\psi(A+n)
 -\tfrac12\bigl[\psi(\alpha+n)+\psi(\alpha'+n)\bigr]+\psi(A+2n+1)
 =\tfrac12C_s+r_n(i,k).
\]
Hence the expectation of the term of index $n$, multiplied by
$w_{ijkl}$, is
\[
 a^{2n}d_{ij}d_{kl}K_n(i,k)K_n(j,l)\bigl[\log(1/a)+C_s+r_n(i,k)
 +r_n(j,l)\bigr].
\]
Summing over crosses gives the displayed main terms.

By Lemma~\ref{pf:hyperbola}(c),
$0\le\mathcal R_m(y)\le\log(1/y)(ABy^2)^{m+1}/(1-ABy^2)$. The right side
increases with $y$ on $(0,a/4]$, because the derivative of
$y^{2m+2}\log(1/y)$ is $y^{2m+1}((2m+2)\log(1/y)-1)>0$ when
$\log(1/y)>1/2$. So
$0\le\mathcal R_m(y)\le\log(4/a)(ABz_0)^{m+1}/(1-ABz_0)$. For crosses
with $w_{ijkl}>0$ we discard the remainder, and for crosses with
$w_{ijkl}<0$ we subtract its upper bound. This proves
\eqref{pf:collected-overlap}.
\end{proof}

With $m=1$ the bound keeps the exact term of order $a^2$ and leaves a
signed error of order $a^4\log(1/a)$. It is a lower bound for $I_\C$
only.

\subsection{An Upper Bound for the Pair Mass}

\begin{proposition}\label{pf:mass}
Let $s>0$ with $q=sp-1>0$, let $a>0$, and let $P$ be positive on
$[0,1]^2$. Then
\begin{equation}\label{pf:polynomial-mass}
 A_\C=\frac{\pi^2}{a^2q^2}\,\bar A_\C,\qquad
 \bar A_\C=\mathbb E\,P(T,T')^p,
\end{equation}
where $T$ and $T'$ are independent with density $qt^{q-1}$ on $(0,1)$.
Partition $[0,1]^2$ into rectangles $\mathcal Q=[l,r]\times[l',r']$, and
use on $\mathcal Q$ the coordinates $x=(T-l)/(r-l)$ and
$x'=(T'-l')/(r'-l')$. Suppose that the Bernstein coefficients of the
restriction of $P$ to $\mathcal Q$, in these coordinates, lie in
$[m_{\mathcal Q},M_{\mathcal Q}]$ with $m_{\mathcal Q}>0$, and put
\[
 c_{\mathcal Q}=\frac{m_{\mathcal Q}+M_{\mathcal Q}}2,\qquad
 \rho_{\mathcal Q}=\frac{M_{\mathcal Q}-m_{\mathcal Q}}
 {M_{\mathcal Q}+m_{\mathcal Q}},\qquad
 P=c_{\mathcal Q}(1+W_{\mathcal Q})\ \text{on }\mathcal Q .
\]
Define
\begin{equation}\label{pf:subrectangle-moments}
 M_i(l,r)=\int_l^rqv^{q-1}\Bigl(\frac{v-l}{r-l}\Bigr)^idv
 =\frac q{(r-l)^i}\sum_{j=0}^i\binom ij(-l)^{i-j}
 \frac{r^{q+j}-l^{q+j}}{q+j},
\end{equation}
and, writing $W_{\mathcal Q}^k=\sum_{i,j}c^{(k)}_{ij}x^ix'^{\,j}$,
\[
 \Sigma_{\mathcal Q,N}=c_{\mathcal Q}^p\sum_{k=0}^N\binom pk
 \sum_{i,j}c^{(k)}_{ij}M_i(l,r)M_j(l',r').
\]
Then for every $N\ge1$,
\begin{equation}\label{pf:subrectangle-mass}
 \Bigl|\bar A_\C-\sum_{\mathcal Q}\Sigma_{\mathcal Q,N}\Bigr|
 \le\sum_{\mathcal Q}c_{\mathcal Q}^p
 \frac{|\binom p{N+1}|\rho_{\mathcal Q}^{N+1}}{1-\rho_{\mathcal Q}}
 M_0(l,r)M_0(l',r').
\end{equation}
\end{proposition}

\begin{proof}
The substitution $t=(1+a|z|^2)^{-1}$ gives $dz=\pi\,dt/(at^2)$ on $\C$.
Hence
\[
 A_\C=\frac{\pi^2}{a^2}\int_0^1\!\!\int_0^1t^{sp-2}t'^{\,sp-2}
 P(t,t')^p\,dt\,dt',
\]
which is \eqref{pf:polynomial-mass} since $sp-2=q-1$. The Bernstein
convex-hull property gives $|W_{\mathcal Q}|\le\rho_{\mathcal Q}<1$ on
$\mathcal Q$. For $1<p<2$ the binomial coefficients satisfy
$|\binom p{k+1}|/|\binom pk|=|p-k|/(k+1)<1$ for $k\ge1$, so the binomial
series of $(1+W)^p$ converges uniformly on $|W|\le\rho_{\mathcal Q}$,
and its tail after index $N\ge1$ is at most
$|\binom p{N+1}|\rho_{\mathcal Q}^{N+1}/(1-\rho_{\mathcal Q})$.
Integrate against the density on $\mathcal Q$, whose total mass is
$M_0(l,r)M_0(l',r')$. The moment formula follows by expanding
$(v-l)^i$; the convention $0^q=0$ is valid since $q>0$.
\end{proof}

When $p$, $q$, the partition and the coefficients of $P$ are rational,
the restrictions, the Bernstein coefficients, the powers
$W_{\mathcal Q}^k$ and the binomial coefficients are exact rationals,
computed by integer polynomial arithmetic after clearing denominators.
The endpoint powers $r^{q+j}$ and the factors $c_{\mathcal Q}^p$ need
not be rational; they are enclosed by interval evaluation of
$\exp(y\log x)$. Exact rational algebra followed by interval evaluation
keeps the bound valid despite the cancellation in
\eqref{pf:subrectangle-moments}. If $P$ is symmetric and both axes are
partitioned alike, a rectangle and its reflection contribute equally,
so off-diagonal rectangles may be counted twice and diagonal ones once.

\subsection{The Pair Functional}

\begin{corollary}\label{pf:functional}
Under the hypotheses of Propositions~\ref{pf:overlap} and~\ref{pf:mass},
let $Z_*$ be the right side of \eqref{pf:collected-overlap} for some
$m$, and suppose $Z_*>0$. If $\bar A_\C^*\ge\bar A_\C$, then
\begin{equation}\label{pf:collected-functional}
 J_\C\ge\log Z_*+2(1+\delta)\log q-2\delta\log\pi
 +2\delta\log a-(1+\delta)\log\bar A_\C^* .
\end{equation}
\end{corollary}

\begin{proof}
By Proposition~\ref{pf:overlap}, $I_\C\ge\pi^2a^{-2}Z_*$, and by
\eqref{pf:polynomial-mass}, $A_\C\le\pi^2a^{-2}q^{-2}\bar A_\C^*$.
Substitute into $J_\C=\log I_\C-(1+\delta)\log A_\C$.
\end{proof}

Interval evaluation of the right side of \eqref{pf:collected-functional}
gives a certified lower bound for $J_\C$; the upper endpoint of such an
enclosure is not an upper bound for $J_\C$. For rational $s$, all
digamma values in Proposition~\ref{pf:overlap} reduce to rational
corrections and to $C_s=2/\eta_0+\psi(\eta_0)-2\psi(s)+\psi(1)$. The
digamma function at rational arguments is enclosed as follows.

\begin{lemma}\label{pf:digamma}
For $x>0$ and integers $L\ge0$ and $m\ge1$, put $z=x+L$. Then
\[
 \Bigl|\psi(x)-\Bigl(\log z-\frac1{2z}
 -\sum_{k=1}^{m-1}\frac{B_{2k}}{2kz^{2k}}
 -\sum_{j=0}^{L-1}\frac1{x+j}\Bigr)\Bigr|
 \le\frac{|B_{2m}|}{2mz^{2m}},
\]
where $B_{2k}$ are the Bernoulli numbers.
\end{lemma}

\begin{proof}
By Binet's formula \cite[Equation~(5.9.15)]{DLMF},
\[
 \psi(z)=\log z-\frac1{2z}
 -2\int_0^\infty\frac{t\,dt}{(t^2+z^2)(e^{2\pi t}-1)} .
\]
Write
\[
 \frac1{t^2+z^2}=\sum_{k=1}^{m-1}(-1)^{k-1}\frac{t^{2k-2}}{z^{2k}}
 +(-1)^{m-1}\frac{t^{2m-2}}{z^{2m-2}(t^2+z^2)},
\]
and use $\int_0^\infty t^{2k-1}(e^{2\pi t}-1)^{-1}dt=|B_{2k}|/(4k)$ and
$(-1)^{k-1}|B_{2k}|=B_{2k}$. The remaining integral has a fixed sign and
is at most $|B_{2m}|/(2mz^{2m})$ in absolute value, because
$(t^2+z^2)^{-1}\le z^{-2}$. Finally apply the recurrence
$\psi(x+1)=\psi(x)+1/x$ $L$ times.
\end{proof}

The certificate evaluates the right side of this lemma with exact
rational arguments and Bernoulli numbers, $L=64$ and $m=16$, in outward
rounded interval arithmetic, without rounding the rational parameters of
the profile.

\section{The Final Certificate}
\label{sec:certificate}

We apply Proposition~\ref{fw:transfer} to the fields of
Theorem~\ref{tw:field-family}, with the ceiling of
Theorem~\ref{an:ceiling} and the rational data below. Every parameter
written as a finite decimal is exact. The inequalities between decimals
displayed in this section, and the entries of
Tables~\ref{cert:tab-local} and~\ref{cert:tab-budget}, are computed and
printed by the supplementary program \texttt{geom241.py}
(Section~\ref{cert:computations}); the decimals are rounded outward from
interval enclosures and are never substituted into later calculations.

\subsection{Parameters and Local Data}
Set
\[
 \delta=\frac{4273}{100000}=\deltastar,\qquad p=\frac2{1+\delta}
 =\frac{200000}{104273},\qquad C=\frac{4871285}{10^8}=\ceiling,
\]
\[
 \theta_*=\frac{65535}{131072}.
\]
At the compact coordinates we use the Gaussian of
Lemma~\ref{pf:gaussian}, $f_\R^p=e^{-a_\R|z|^2}$ with
$a_\R=2p\delta=17092/104273<1$. At the pair coordinates we use the
profile $g_\C$ of Section~\ref{pf:section}, with the coefficients
\eqref{pf:bernstein-witness}, $s=1.14600585$ and
$a=0.0000128633241758$. For the Fourier bound we take
$\rho=10^{-3}$. For the mass we partition $[0,1]^2$ into the $64$
squares of side $1/8$ with rational corners and use the binomial degree
$N=18$ in Proposition~\ref{pf:mass}. For the overlap we use $m=1$ in
Proposition~\ref{pf:overlap}.

The selected rational primes are $\mathcal R=\{2,3,5,7,29\}$, with the
uniform local types of Theorem~\ref{tw:field-family}. At every place of
$F$ above $r\in\mathcal R$ we use the shell profile with product
weights $w_{ij}=w_{r,i}w_{r,j}$, in the notation of
\eqref{fw:product-shells} with $x=x'=(w_{r,i})_i$, and with
$Q=Q_r=r^{f_r}$ and $k=k_r$. Table~\ref{cert:tab-local} lists the local
data and Table~\ref{cert:tab-weights} the exact weights. At $7$ and $29$
the weights of the outer shells are zero, so these profiles have only
one and two nontrivial shells, and their gain over the hard window is
small (Table~\ref{cert:tab-local}). By \eqref{fw:product-shells}, the
local functional is
\begin{equation}\label{cert:finite-factor}
 \log\mathcal F_{\delta,r}
 =-\delta k_r\log Q_r+\log\bigl((k_r+1)S_r^2+2S_rR_r\bigr)
 -2(1+\delta)\log P_r,
\end{equation}
where $P_r=\sum_im_iw_{r,i}^p$, $S_r=\sum_im_iw_{r,i}^2$ and
$R_r=\sum_{i,i'}\Delta_{ii'}w_{r,i}w_{r,i'}$, with the shell masses
$m_i$ and the matrix $\Delta$ of Section~\ref{fw:section} for $Q=Q_r$.

\begin{table}[htbp]
\centering\small
\caption{Local data at the selected primes. Here $Q_r=r^{f_r}$, the
profile at $r$ uses the weights $w_{r,0},\dots,w_{r,n_r-1}$, and the
last two columns give $\log\mathcal F_{\delta,r}/(e_rf_r)$ for the shell
profile and for the hard window, rounded to ten decimals.}
\label{cert:tab-local}
\begin{tabular}{rrrrrrrr}
\toprule
$r$ & $e_r$ & $f_r$ & $Q_r$ & $k_r$ & $n_r$ & shells & hard window\\
\midrule
$2$ & $8$ & $4$ & $16$ & $7$ & $6$ & $0.0393401103$ & $0.0390666415$\\
$3$ & $2$ & $2$ & $9$ & $9$ & $6$ & $0.3675484302$ & $0.3643996093$\\
$5$ & $2$ & $2$ & $25$ & $6$ & $6$ & $0.2817099488$ & $0.2801636913$\\
$7$ & $1$ & $8$ & $5764801$ & $1$ & $2$ & $0.0034946608$ & $0.0034946569$\\
$29$ & $1$ & $4$ & $707281$ & $1$ & $3$ & $0.0294023321$ & $0.0294022443$\\
\bottomrule
\end{tabular}
\end{table}

\begin{table}[htbp]
\centering\small
\caption{The exact shell weights $w_{r,i}$ for $i\ge1$; all $w_{r,0}=1$.
The last column gives the weights rounded to six significant digits,
for orientation only; the certificate uses the exact fractions.}
\label{cert:tab-weights}
\begin{tabular}{rrlr}
\toprule
$r$ & $i$ & $w_{r,i}$ & $\approx w_{r,i}$\\
\midrule
2 & 1 & $2854206002066/59041214756005$ & $4.83426\cdot10^{-2}$\\
2 & 2 & $198349814036/95709219081499$ & $2.07242\cdot10^{-3}$\\
2 & 3 & $7762557303/86561984122975$ & $8.96763\cdot10^{-5}$\\
2 & 4 & $109894931/27828186078256$ & $3.94905\cdot10^{-6}$\\
2 & 5 & $17033757/96603874664399$ & $1.76326\cdot10^{-7}$\\
\midrule
3 & 1 & $5693537714695/63764117760919$ & $8.92906\cdot10^{-2}$\\
3 & 2 & $567838527095/78237129879259$ & $7.25792\cdot10^{-3}$\\
3 & 3 & $43356430379/72943614524859$ & $5.94383\cdot10^{-4}$\\
3 & 4 & $2265261068/45953937834745$ & $4.92942\cdot10^{-5}$\\
3 & 5 & $74149907/18018553422030$ & $4.11520\cdot10^{-6}$\\
\midrule
5 & 1 & $1343081168627/46386563301209$ & $2.89541\cdot10^{-2}$\\
5 & 2 & $61788974773/82067065665364$ & $7.52908\cdot10^{-4}$\\
5 & 3 & $1389672153/69896165011649$ & $1.98820\cdot10^{-5}$\\
5 & 4 & $32644248/60812698834313$ & $5.36800\cdot10^{-7}$\\
5 & 5 & $511565/34657937034136$ & $1.47604\cdot10^{-8}$\\
\midrule
7 & 1 & $2309805/69240611498956$ & $3.33591\cdot10^{-8}$\\
\midrule
29 & 1 & $36498535/95956876917664$ & $3.80364\cdot10^{-7}$\\
29 & 2 & $4/29090285629613$ & $1.37503\cdot10^{-13}$\\
\bottomrule
\end{tabular}
\end{table}


The weights are positive and nonincreasing, and the coefficients
$\beta_{ij}$ are positive; these are exact rational comparisons. The
rational number $Z(g_r)Q_r^{-k_r}$ is computed twice for each $r$: from
\eqref{fw:product-shells}, and by writing the profile as the
nonnegative combination
$\sum_{i,j}(w_{r,i}-w_{r,i+1})(w_{r,j}-w_{r,j+1})\mathbf 1_{\mathcal
B_i\times\varpi^{-k_r}\mathcal B_j}$ of indicators of products of balls,
with $w_{r,n_r}=0$, and applying \eqref{fw:rectangle-gram}. The two
exact values agree.

\subsection{Certified Enclosures}
Let $q=sp-1=12492817/10427300$, and let $T$ and $T'$ be independent
with density $qt^{q-1}$ on $(0,1)$. Proposition~\ref{pf:mass}, with the
data above, gives
\begin{equation}\label{cert:mass}
 38.7888697256861058<\bar A_\C=\mathbb E\,P(T,T')^p<38.7888697256861071 .
\end{equation}
Here every local Bernstein coefficient is positive, every local ratio
$\rho_{\mathcal Q}$ is less than $0.332$, and the tail allowances in
\eqref{pf:subrectangle-mass} add up to less than $5.9\cdot10^{-16}$. The
threshold conditions \eqref{pf:cross-threshold} hold for every cross,
and Proposition~\ref{pf:overlap} with $m=1$ gives
\begin{equation}\label{cert:overlap}
 348.4186894704059951<Z_*<348.4186894704059952 ,
\end{equation}
where $Z_*$ is the right side of \eqref{pf:collected-overlap}; its
signed remainder $E_{-,1}$ is less than $6.1\cdot10^{-17}$. The upper
endpoint in \eqref{cert:overlap} bounds this expression, not the
overlap $I_\C$. Let $J_\C^*$ be the right side of
\eqref{pf:collected-functional} with $\bar A_\C^*$ the upper endpoint
of the enclosure \eqref{cert:mass}. By Corollary~\ref{pf:functional},
$J_\C\ge J_\C^*$, and
\[
 1.3556525171042125<J_\C^*<1.3556525171042127 .
\]
Formula~\eqref{pf:gaussian-compact} gives
\[
 -0.2107595968649958<J_\R<-0.2107595968649957 .
\]
For the finite places, \eqref{cert:finite-factor} and the data of
Table~\ref{cert:tab-local} give
\[
 0.7214954821109038790<\sum_{r\in\mathcal R}
 \frac{\log\mathcal F_{\delta,r}}{e_rf_r}<0.7214954821109038791 ,
\]
against $J-\delta H<0.7165268434$ for the hard windows. We also have
\begin{gather*}
 15.6917787983405342<H<15.6917787983405343,\\
 5.6560047235563094<\ell<5.6560047235563095 ,
\end{gather*}
where $H=\sum_rk_r\log r/e_r$ and
$\ell=\frac94\log2+\frac12\log3615$.

For the Fourier condition, the rational tube quantity of
Lemma~\ref{pf:tube} satisfies $q_0<0.026<1$, and
\[
 \log K_\C<0.0491254364,\qquad \sigma_\C>1.7518755886,\qquad
 \mu=2e^{H/2-\ell}>17.8683654522 .
\]
Since $K_\C<4$ and $\sigma_\C>1$, Proposition~\ref{pf:admissible} shows
that $(f_\R,g_\C)$ is admissible with Fourier constants $M=2$ and
$\sigma=1$. The right side of \eqref{geo:poisson-condition} is then
$\log50+2<5.9121$, so \eqref{fw:fourier-condition}, which by
\eqref{fw:uniform} is the same condition with $H_f=H$, holds with a
wide margin: although $\mu$ is not large, only
$\mu\ge\log M+2\log5+2$ is needed.

Define the conservative margin
\[
 \begin{split}
 M_*(\theta)={}&\sum_{r\in\mathcal R}
 \frac{\log\mathcal F_{\delta,r}}{e_rf_r}
 -(1/2-\delta)\ell-C+(1-\delta)\log2+(1-\theta)\log\pi\\
 &+(1-2\theta)J_\R+\theta J_\C^* .
 \end{split}
\]
By \eqref{fw:uniform}, $J_\C\ge J_\C^*$ and $\theta\ge0$, it is a lower
bound for \eqref{fw:margin}. It is affine in $\theta$ with slope
\[
 0.6324418249848039<-\log\pi-2J_\R+J_\C^*<0.6324418249848040 ,
\]
so on $[\theta_*,1/2)$ its minimum is at $\theta_*$, and
\begin{equation}\label{cert:margin}
 M_*(\theta_*)>0.000176730033534598 .
\end{equation}
Table~\ref{cert:tab-budget} lists the terms of $M_*(\theta_*)$. The
proof of Proposition~\ref{fw:transfer} only needs some $\eta>0$ with
$4\eta$ below the margin. For illustration, $\eta=10^{-9}$ gives
\begin{equation}\label{cert:final-margin}
 M_*(\theta_*)-4\eta>0.000176726033534598>0 .
\end{equation}
The margin is small compared with the individual terms. Since
$M_*$ decreases with slope $1$ in $C$ and with slope $1/2-\delta$ in
$\ell$, it would be absorbed by an increase of $1.77\cdot10^{-4}$ in $C$,
or of $3.87\cdot10^{-4}$ in $\ell$. With these profiles and shell
weights optimized in the same way for $\delta=0.0428$, the same
evaluation gives $M_*(\theta_*)<-0.0015$.

\begin{table}[htbp]
\centering\small
\caption{The terms of $M_*(\theta_*)$, rounded to ten decimals.}
\label{cert:tab-budget}
\begin{tabular}{lr}
\toprule
Term & Value\\
\midrule
$\sum_r\log\mathcal F_{\delta,r}/(e_rf_r)$ & $0.7214954821$\\
$-(1/2-\delta)\ell$ & $-2.5863212799$\\
$-C$ & $-0.0487128500$\\
$(1-\delta)\log2$ & $0.6635290015$\\
$(1-\theta_*)\log\pi$ & $0.5723736765$\\
$(1-2\theta_*)J_\R$ & $-0.0000032159$\\
$\theta_*J_\C^*$ & $0.6778159157$\\
\midrule
$M_*(\theta_*)$ & $0.0001767300$\\
\bottomrule
\end{tabular}
\end{table}

\begin{remark}\label{cert:signature-remark}
The certificate uses the bound $\theta\ge\theta_*$ only through the
positivity of $M_*$ on $[\theta_*,1/2)$. The program
\texttt{geom241.py} also encloses the value at $\theta=0$ and the zero
$\theta_0$ of the affine function $M_*$:
\[
 M_*(0)<-0.316,\qquad \theta_0<0.49971293 .
\]
So $M_*(\theta)>0$ for $\theta\ge0.49971293$. By the proof of
Theorem~\ref{tw:field-family}, $b/d=1/(2N_\iota)$, where $N_\iota$ is
the number of conjugates of $\iota_1$ in $\Gal(K/B)$, so that
$\theta=1/2-1/(4N_\iota)$. Hence $N_\iota\ge871$ would suffice, whereas
that proof gives $N_\iota\ge2^{15}$. On the other hand, at $\theta=0$ the margin
\eqref{fw:margin} does not involve $J_\C$ and equals $M_*(0)<0$: with
these profiles, mixed signature is essential.
\end{remark}

\subsection{Proof of the Main Theorem}

\begin{proof}[Proof of Theorem~\ref{thm:main}]
For each $n\ge49$, choose a Galois extension $K_n/B$ of degree $2^n$ as
in Theorem~\ref{tw:field-family} (with $m=2^n$), a complex conjugation
$\iota_1$ at a place of $K_n$ above $v_1$, and put
$F_n=K_n^{\langle\iota_1\rangle}$. Then $d_n=[F_n:\Q]=2^n\to\infty$. By
that theorem, each $K_n/F_n$ satisfies (G1)--(G3) with
$\lambda=2^{9/4}\sqrt{3615}$, has the uniform local types
$(e_r,f_r)_{r\in\mathcal R}$ of Table~\ref{cert:tab-local}, and has
$\theta_*\le\theta_n<1/2$. By Theorem~\ref{an:ceiling},
$d_n^{-1}\log L_{F_n}(1)<C$ for all sufficiently large $n$; we discard
the finitely many other fields.

Let $\mathcal P_n$ be the set of places of $F_n$ above $\mathcal R$,
with the shell profiles of Table~\ref{cert:tab-weights}. Their types
form a fixed set of five. By Proposition~\ref{pf:admissible} and the
enclosures above, $(f_\R,g_\C)$ is admissible with Fourier constants
$M=2$, $\sigma=1$, and \eqref{fw:fourier-condition} holds for every
$n$. By \eqref{fw:uniform}, Corollary~\ref{pf:functional} and the slope
bound, the margin \eqref{fw:margin} of $K_n/F_n$ is at least
$M_*(\theta_n)\ge M_*(\theta_*)$, which is positive by
\eqref{cert:margin}. Proposition~\ref{fw:transfer} therefore gives finite
sets $U_n\subset\R^2$ with $|U_n|\to\infty$ and
$u(U_n)/|U_n|^{1+\delta}\to\infty$, where $1+\delta=1.04273$; these are
the sets of Theorem~\ref{thm:main}, indexed by $n$. In particular
$u(U_n)\ge|U_n|^{1.04273}$ for all large $n$, which gives the last
assertion of Theorem~\ref{thm:main} at the cardinalities $|U_n|$.
\end{proof}

\subsection{Computations}\label{cert:computations}
The program \texttt{geom241.py} evaluates the formulas of
Sections~\ref{fw:section} and~\ref{pf:section} for the data above: the
local functionals \eqref{cert:finite-factor} from the exponents and
weights in the data file \texttt{shells241\_0.04273.json}, and the
quantities of Propositions~\ref{pf:mass} and~\ref{pf:overlap}, of
Lemma~\ref{pf:tube} and of Corollary~\ref{pf:functional} from the
profile data above. For these it uses, unchanged, three certified
modules of an earlier supplementary archive of the author
(Section~\ref{intro:code}), which implement exactly these formulas. It checks
every finite hypothesis used above: the positivity and monotonicity of
the weights, the two exact computations of each $Z(g_r)Q_r^{-k_r}$, the
positivity of the $\beta_{ij}$ and of the local Bernstein coefficients,
the conditions \eqref{pf:cross-threshold}, $q_0<1$, $a_\R\le1$,
$\sigma_\C>1$, $\log K_\C<2\log2$, $\mu>10$ and the positivity of the
slope. It then prints the enclosures stated in this section. The
exponents and weights were produced by the separate optimization
program \texttt{shells241.py}; they enter the certificate only as the
exact rational data of Table~\ref{cert:tab-weights}.

Rational quantities are computed exactly: $q_0$, the coefficients of
the powers $W_{\mathcal Q}^k$ and the binomial coefficients in
Proposition~\ref{pf:mass}, $G_n$ and $\widetilde G_n$, and the shell sums
$S_r$ and $R_r$. The remaining quantities involve the logarithm, the
exponential
(also through the real powers $x^y=\exp(y\log x)$ in
\eqref{pf:subrectangle-moments}, in $c_{\mathcal Q}^p$ and in
$w_{r,i}^p$), the square root in $\sigma_\C$,
the constant $\pi$, and the digamma function, which enters only through
Lemma~\ref{pf:digamma}. They are evaluated in the interval arithmetic of
mpmath \cite{Mpmath}, which computes each arithmetic operation and each
of these elementary functions with the lower endpoint rounded down and
the upper endpoint rounded up, at a working precision of $80$ decimal
digits. Hence every computed interval contains the exact value, and the
decimals stated above, which are coarser than the working precision by
more than fifty orders of magnitude, are rigorous bounds.


\bibliographystyle{amsplain}
\bibliography{references}
\end{document}
